% Suites numériques — Mathématiques Tle Tech — Cours signé OnBuch+
% Programme officiel MINESEC. Compiler avec : tectonic N06-suites-numeriques.tex
\newcommand{\DOCMATIERE}{Mathématiques}
\newcommand{\DOCNIVEAU}{Tle Tech}
\newcommand{\DOCMODULE}{Mathématiques – classe de Terminale technique}
\newcommand{\DOCLECON}{N06}
\newcommand{\DOCTITRE}{Suites numériques}
% OnBuch+ — préambule « cours signés » ENRICHI (compilé avec tectonic / XeLaTeX)
% Système visuel « Pop » d'OnBuch+ : fond crème (jamais blanc), bordures encre
% épaisses, ombres « dures » décalées, titres Archivo Black en capitales.
% Version MATHÉMATIQUES : graphes (pgfplots/tikz), boîtes pédagogiques
% (définition, propriété, méthode, attention, le savais-tu, exercice résolu,
% exercice, TP, bilan), page de garde et signature OnBuch+ ancrée en bas.
%
% Macros attendues dans main.tex AVANT \input{preamble} :
%   \DOCMATIERE  \DOCNIVEAU  \DOCMODULE  \DOCLECON (numéro)  \DOCTITRE
\documentclass[11pt]{article}

\usepackage{fontspec}
\usepackage[french]{babel}
\usepackage[a4paper,margin=2cm,top=3.4cm,bottom=2.8cm,headheight=1.4cm,footskip=1.3cm]{geometry}
\usepackage{amsmath,amssymb,amsfonts}
\usepackage{mathtools} % \xmapsto, \coloneqq... souvent utilisés par les modèles
\usepackage{cancel} % simplifications barrées
% \abs{z} (module, valeur absolue) et \norm{u} (norme), utilisés par les modèles
\providecommand{\abs}{}\let\abs\relax\DeclarePairedDelimiter{\abs}{\lvert}{\rvert}
\providecommand{\norm}{}\let\norm\relax\DeclarePairedDelimiter{\norm}{\lVert}{\rVert}
\usepackage[table]{xcolor} % [table] : fournit aussi \rowcolors (alternance de lignes)
\usepackage{pagecolor}
\usepackage{tikz}
\usetikzlibrary{arrows.meta,positioning,calc,shapes.geometric,shapes.misc,shapes.arrows,decorations.pathmorphing,decorations.pathreplacing,decorations.markings,patterns,shadows,mindmap,trees,fit,backgrounds,matrix,intersections,angles,quotes}
\usepackage{pgfplots}
\usepackage{pgf-pie} % diagrammes circulaires \pie (statistiques)
\pgfplotsset{compat=1.18}
\usepgfplotslibrary{fillbetween,groupplots} % aires entre courbes (calcul intégral)
\usepackage{enumitem}
\usepackage{fancyhdr}
% [most] charge toutes les bibliothèques tcolorbox (listings, minted, raster,
% documentation...) et gonfle inutilement la mémoire TeX sur un document long
% (~300 boîtes) : on ne charge que ce qui est réellement utilisé.
\usepackage[skins,breakable]{tcolorbox}
\usepackage{graphicx}
\usepackage{colortbl}
\usepackage{booktabs}
\usepackage{tabularx}
\usepackage{diagbox} % case d'en-tête diagonale des tableaux à double entrée
% Colonnes extensibles centrée / gauche / droite (C, L, R), souvent utilisées par les modèles
\newcolumntype{C}{>{\centering\arraybackslash}X}
\newcolumntype{L}{>{\raggedright\arraybackslash}X}
\newcolumntype{R}{>{\raggedleft\arraybackslash}X}
\usepackage{array}
\usepackage{multicol}
\usepackage{siunitx}
% parse-numbers=false : accepte aussi \SI{2,5 \times 10^{-3}}{...}
\sisetup{locale=FR,output-decimal-marker={,},group-minimum-digits=4,parse-numbers=false}
\DeclareSIUnit\ohm{\ensuremath{\symup{\Omega}}} % U+2126 absent de la police
\usepackage{qrcode}
% \degree : commande fréquemment utilisée par les modèles pour un angle ou une
% température en dehors de \SI{}{} (non fournie par les paquets chargés).
\providecommand{\degree}{^{\circ}}
% \qed (fin de démonstration), utilisé par les modèles sans amsthm
\providecommand{\qed}{\hfill$\square$}
% \barre : double barre des valeurs interdites dans un tableau de variations (tkz-tab)
\providecommand{\barre}{\ensuremath{\|}}
% environnement proof (amsthm non chargé) utilisé par les modèles
\ifdefined\proof\else\newenvironment{proof}[1][Démonstration]{\par\noindent\textit{#1.}\ }{\qed\par}\fi
\usepackage{lastpage}
\usepackage{hyperref}
\hypersetup{hidelinks}

% ============================================================
% Palette « Pop » OnBuch+ — mêmes hex que la classe P (plus_ui.dart)
% ============================================================
\definecolor{poporange}{HTML}{F59321}
\definecolor{popdark}{HTML}{E07A0C}
\definecolor{popcream}{HTML}{FFF6E8}
\definecolor{popink}{HTML}{1C1714}
\definecolor{poppurple}{HTML}{7A5AE0}
\definecolor{popgreen}{HTML}{2E9E6A}
\definecolor{poppink}{HTML}{E0457B}
\definecolor{popblue}{HTML}{2F7DE1}
\definecolor{popgold}{HTML}{C98A12}
\definecolor{popmuted}{HTML}{8A7F76}
\definecolor{popline}{HTML}{EADFD2}
\definecolor{popgray}{HTML}{8A7F76}
\colorlet{popgrey}{popgray}
% Alias tolérants (le modèle nomme parfois les couleurs différemment)
\colorlet{popwhite}{white}
\colorlet{popblack}{popink}
\colorlet{popred}{poppink}
\colorlet{popyellow}{popgold}
% Teintes claires pour fonds de tableaux / zones
\colorlet{popblueL}{popblue!10!white}
\colorlet{poporangeL}{poporange!14!white}
\colorlet{popgreenL}{popgreen!12!white}
\colorlet{poppurpleL}{poppurple!12!white}
\colorlet{poppinkL}{poppink!12!white}
\colorlet{popgoldL}{popgold!14!white}

\pagecolor{popcream}

% ============================================================
% Typographie
% ============================================================
\defaultfontfeatures{Ligatures=TeX}
\setmainfont{PlusJakartaSans-Medium.ttf}[Path=../../../fonts/,BoldFont=PlusJakartaSans-Bold.ttf]
% Maths en sans-serif (Fira Math) avec chiffres et lettres droites de Plus
% Jakarta, pour que les formules se fondent dans le texte.
\usepackage[math-style=ISO,bold-style=ISO]{unicode-math}
\setmathfont{FiraMath-Regular.otf}[Path=../../../fonts/]
\setmathfont{PlusJakartaSans-Medium.ttf}[Path=../../../fonts/,range={up/num,up/{Latin,latin},\mathup}]
\usepackage{icomma} % virgule décimale sans espace en mode math
% Caractères mathématiques Unicode absents des polices de texte (Plus Jakarta,
% Archivo Black) : rendus par leur équivalent mathématique, en texte comme en
% formule — sinon ils s'affichent en carré vide (ex. « Arithmétique dans ℤ »).
\usepackage{newunicodechar}
\newunicodechar{ℝ}{\ensuremath{\mathbb{R}}}
\newunicodechar{ℤ}{\ensuremath{\mathbb{Z}}}
\newunicodechar{ℂ}{\ensuremath{\mathbb{C}}}
\newunicodechar{ℕ}{\ensuremath{\mathbb{N}}}
\newunicodechar{ℚ}{\ensuremath{\mathbb{Q}}}
\newunicodechar{𝔻}{\ensuremath{\mathbb{D}}}
\newunicodechar{α}{\ensuremath{\alpha}}
\newunicodechar{β}{\ensuremath{\beta}}
\newunicodechar{γ}{\ensuremath{\gamma}}
\newunicodechar{δ}{\ensuremath{\delta}}
\newunicodechar{ε}{\ensuremath{\varepsilon}}
\newunicodechar{θ}{\ensuremath{\theta}}
\newunicodechar{λ}{\ensuremath{\lambda}}
\newunicodechar{μ}{\ensuremath{\mu}}
\newunicodechar{σ}{\ensuremath{\sigma}}
\newunicodechar{τ}{\ensuremath{\tau}}
\newunicodechar{φ}{\ensuremath{\varphi}}
\newunicodechar{ψ}{\ensuremath{\psi}}
\newunicodechar{ω}{\ensuremath{\omega}}
\newunicodechar{Σ}{\ensuremath{\Sigma}}
% majuscules produites par \MakeUppercase dans les titres (α-aminés -> Α)
\newunicodechar{Α}{\ensuremath{\alpha}}
\newunicodechar{Β}{\ensuremath{\beta}}
\newunicodechar{Γ}{\ensuremath{\Gamma}}
\newunicodechar{Δ}{\ensuremath{\Delta}}
\newunicodechar{Ε}{\ensuremath{\varepsilon}}
\newunicodechar{Θ}{\ensuremath{\Theta}}
\newunicodechar{Λ}{\ensuremath{\Lambda}}
\newunicodechar{Μ}{\ensuremath{\mu}}
\newunicodechar{Τ}{\ensuremath{\tau}}
\newunicodechar{Φ}{\ensuremath{\Phi}}
\newunicodechar{Ψ}{\ensuremath{\Psi}}
\newunicodechar{Ω}{\ensuremath{\Omega}}
\newunicodechar{↦}{\ensuremath{\mapsto}}
\newunicodechar{∈}{\ensuremath{\in}}
\newunicodechar{∉}{\ensuremath{\notin}}
\newunicodechar{∧}{\ensuremath{\wedge}}
\newunicodechar{⇔}{\ensuremath{\Leftrightarrow}}
\newunicodechar{⟺}{\ensuremath{\Longleftrightarrow}}
\newunicodechar{⇒}{\ensuremath{\Rightarrow}}
\newunicodechar{⟹}{\ensuremath{\Longrightarrow}}
\newunicodechar{∘}{\ensuremath{\circ}}
\newunicodechar{∪}{\ensuremath{\cup}}
\newunicodechar{∩}{\ensuremath{\cap}}
\newunicodechar{≡}{\ensuremath{\equiv}}
\newunicodechar{⊂}{\ensuremath{\subset}}
\newunicodechar{⊥}{\ensuremath{\perp}}
\newunicodechar{∃}{\ensuremath{\exists}}
\newunicodechar{∀}{\ensuremath{\forall}}
\newunicodechar{∛}{\ensuremath{\sqrt[3]{\,}}}
\newunicodechar{∜}{\ensuremath{\sqrt[4]{\,}}}
\newunicodechar{⃗}{\ensuremath{{}^{\to}}}
\newunicodechar{ⁿ}{\ifmmode^{n}\else\textsuperscript{\ensuremath{n}}\fi}
\newunicodechar{ˣ}{\ifmmode^{x}\else\textsuperscript{\ensuremath{x}}\fi}
\newunicodechar{ᵇ}{\ifmmode^{b}\else\textsuperscript{\ensuremath{b}}\fi}
\newunicodechar{ʳ}{\ifmmode^{r}\else\textsuperscript{\ensuremath{r}}\fi}
\newunicodechar{ᵘ}{\ifmmode^{u}\else\textsuperscript{\ensuremath{u}}\fi}
\newunicodechar{ʸ}{\ifmmode^{y}\else\textsuperscript{\ensuremath{y}}\fi}
\newunicodechar{ᵃ}{\ifmmode^{a}\else\textsuperscript{\ensuremath{a}}\fi}
\newunicodechar{ᵉ}{\ifmmode^{e}\else\textsuperscript{\ensuremath{e}}\fi}
\newunicodechar{ᵏ}{\ifmmode^{k}\else\textsuperscript{\ensuremath{k}}\fi}
\newunicodechar{ᵖ}{\ifmmode^{p}\else\textsuperscript{\ensuremath{p}}\fi}
\newunicodechar{ᴺ}{\ifmmode^{N}\else\textsuperscript{\ensuremath{N}}\fi}
\newunicodechar{ᶜ}{\ifmmode^{c}\else\textsuperscript{\ensuremath{c}}\fi}
\newunicodechar{ʰ}{\ifmmode^{h}\else\textsuperscript{\ensuremath{h}}\fi}
\newunicodechar{ᵠ}{\ifmmode^{q}\else\textsuperscript{\ensuremath{q}}\fi}
\newunicodechar{ₙ}{\ifmmode_{n}\else\textsubscript{\ensuremath{n}}\fi}
\newunicodechar{ₐ}{\ifmmode_{a}\else\textsubscript{\ensuremath{a}}\fi}
\newunicodechar{ᵢ}{\ifmmode_{i}\else\textsubscript{\ensuremath{i}}\fi}
\newunicodechar{ᵣ}{\ifmmode_{r}\else\textsubscript{\ensuremath{r}}\fi}
\newunicodechar{ₖ}{\ifmmode_{k}\else\textsubscript{\ensuremath{k}}\fi}
\newunicodechar{ᵦ}{\ifmmode_{\beta}\else\textsubscript{\ensuremath{\beta}}\fi}
\newunicodechar{ₓ}{\ifmmode_{x}\else\textsubscript{\ensuremath{x}}\fi}
\newfontfamily\popdisplay{ArchivoBlack-Regular.ttf}[Path=../../../fonts/]
\newfontfamily\poplogo{SpaceGrotesk-Bold.ttf}[Path=../../../fonts/]
\newfontfamily\popsemi{PlusJakartaSans-SemiBold.ttf}[Path=../../../fonts/]
\newfontfamily\popxbold{PlusJakartaSans-ExtraBold.ttf}[Path=../../../fonts/]
\newfontfamily\popxboldls{PlusJakartaSans-ExtraBold.ttf}[Path=../../../fonts/,LetterSpace=3.0]

\setlist[enumerate]{leftmargin=1.7em,itemsep=5pt,topsep=4pt}
\setlist[itemize]{leftmargin=1.7em,itemsep=4pt,topsep=4pt,label={\color{poporange}\textbullet}}
\setlength{\parindent}{0pt}
\setlength{\parskip}{5pt}
\renewcommand{\arraystretch}{1.25}

% pgfplots : style Pop par défaut
\pgfplotsset{
  every axis/.append style={axis line style={popink,line width=1pt},tick style={popink},
    grid style={popline,line width=0.5pt},label style={font=\small},tick label style={font=\footnotesize}},
  popcurve/.style={poporange,line width=1.8pt,no markers,smooth},
}
% Même style disponible dans un \draw TikZ simple (hors pgfplots)
\tikzset{popcurve/.style={poporange,line width=1.8pt}}

% ============================================================
% Pill / squiggle
% ============================================================
\newcommand{\poppill}[3]{%
  \tikz[baseline=-0.55ex]\node[draw=popink,line width=1.2pt,rounded corners=2.6pt,%
    fill=#2,inner xsep=6.5pt,inner ysep=3pt]{%
    \popxboldls\fontsize{7}{7}\selectfont\color{#3}\MakeUppercase{#1}};%
}
\newcommand{\popsquiggle}[1]{%
  \tikz[baseline=-0.4ex]{
    \draw[#1,line width=2.6pt,line cap=round]
      plot [smooth] coordinates {(0,0.09) (0.34,0.01) (0.68,0.21) (1.02,0.01) (1.36,0.10)};
  }%
}
% Mot-clé surligné (marqueur orange sous le texte)
\newcommand{\cle}[1]{{\popxbold\color{popdark}#1}}

% ============================================================
% En-tête / pied
% ============================================================
\pagestyle{fancy}
\fancyhf{}
\renewcommand{\headrulewidth}{0pt}
\renewcommand{\footrulewidth}{0pt}
\lhead{\raisebox{-0.15ex}{\poplogo\fontsize{13}{13}\selectfont\color{popink}OnBuch\color{poporange}+}}
\rhead{\poppill{\DOCMATIERE\ \textbullet\ \DOCNIVEAU\ \textbullet\ Leçon \DOCLECON}{white}{popink}}
\lfoot{\popxbold\fontsize{7.6}{7.6}\selectfont\color{popmuted}\MakeUppercase{Cours signé OnBuch+}\ \textbullet\ {\popsemi\DOCTITRE}}
\rfoot{\tikz[baseline=-0.6ex]\node[rounded corners=8.5pt,fill=popink,minimum height=17pt,minimum width=17pt,inner xsep=5pt,inner sep=0]{\popxbold\fontsize{8}{8}\selectfont\color{popcream}\thepage\,/\,\pageref*{LastPage}};}

% ============================================================
% Titres
% ============================================================
\newcommand{\chaptitre}[2]{%
  \begin{tcolorbox}[enhanced,colback=poporange,colframe=popink,boxrule=2.6pt,arc=11pt,
      drop shadow=popink,left=15pt,right=15pt,top=12pt,bottom=11pt,before skip=3pt,after skip=18pt]
    \poppill{#2}{popink}{popcream}\par\vspace{9pt}
    {\raggedright\hyphenpenalty=10000\exhyphenpenalty=10000\popdisplay\fontsize{22}{24}\selectfont\color{popcream}\MakeUppercase{#1}\par}
  \end{tcolorbox}
}
% Section numérotée : pastille orange avec le numéro + titre Archivo
\newcounter{popsec}
\newcommand{\coursec}[1]{%
  \stepcounter{popsec}\setcounter{popsub}{0}%
  \par\vspace{16pt}\noindent
  \begin{minipage}{\linewidth}
  \tikz[baseline=-0.7ex]\node[draw=popink,line width=1.6pt,fill=poporange,rounded corners=4pt,minimum size=22pt,inner sep=0]{\popdisplay\fontsize{12}{12}\selectfont\color{popcream}\thepopsec};%
  \hspace{8pt}{\popdisplay\fontsize{15}{17}\selectfont\color{popink}\MakeUppercase{#1}}\par
  \vspace{4pt}\noindent{\color{poporange}\rule{36pt}{3.6pt}}
  \end{minipage}\par\nopagebreak\vspace{8pt}
}
\newcounter{popsub}
\newcommand{\courssub}[1]{%
  \stepcounter{popsub}%
  \par\vspace{9pt}\noindent{\popxbold\fontsize{11.5}{13}\selectfont\color{popink}{\color{poporange}\thepopsec.\thepopsub}\hspace{6pt}#1}\par\nopagebreak\vspace{4pt}
}

% ============================================================
% Boîtes pédagogiques — même recette : fond blanc, bordure épaisse colorée,
% ombre dure encre, badge plein. Titre optionnel [#1].
% ============================================================
\tcbset{popbase/.style={breakable,enhanced,colback=white,boxrule=2.3pt,arc=9pt,
  shadow={3pt}{-3pt}{0pt}{popink},left=14pt,right=14pt,top=11pt,bottom=11pt,before skip=11pt,after skip=15pt,
  fontupper=\popsemi\fontsize{10.6}{14.5}\selectfont\color{popink},
  fontlower=\popsemi\fontsize{10.6}{14.5}\selectfont\color{popink}}}
% Coche dessinée (le glyphe ✓ n'existe pas dans les polices du document)
\newcommand{\popcheck}[1]{\tikz[baseline=-0.5ex]\draw[#1,line width=1.6pt,line cap=round,line join=round](-0.13,0.0)--(-0.04,-0.09)--(0.13,0.1);}
\newcommand{\popboxhead}[3]{\poppill{#1}{#2}{white}\ifx\relax#3\relax\else\hspace{6pt}{\popxbold\fontsize{10.6}{12}\selectfont\color{popink}#3}\fi\par\vspace{7pt}}

\newtcolorbox{aretenir}[1][]{popbase,colframe=popgreen,before upper={\popboxhead{À retenir}{popgreen}{#1}}}
\newtcolorbox{exemplebox}[1][]{popbase,colframe=poporange,before upper={\popboxhead{Exemple}{poporange}{#1}}}
\newtcolorbox{definition}[1][]{popbase,colframe=popblue,before upper={\popboxhead{Définition}{popblue}{#1}}}
\newtcolorbox{propriete}[1][]{popbase,colframe=poppurple,before upper={\popboxhead{Propriété}{poppurple}{#1}}}
\newtcolorbox{methode}[1][]{popbase,colframe=popgold,before upper={\popboxhead{Méthode}{popgold}{#1}}}
\newtcolorbox{attention}[1][]{popbase,colframe=poppink,before upper={\popboxhead{Attention}{poppink}{#1}}}
\newtcolorbox{experience}[1][]{popbase,colframe=popblue,colback=popblueL,before upper={\popboxhead{Expérience}{popblue}{#1}}}
% « Le savais-tu ? » : carte violette pleine (contexte, culture, Cameroun)
\newtcolorbox{savaistu}[1][]{popbase,colback=poppurple,colframe=popink,
  fontupper=\popsemi\fontsize{10.4}{14.2}\selectfont\color{white},
  before upper={\poppill{Le savais-tu\,?}{white}{poppurple}\ifx\relax#1\relax\else\hspace{6pt}{\popxbold\color{white}#1}\fi\par\vspace{7pt}}}
% Exercice résolu : énoncé en haut, \tcblower puis la solution sur fond crème
\newcounter{popexo}\newcounter{popexr}
\newtcolorbox{exoresolu}[1][]{popbase,colframe=poporange,colbacklower=popcream,
  segmentation style={popink,line width=1.2pt,dashed},
  before upper={\stepcounter{popexr}\popboxhead{Exercice résolu \thepopexr}{poporange}{#1}},
  before lower={\poppill{Solution}{popink}{popcream}\par\vspace{6pt}}}
% Exercice (non résolu en place) — la correction est dans la partie Corrigés
\newtcolorbox{exercice}[1][]{popbase,colframe=popink,
  before upper={\stepcounter{popexo}\popboxhead{Exercice \thepopexo}{popink}{#1}}}
% Corrigé
\newtcolorbox{corrige}[1][]{popbase,colframe=popgreen,colback=popgreenL,
  before upper={\popboxhead{Corrigé}{popgreen}{#1}}}
% Objectifs / prérequis / bilan
\newtcolorbox{objectifs}{popbase,colframe=popink,colback=poporangeL,
  before upper={\popboxhead{Objectifs de la leçon}{popink}{}}}
\newtcolorbox{prerequis}{popbase,colframe=popink,colback=popblueL,
  before upper={\popboxhead{Prérequis}{popblue}{}}}
\newtcolorbox{bilan}{popbase,colframe=popink,colback=white,boxrule=2.8pt,
  before upper={\popboxhead{Fiche bilan}{poporange}{L'essentiel à connaître pour l'examen}}}
% Figure encadrée avec légende
\newtcolorbox{popfigure}[1][]{enhanced,breakable=false,colback=white,colframe=popline,boxrule=1.4pt,arc=8pt,
  left=10pt,right=10pt,top=10pt,bottom=8pt,before skip=10pt,after skip=12pt,halign=center,
  lower separated=false,halign lower=center,
  fontlower=\popsemi\fontsize{9}{11.5}\selectfont\color{popmuted},
  #1}
\newcommand{\legende}[1]{\par\vspace{6pt}{\popsemi\fontsize{9}{11.5}\selectfont\color{popmuted}#1}}

% ============================================================
% Signature OnBuch+
% ============================================================
\newcommand{\poplogoinline}[1]{{\poplogo\fontsize{#1}{#1}\selectfont\color{popink}OnBuch\color{poporange}+}}
\newcommand{\signatureonbuch}{%
  \begin{tcolorbox}[enhanced,colback=popink,colframe=popink,boxrule=2.2pt,arc=10pt,
      drop shadow=poporange,left=14pt,right=14pt,top=10pt,bottom=10pt,before skip=0pt,after skip=0pt]
    \tikz[baseline=-0.7ex]\node[circle,fill=popgreen,draw=popcream,line width=1.3pt,minimum size=22pt,inner sep=0]{\popcheck{white}};%
    \hspace{9pt}\begin{minipage}[c]{0.62\linewidth}
      {\popxbold\fontsize{12}{13}\selectfont\color{popcream}Signé\ \textbullet\ {\poplogo OnBuch\color{poporange}+}}\par\vspace{2pt}
      {\raggedright\popsemi\fontsize{8}{10.5}\selectfont\color{popline}\MakeUppercase{Équipe pédagogique OnBuch+}\\\MakeUppercase{Conforme au programme officiel MINESEC}\par}
    \end{minipage}\hfill
    {\poplogo\fontsize{22}{22}\selectfont\color{popcream}OnBuch\color{poporange}+}
  \end{tcolorbox}%
}

% ============================================================
% Page de garde
% #1 titre · #2 matière · #3 niveau · #4 module · #5 n° leçon · #6 durée ·
% #7 description / objectifs (texte ou itemize)
% ============================================================
\newcommand{\pagegarde}[7]{%
  \thispagestyle{empty}
  \newgeometry{a4paper,margin=1.7cm,top=1.6cm,bottom=1.4cm}%
  \begin{tikzpicture}[remember picture,overlay]
    \fill[poporange!18] ([xshift=-3.2cm,yshift=-3.6cm]current page.north east) circle (4.6cm);
    \fill[poppurple!14] ([xshift=2.4cm,yshift=7.6cm]current page.south west) circle (3.8cm);
  \end{tikzpicture}%
  {\poplogo\fontsize{30}{30}\selectfont\color{popink}OnBuch\color{poporange}+}\hfill
  \raisebox{0.6ex}{\poppill{Cours signé \textbullet\ Premium}{popink}{popcream}}\par
  \vspace{1pt}\popsquiggle{poppurple}\par
  \vspace{8pt}
  \begin{tcolorbox}[enhanced,colback=poporange,colframe=popink,boxrule=2.8pt,arc=13pt,
      drop shadow=popink,left=18pt,right=18pt,top=12pt,bottom=13pt,after skip=10pt]
    \poppill{#2 \textbullet\ #3}{popink}{popcream}\hspace{5pt}\poppill{#4}{white}{popink}\par\vspace{8pt}
    {\popdisplay\fontsize{14}{14}\selectfont\color{popink}LEÇON #5}\par\vspace{4pt}
    {\raggedright\hyphenpenalty=10000\exhyphenpenalty=10000\popdisplay\fontsize{25}{27}\selectfont\color{popcream}\MakeUppercase{#1}\par}
    \vspace{8pt}
    {\popxbold\fontsize{9.5}{9.5}\selectfont\color{popink}\MakeUppercase{Durée conseillée : #6}}
  \end{tcolorbox}
  \begin{tcolorbox}[enhanced,colback=white,colframe=popink,boxrule=2.2pt,arc=10pt,
      drop shadow=popink,left=16pt,right=16pt,top=10pt,bottom=10pt,after skip=8pt]
    \poppill{Dans cette leçon}{poppurple}{white}\par\vspace{5pt}
    \popsemi\fontsize{9.3}{12.6}\selectfont\color{popink}#7
  \end{tcolorbox}
  \noindent\begin{minipage}[t]{0.487\linewidth}
    \begin{tcolorbox}[enhanced,colback=popgreen,colframe=popink,boxrule=2.2pt,arc=10pt,
        drop shadow=popink,left=11pt,right=11pt,top=10pt,bottom=10pt,height=3.1cm,valign=center]
      \tcbox[colback=white,colframe=popink,boxrule=1.3pt,arc=5pt,boxsep=0pt,nobeforeafter,
        left=3pt,right=3pt,top=3pt,bottom=3pt,valign=center]{\qrcode[height=1.9cm]{https://play.google.com/store/apps/details?id=com.luvvix.onbuch&hl=fr}}%
      \hspace{8pt}\begin{minipage}[c]{0.5\linewidth}
        {\popdisplay\fontsize{11}{12.5}\selectfont\color{white}\MakeUppercase{Télécharge OnBuch}}\par\vspace{4pt}
        {\raggedright\popsemi\fontsize{8.4}{11}\selectfont\color{white}Gratuit \textbullet\ Google Play\\Tuteur IA, annales, concours, résultats\par}
      \end{minipage}
    \end{tcolorbox}
  \end{minipage}\hfill
  \begin{minipage}[t]{0.487\linewidth}
    \begin{tcolorbox}[enhanced,colback=poppurple,colframe=popink,boxrule=2.2pt,arc=10pt,
        drop shadow=popink,left=11pt,right=11pt,top=10pt,bottom=10pt,height=3.1cm,valign=center]
      \tikz[baseline=-0.7ex]\node[circle,fill=white,draw=popink,line width=1.3pt,minimum size=1.6cm,inner sep=0]{\popdisplay\fontsize{24}{24}\selectfont\color{poppurple}+};%
      \hspace{8pt}\begin{minipage}[c]{0.6\linewidth}
        {\popdisplay\fontsize{11}{12.5}\selectfont\color{white}\MakeUppercase{Passe à OnBuch+}}\par\vspace{4pt}
        {\raggedright\popsemi\fontsize{8.4}{11}\selectfont\color{white}Tous les cours signés, Tuteur IA illimité, fiches PDF — onglet OnBuch+ de l'app\par}
      \end{minipage}
    \end{tcolorbox}
  \end{minipage}
  \vspace{8pt}
  \popcontenu
  \vfill
  \signatureonbuch
  \restoregeometry
  \newpage
}

% Rangée « Ce cours contient » (page de garde)
\newcommand{\poptile}[3]{% #1 couleur #2 gros texte #3 légende
  \begin{tcolorbox}[enhanced,colback=#1,colframe=popink,boxrule=2pt,arc=9pt,shadow={2.5pt}{-2.5pt}{0pt}{popink},
      left=6pt,right=6pt,top=9pt,bottom=9pt,halign=center,valign=center,height=1.75cm,width=\linewidth,nobeforeafter]
    {\popdisplay\fontsize{12.5}{13}\selectfont\color{white}\MakeUppercase{#2}}\par\vspace{3pt}
    {\centering\popsemi\fontsize{7.8}{9.5}\selectfont\color{white}#3\par}
  \end{tcolorbox}}
\newcommand{\popcontenu}{%
  \par\noindent{\popxboldls\fontsize{7.5}{7.5}\selectfont\color{popmuted}CE COURS SIGNÉ CONTIENT}\par\vspace{7pt}
  \noindent\begin{minipage}[t]{0.235\linewidth}\poptile{popblue}{Cours}{complet, illustré, pas à pas}\end{minipage}\hfill
  \begin{minipage}[t]{0.235\linewidth}\poptile{poporange}{Méthodes}{et exercices résolus}\end{minipage}\hfill
  \begin{minipage}[t]{0.235\linewidth}\poptile{poppink}{Exercices}{type examen + corrigés}\end{minipage}\hfill
  \begin{minipage}[t]{0.235\linewidth}\poptile{popgreen}{Fiche bilan}{l'essentiel à retenir}\end{minipage}\par}

% Sommaire
\newtcolorbox{sommaire}{enhanced,colback=white,colframe=popink,boxrule=2.2pt,arc=10pt,
  drop shadow=popink,left=18pt,right=18pt,top=13pt,bottom=13pt,before skip=6pt,after skip=20pt,
  before upper={\poppill{Sommaire}{poppurple}{white}\par\vspace{9pt}\popsemi\fontsize{10.6}{14.5}\selectfont\color{popink}}}

% Signature de fin de cours — poussée tout en bas de la dernière page
\newcommand{\findecours}{%
  \par\vfill\nopagebreak
  \begin{center}\popsquiggle{poporange}\end{center}\vspace{4pt}
  \signatureonbuch
}
\AtBeginDocument{\def\llbracket{\mathopen{[\mkern-3.5mu[}}\def\rrbracket{\mathclose{]\mkern-3.5mu]}}} % intervalles d'entiers (glyphe ⟦ absent de la police)
% Glyphes absents de Fira Math : redéfinis à partir de glyphes disponibles
\AtBeginDocument{%
  \def\setminus{\mathbin{\backslash}}\let\smallsetminus\setminus
  \def\vdots{\mathord{\vbox{\baselineskip3.2pt\lineskiplimit0pt\kern4pt\hbox{.}\hbox{.}\hbox{.}}}}%
  \def\ddots{\mathinner{\mkern1mu\raise6.5pt\hbox{.}\mkern2mu\raise3.6pt\hbox{.}\mkern2mu\raise0.7pt\hbox{.}\mkern1mu}}%
  \def\triangle{\mathord{\scalebox{0.9}{$\symup{\Delta}$}}}%
  \def\nabla{\mathord{\raisebox{\depth}{\scalebox{1}[-1]{$\symup{\Delta}$}}}}%
  \def\checkmark{\popcheck{popdark}}%
  \def\star{\mathbin{\ast}}\def\bigstar{\mathord{\ast}}%
  \def\diamond{\mathbin{\scalebox{0.6}{$\Diamond$}}}%
  \def\bigcup{\mathop{\mathchoice{\raisebox{-0.3em}{\scalebox{1.7}{$\cup$}}}{\scalebox{1.25}{$\cup$}}{\cup}{\cup}}\displaylimits}%
  \def\bigcap{\mathop{\mathchoice{\raisebox{-0.3em}{\scalebox{1.7}{$\cap$}}}{\scalebox{1.25}{$\cap$}}{\cap}{\cap}}\displaylimits}%
  \def\lesseqqgtr{\mathrel{\lessgtr}}\def\bigoplus{\mathop{\oplus}\displaylimits}%
}
\providecommand{\ssi}{si et seulement si} % abréviation fréquente
\AtBeginDocument{\providecommand{\wideparen}{\overparen}} % arc de cercle

%% FIGFIX-BEGIN v5 — correctifs globaux des figures (docs/onbuch-generation/tools/figfix)
\usepackage{adjustbox}\usepackage{etoolbox}\tcbuselibrary{raster}
% toute figure TikZ/pgfplots plus large que la ligne est réduite à la largeur disponible
\BeforeBeginEnvironment{tikzpicture}{\begin{adjustbox}{max width=\linewidth}}
\AfterEndEnvironment{tikzpicture}{\end{adjustbox}}
% légendes pgfplots lisibles par-dessus les courbes
\pgfplotsset{every axis/.append style={legend style={fill=white,fill opacity=0.92,text opacity=1,draw=black!15,inner sep=3pt}}}
% étiquettes d'arêtes (syntaxe edge["..."]) avec fond blanc
\tikzset{every edge quotes/.append style={fill=white,inner sep=1.2pt}}
%% FIGFIX-END
\begin{document}

\pagegarde{Suites numériques}{Mathématiques}{Tle Tech}{Mathématiques – classe de Terminale technique}{N06}{8 séances (fiche de progression harmonisée)}{Cette leçon structure l'étude rigoureuse des suites numériques : raisonnement par récurrence, monotonie, bornitude, convergence et théorèmes fondamentaux. Elle prépare l'élève à modéliser et résoudre des problèmes concrets par l'outil suite récurrente, compétence clé pour le Probatoire et la poursuite d'études techniques.
\vspace{4pt}
\begin{multicols}{2}\small
\begin{itemize}
\item À la fin de cette leçon, je sais construire et rédiger une démonstration par récurrence sur ℕ pour prouver une propriété héréditaire.
\item Je sais déterminer la monotonie d'une suite (croissante, décroissante, constante) par la différence, le quotient ou l'étude de la fonction f définissant la récurrence.
\item Je sais majorer ou minorer une suite et en déduire sa bornitude.
\item J'applique le théorème de convergence des suites monotones bornées pour affirmer l'existence d'une limite finie.
\item Je calcule la limite d'une suite convergente définie par U\_{n+1} = f(U\_n) en résolvant l'équation f(l) = l.
\item J'utilise les opérations sur les limites (somme, produit, quotient, composition) pour déterminer des limites de suites complexes.
\end{itemize}\end{multicols}}

\begin{sommaire}
\begin{multicols}{2}
\begin{enumerate}
\item 1. Raisonnement par récurrence : principe et rédaction
\item 2. Monotonie d'une suite : méthodes d'étude
\item 3. Bornitude et Convergence : Théorèmes fondamentaux
\item 4. Calcul de limites : Opérations et suites de référence
\item 5. Modélisation et applications concrètes (Vie courante \& Technique)
\item 6. Synthèse et entraînement type Probatoire
\item Activité d'intégration
\item Méthodes et exercices résolus
\item Exercices
\item Corrigés des exercices
\item Fiche bilan
\end{enumerate}
\end{multicols}
\end{sommaire}

\begin{objectifs}
\begin{itemize}
\item À la fin de cette leçon, je sais construire et rédiger une démonstration par récurrence sur ℕ pour prouver une propriété héréditaire.
\item Je sais déterminer la monotonie d'une suite (croissante, décroissante, constante) par la différence, le quotient ou l'étude de la fonction f définissant la récurrence.
\item Je sais majorer ou minorer une suite et en déduire sa bornitude.
\item J'applique le théorème de convergence des suites monotones bornées pour affirmer l'existence d'une limite finie.
\item Je calcule la limite d'une suite convergente définie par U\_{n+1} = f(U\_n) en résolvant l'équation f(l) = l.
\item J'utilise les opérations sur les limites (somme, produit, quotient, composition) pour déterminer des limites de suites complexes.
\item Je modélise un problème concret (épargne, population, stock, pharmacocinétique) par une suite récurrente et j'interprète le résultat dans le contexte.
\item Je rédige une réponse complète, justifiée et structurée, adaptée aux attendus de l'épreuve du Probatoire.
\end{itemize}
\end{objectifs}

\begin{prerequis}
\begin{itemize}
\item Définition d'une suite numérique (explicite et par récurrence) - Première/Debut Terminale.
\item Notion de limite d'une suite (définition intuitive, suites de référence : 1/n, q\^{}n) - Première.
\item Propriétés des fonctions de référence (affine, racine carrée, inverse, exponentielle, logarithme) - Première/Terminale.
\item Calcul algébrique : résolution d'équations et d'inéquations du 1er et 2nd degré, manipulation de fractions.
\item Logique élémentaire : implication, équivalence, contraposée, quantification (pour tout, il existe).
\end{itemize}
\end{prerequis}

\newpage

\coursec{Raisonnement par récurrence : principe et rédaction}

\textbf{Situation d'accroche.} Monsieur Kamga, gérant d'une petite entreprise de transformation de manioc à Bafoussam, observe que sa production mensuelle augmente de 5\,\% chaque mois grâce à l'acquisition d'une nouvelle machine, mais il doit aussi rembourser un prêt qui réduit son stock net de 2 tonnes fixes par mois. Partant d'un stock initial de 100 tonnes, il se demande si sa production finira par \emph{exploser} ou se \emph{stabiliser}, et au bout de combien de mois il dépassera le seuil de rentabilité de 200 tonnes.

Si l'on note $u_n$ le stock (en tonnes) au bout de $n$ mois, la situation se traduit par :
\[ u_0 = 100 \qquad \text{et} \qquad u_{n+1} = 1{,}05\,u_n - 2 \quad \text{pour tout } n \in \mathbb{N}. \]
C'est une \cle{suite définie par récurrence}. Pour répondre à Monsieur Kamga, il faudra établir des propriétés valables \emph{pour tous} les mois : par exemple \og le stock reste toujours supérieur ou égal à 100 tonnes \fg{} ou \og le stock est croissant \fg. Comment prouver une propriété pour une infinité de valeurs de $n$ sans pouvoir toutes les vérifier une par une ?

\textbf{Problématique.} Comment démontrer rigoureusement qu'une propriété $P(n)$ est vraie pour \emph{tous} les entiers naturels, en ne faisant qu'un nombre fini de vérifications ?

\courssub{Le principe de récurrence : intuition et énoncé}

\textbf{Intuition : l'effet domino.} Imagine une file infinie de dominos debout, numérotés $0, 1, 2, 3, \ldots$ Pour être certain que \emph{tous} les dominos tomberont, il suffit de deux choses :
\begin{itemize}
\item faire tomber le \textbf{premier} domino (le domino numéro 0) ;
\item s'assurer que chaque domino, en tombant, fait tomber le \textbf{suivant}.
\end{itemize}
Si ces deux conditions sont réunies, alors le domino 0 tombe, donc le 1 tombe, donc le 2 tombe, et ainsi de suite : \emph{tous} tombent. C'est exactement l'idée du raisonnement par récurrence.

\begin{definition}[Principe de récurrence (énoncé formel)]
Soit $P(n)$ une propriété dépendant d'un entier naturel $n$. Si les deux conditions suivantes sont satisfaites :
\begin{itemize}
\item \textbf{Initialisation :} $P(0)$ est vraie ;
\item \textbf{Hérédité :} pour tout entier $n \in \mathbb{N}$, si $P(n)$ est vraie alors $P(n+1)$ est vraie ;
\end{itemize}
alors la propriété $P(n)$ est vraie pour \textbf{tout} entier naturel $n \in \mathbb{N}$.

Plus généralement, si $P(n_0)$ est vraie et si l'hérédité est vraie pour tout $n \geq n_0$, alors $P(n)$ est vraie pour tout $n \geq n_0$ (\cle{initialisation décalée}).
\end{definition}

Ce principe est un \textbf{axiome} des mathématiques (lié à la construction de $\mathbb{N}$ par Peano) : on ne le démontre pas, on l'admet comme règle fondamentale de raisonnement. Il est d'ailleurs équivalent à une propriété remarquable des entiers naturels.

\begin{propriete}[Propriété admise : bon ordre de $\mathbb{N}$]
Toute partie non vide de $\mathbb{N}$ admet un plus petit élément. Cette propriété est \emph{équivalente} au principe de récurrence : admettre l'une revient à admettre l'autre.
\end{propriete}

\begin{experience}[Pourquoi l'équivalence est plausible]
Supposons qu'une propriété $P(n)$ soit initialisée et héréditaire, mais qu'il existe malgré tout un entier pour lequel elle est fausse. L'ensemble $E$ des entiers pour lesquels $P$ est fausse serait non vide : il admettrait donc un plus petit élément $m$. Or $m \neq 0$ (initialisation), donc $m - 1 \in \mathbb{N}$ et $P(m-1)$ est vraie (car $m$ est le plus petit élément de $E$). Par hérédité, $P(m)$ serait vraie : contradiction. Donc $E$ est vide et $P$ est vraie partout.
\end{experience}

La figure suivante visualise la chaîne infinie d'implications qui découle de l'initialisation et de l'hérédité.

\begin{popfigure}
\begin{tikzpicture}[node distance=0.55cm and 1.1cm,
  boite/.style={draw=popblue, fill=popblueL, rounded corners=2pt, thick, minimum height=0.8cm, inner sep=4pt, font=\small},
  vrai/.style={draw=popgreen, fill=popgreenL, rounded corners=2pt, thick, minimum height=0.8cm, inner sep=4pt, font=\small},
  fleche/.style={-{Stealth[length=2.5mm]}, thick, poporange}]
\node[boite] (p0) {$P(0)$ vraie};
\node[boite, right=of p0] (p1) {$P(1)$ vraie};
\node[boite, right=of p1] (p2) {$P(2)$ vraie};
\node[boite, right=of p2] (p3) {$P(3)$ vraie};
\node[right=of p3, font=\Large] (pts) {$\cdots$};
\node[vrai, right=of pts] (pn) {$\forall n,\; P(n)$ vraie};
\draw[fleche] (p0) -- node[above, font=\scriptsize]{hérédité} (p1);
\draw[fleche] (p1) -- node[above, font=\scriptsize]{hérédité} (p2);
\draw[fleche] (p2) -- node[above, font=\scriptsize]{hérédité} (p3);
\draw[fleche] (p3) -- (pts);
\draw[fleche] (pts) -- (pn);
\node[below=0.35cm of p0, font=\scriptsize, text=popdark]{initialisation};
\draw[fleche, popgreen] (p0.south) -- ++(0,-0.28);
\end{tikzpicture}
\legende{Chaîne d'implications du raisonnement par récurrence : une fois le premier domino tombé (initialisation), l'hérédité fait tomber tous les suivants.}
\end{popfigure}

\courssub{La rédaction type d'une démonstration par récurrence}

\begin{methode}[Rédiger une démonstration par récurrence]
Pour montrer qu'une propriété $P(n)$ est vraie pour tout $n \geq n_0$ :
\begin{enumerate}
\item \textbf{Énoncer clairement la propriété} $P(n)$ à démontrer (par exemple : \og Soit $P(n)$ : $u_n \geq 100$ \fg).
\item \textbf{Initialisation :} vérifier que $P(n_0)$ est vraie, par un calcul direct. Rédiger : \og Pour $n = n_0$, \ldots donc $P(n_0)$ est vraie. \fg
\item \textbf{Hérédité :} soit $n \geq n_0$ un entier fixé. \textbf{Supposer} que $P(n)$ est vraie (c'est l'\cle{hypothèse de récurrence}) et \textbf{démontrer}, à l'aide de cette hypothèse, que $P(n+1)$ est vraie. Rédiger : \og Supposons $P(n)$ vraie, c'est-à-dire \ldots Montrons que $P(n+1)$ est vraie. \fg
\item \textbf{Conclusion :} invoquer le principe de récurrence. Rédiger : \og Par le principe de récurrence, $P(n)$ est vraie pour tout entier $n \geq n_0$. \fg
\end{enumerate}
\end{methode}

\begin{popfigure}
\begin{tikzpicture}[node distance=0.5cm,
  etape/.style={draw=popdark, thick, rounded corners=3pt, fill=popcream, text width=3.3cm, minimum height=1.5cm, align=center, font=\small},
  fleche/.style={-{Stealth[length=3mm]}, very thick, poporange}]
\node[etape, align=center] (init){\textcolor{popblue}{\textbf{1. Initialisation}}\\[2pt] Vérifier $P(n_0)$ par un calcul direct};
\node[etape, right=1.1cm of init, align=center] (hered){\textcolor{popblue}{\textbf{2. Hérédité}}\\[2pt] Supposer $P(n)$, prouver $P(n+1)$};
\node[etape, right=1.1cm of hered, align=center] (concl){\textcolor{popblue}{\textbf{3. Conclusion}}\\[2pt] $P(n)$ vraie pour tout $n \geq n_0$};
\draw[fleche] (init) -- (hered);
\draw[fleche] (hered) -- (concl);
\end{tikzpicture}
\legende{Les trois étapes obligatoires de la rédaction d'une démonstration par récurrence.}
\end{popfigure}

\begin{attention}[Erreurs fréquentes à éviter absolument]
\begin{itemize}
\item \textbf{Oublier l'initialisation.} L'hérédité seule ne suffit pas ! Exemple : la propriété $P(n)$ : \og $3^n$ est pair \fg{} est héréditaire (si $3^n$ est pair, alors $3^{n+1} = 3 \times 3^n$ est pair), mais $P(0)$ est fausse ($3^0 = 1$ est impair). La propriété est donc fausse pour tout $n$.
\item \textbf{Confondre hypothèse et conclusion.} Dans l'hérédité, on \emph{suppose} $P(n)$ et on \emph{démontre} $P(n+1)$. Écrire \og supposons que $P(n)$ est vraie pour tout $n$ \fg{} revient à supposer ce qu'on veut prouver : c'est un \cle{raisonnement circulaire}, sanctionné à l'examen.
\item \textbf{Ne pas utiliser l'hypothèse de récurrence.} Si la preuve de $P(n+1)$ n'utilise jamais $P(n)$, ce n'est pas une récurrence (ou la rédaction est fautive).
\item \textbf{Initialiser au mauvais rang.} Si la propriété n'a de sens que pour $n \geq 1$ (par exemple une somme de 1 à $n$), l'initialisation doit se faire en $n = 1$, pas en $n = 0$.
\end{itemize}
\end{attention}

\courssub{Exemples guidés et applications classiques}

\begin{exoresolu}[Somme des premiers entiers]
\textbf{Énoncé.} Démontrer par récurrence que pour tout entier $n \geq 1$ :
\[ \sum_{k=1}^{n} k = 1 + 2 + \cdots + n = \frac{n(n+1)}{2}. \]
\tcblower
\textbf{Solution détaillée.} Soit $P(n)$ la propriété : $\displaystyle\sum_{k=1}^{n} k = \frac{n(n+1)}{2}$.

\textbf{Initialisation.} Pour $n = 1$ : $\displaystyle\sum_{k=1}^{1} k = 1$ et $\dfrac{1 \times (1+1)}{2} = \dfrac{2}{2} = 1$. Donc $P(1)$ est vraie.

\textbf{Hérédité.} Soit $n \geq 1$ un entier fixé. Supposons $P(n)$ vraie, c'est-à-dire $\displaystyle\sum_{k=1}^{n} k = \frac{n(n+1)}{2}$. Montrons que $P(n+1)$ est vraie. On a :
\begin{align*}
\sum_{k=1}^{n+1} k &= \underbrace{1 + 2 + \cdots + n}_{=\,\frac{n(n+1)}{2} \text{ par hypothèse}} + (n+1) \\
&= \frac{n(n+1)}{2} + (n+1) = \frac{n(n+1) + 2(n+1)}{2} \\
&= \frac{(n+1)(n+2)}{2}.
\end{align*}
C'est exactement $P(n+1)$. Donc l'hérédité est établie.

\textbf{Conclusion.} Par le principe de récurrence, pour tout entier $n \geq 1$, $\displaystyle\sum_{k=1}^{n} k = \frac{n(n+1)}{2}$.
\end{exoresolu}

\begin{exoresolu}[Divisibilité : exemple guidé]
\textbf{Énoncé.} Montrer que pour tout entier $n \in \mathbb{N}^*$, le nombre $3^n - 1$ est divisible par 2.
\tcblower
\textbf{Solution détaillée.} Soit $P(n)$ la propriété : \og il existe un entier $k \in \mathbb{N}$ tel que $3^n - 1 = 2k$ \fg.

\textbf{Initialisation.} Pour $n = 1$ : $3^1 - 1 = 2 = 2 \times 1$. Donc $P(1)$ est vraie (avec $k = 1$).

\textbf{Hérédité.} Soit $n \geq 1$ fixé. Supposons $P(n)$ vraie : il existe $k \in \mathbb{N}$ tel que $3^n - 1 = 2k$, c'est-à-dire $3^n = 2k + 1$. Alors :
\[ 3^{n+1} - 1 = 3 \times 3^n - 1 = 3(2k+1) - 1 = 6k + 3 - 1 = 6k + 2 = 2(3k+1). \]
Or $3k + 1$ est un entier naturel, donc $3^{n+1} - 1$ est divisible par 2 : $P(n+1)$ est vraie.

\textbf{Conclusion.} Par le principe de récurrence, $3^n - 1$ est divisible par 2 pour tout $n \in \mathbb{N}^*$.
\end{exoresolu}

\begin{exemplebox}[Inégalité avec initialisation décalée]
Montrons que pour tout $n \geq 3$, on a $2^n \geq n + 5$.

\textbf{Initialisation.} Pour $n = 3$ : $2^3 = 8$ et $3 + 5 = 8$, donc $2^3 \geq 3 + 5$ : la propriété est vraie au rang $3$ (elle est d'ailleurs fausse pour $n = 0, 1, 2$ : l'initialisation décalée est ici indispensable).

\textbf{Hérédité.} Soit $n \geq 3$ fixé tel que $2^n \geq n + 5$. Alors :
\[ 2^{n+1} = 2 \times 2^n \geq 2(n+5) = 2n + 10 = (n + 6) + (n + 4) \geq n + 6 = (n+1) + 5, \]
car $n + 4 \geq 0$. La propriété est donc héréditaire à partir du rang 3.

\textbf{Conclusion.} Pour tout $n \geq 3$, $2^n \geq n + 5$.
\end{exemplebox}

\begin{exemplebox}[Minoration d'une suite récurrente — retour à Monsieur Kamga]
Considérons la suite du stock : $u_0 = 100$ et $u_{n+1} = 1{,}05\,u_n - 2$. Montrons que pour tout $n \in \mathbb{N}$, $u_n \geq 100$.

\textbf{Initialisation.} $u_0 = 100 \geq 100$ : vrai.

\textbf{Hérédité.} Soit $n \in \mathbb{N}$ fixé tel que $u_n \geq 100$. Alors :
\[ u_{n+1} = 1{,}05\,u_n - 2 \geq 1{,}05 \times 100 - 2 = 105 - 2 = 103 \geq 100. \]
Donc $u_{n+1} \geq 100$ : la propriété est héréditaire.

\textbf{Conclusion.} Pour tout $n \in \mathbb{N}$, $u_n \geq 100$ : le stock de Monsieur Kamga ne descendra jamais sous 100 tonnes. C'est un premier pas vers l'étude complète de sa rentabilité (qui sera poursuivie dans les sections suivantes).
\end{exemplebox}

\courssub{Variantes du raisonnement par récurrence}

\begin{definition}[Récurrence forte]
Pour démontrer que $P(n)$ est vraie pour tout $n \geq n_0$, on peut remplacer l'hérédité classique par l'\textbf{hérédité forte} : pour tout $n \geq n_0$, si $P(k)$ est vraie \emph{pour tous} les entiers $k$ tels que $n_0 \leq k \leq n$, alors $P(n+1)$ est vraie.

On l'utilise lorsque la preuve au rang $n+1$ nécessite la propriété à un rang autre que $n$ (par exemple aux rangs $n$ \emph{et} $n-1$, comme pour les suites définies par $u_{n+2} = f(u_{n+1}, u_n)$). Dans ce cas, on initialise souvent aux \textbf{deux premiers rangs}.
\end{definition}

\begin{exemplebox}[Récurrence d'ordre 2]
Soit $(u_n)$ définie par $u_0 = 1$, $u_1 = 3$ et $u_{n+2} = u_{n+1} + 2u_n$. Montrons que $u_n \leq 3^n$ pour tout $n \in \mathbb{N}$.

\textbf{Initialisation (deux rangs).} $u_0 = 1 \leq 3^0 = 1$ et $u_1 = 3 \leq 3^1 = 3$ : vrai aux rangs 0 et 1.

\textbf{Hérédité forte.} Soit $n \geq 0$ fixé. Supposons la propriété vraie aux rangs $n$ et $n+1$ : $u_n \leq 3^n$ et $u_{n+1} \leq 3^{n+1}$. Alors :
\[ u_{n+2} = u_{n+1} + 2u_n \leq 3^{n+1} + 2 \times 3^n = 3 \times 3^n + 2 \times 3^n = 5 \times 3^n \leq 9 \times 3^n = 3^{n+2}. \]
\textbf{Conclusion.} Par récurrence forte, $u_n \leq 3^n$ pour tout $n \in \mathbb{N}$.
\end{exemplebox}

\begin{aretenir}[L'essentiel de la section]
\begin{itemize}
\item Le \cle{raisonnement par récurrence} repose sur deux piliers : l'\textbf{initialisation} (la propriété est vraie au premier rang) et l'\textbf{hérédité} (si elle est vraie au rang $n$, elle l'est au rang $n+1$).
\item La rédaction comporte toujours trois étapes annoncées explicitement : \textbf{Initialisation}, \textbf{Hérédité}, \textbf{Conclusion}.
\item Dans l'hérédité, on \emph{suppose} $P(n)$ et on \emph{démontre} $P(n+1)$ : ne jamais supposer la conclusion.
\item Variantes : initialisation décalée ($n \geq n_0$), récurrence forte (on suppose la propriété vraie jusqu'au rang $n$).
\item Applications classiques au Probatoire/Baccalauréat technique : formules de sommes, inégalités, divisibilité, majoration ou minoration de suites récurrentes.
\end{itemize}
\end{aretenir}

\begin{savaistu}[Un outil né avec les fondements des mathématiques]
Le raisonnement par récurrence est l'outil fondamental pour prouver les formules de sommes (comme celle de Gauss pour $1 + 2 + \cdots + n$) et les propriétés des suites définies par récurrence. Il a été formalisé à la fin du XIX\textsuperscript{e} siècle par le mathématicien italien Giuseppe Peano, dans ses célèbres axiomes construisant l'ensemble $\mathbb{N}$. Une anecdote célèbre raconte que le jeune Carl Friedrich Gauss, écolier, aurait retrouvé en quelques secondes la formule $\frac{n(n+1)}{2}$ lorsque son instituteur demanda de calculer $1 + 2 + \cdots + 100$ : il regroupa les termes par paires $(1 + 100), (2 + 99), \ldots$ chacune valant 101. Aujourd'hui, c'est ce même type de raisonnement qui permet aux ingénieurs de garantir qu'un algorithme répété des millions de fois (contrôle de machines, calculs de mensualités de prêts comme celui de Monsieur Kamga) produit toujours le résultat attendu.
\end{savaistu}

\begin{exercice}[Entraînement]
\begin{enumerate}
\item Démontrer par récurrence que pour tout $n \geq 1$ : $\displaystyle\sum_{k=1}^{n} k^2 = \frac{n(n+1)(2n+1)}{6}$.
\item Démontrer que pour tout $n \in \mathbb{N}$, $4^n - 1$ est divisible par 3.
\item Soit $(u_n)$ la suite définie par $u_0 = 2$ et $u_{n+1} = \dfrac{1}{2}u_n + 3$. Démontrer que pour tout $n \in \mathbb{N}$, $u_n \leq 6$.
\end{enumerate}
\end{exercice}

\begin{corrige}[Exercice 1]
\textbf{1.} Soit $P(n)$ : $\sum_{k=1}^{n} k^2 = \frac{n(n+1)(2n+1)}{6}$.

\emph{Initialisation.} Pour $n = 1$ : $1^2 = 1$ et $\frac{1 \times 2 \times 3}{6} = 1$. $P(1)$ est vraie.

\emph{Hérédité.} Soit $n \geq 1$ fixé, supposons $P(n)$ vraie. Alors :
\begin{align*}
\sum_{k=1}^{n+1} k^2 &= \frac{n(n+1)(2n+1)}{6} + (n+1)^2 = \frac{(n+1)\big[n(2n+1) + 6(n+1)\big]}{6} \\
&= \frac{(n+1)(2n^2 + 7n + 6)}{6} = \frac{(n+1)(n+2)(2n+3)}{6},
\end{align*}
car $2n^2 + 7n + 6 = (n+2)(2n+3)$. C'est $P(n+1)$.

\emph{Conclusion.} Par récurrence, la formule est vraie pour tout $n \geq 1$.

\textbf{2.} Soit $P(n)$ : \og il existe $k \in \mathbb{N}$ tel que $4^n - 1 = 3k$ \fg.

\emph{Initialisation.} $4^0 - 1 = 0 = 3 \times 0$ : $P(0)$ vraie.

\emph{Hérédité.} Soit $n \in \mathbb{N}$ fixé, supposons $4^n - 1 = 3k$ avec $k \in \mathbb{N}$. Alors $4^{n+1} - 1 = 4 \times 4^n - 1 = 4(3k + 1) - 1 = 12k + 3 = 3(4k + 1)$ : $P(n+1)$ vraie.

\emph{Conclusion.} $4^n - 1$ est divisible par 3 pour tout $n \in \mathbb{N}$.

\textbf{3.} Soit $P(n)$ : $u_n \leq 6$.

\emph{Initialisation.} $u_0 = 2 \leq 6$ : vraie.

\emph{Hérédité.} Soit $n \in \mathbb{N}$ fixé, supposons $u_n \leq 6$. Alors $u_{n+1} = \frac{1}{2}u_n + 3 \leq \frac{1}{2} \times 6 + 3 = 6$ : $P(n+1)$ vraie.

\emph{Conclusion.} Pour tout $n \in \mathbb{N}$, $u_n \leq 6$.
\end{corrige}

\coursec{Monotonie d'une suite : méthodes d'étude}

Dans la section précédente, nous avons appris à démontrer une propriété valable pour tous les termes d'une suite grâce au \cle{raisonnement par récurrence}. Nous allons maintenant utiliser cet outil pour répondre à une question très naturelle : une suite \emph{monte}-t-elle, \emph{descend}-t-elle, ou reste-t-elle stable ? C'est le problème de la \cle{monotonie}. Au quotidien technique, cette question est partout : le solde d'un compte épargne augmente-t-il ? La température d'un four en refroidissement diminue-t-elle toujours ? La production mensuelle d'un atelier progresse-t-elle régulièrement ?

\courssub{Les définitions : croître, décroître, stagner}

\textbf{Intuition.} Imaginez les termes d'une suite comme les marches d'un escalier. Si chaque marche est plus haute que la précédente, l'escalier monte : la suite est \emph{croissante}. Si chaque marche est plus basse, il descend : la suite est \emph{décroissante}. Si toutes les marches sont au même niveau, c'est un palier : la suite est \emph{constante}.

\begin{definition}[Suites monotones]
Soit $(u_n)$ une suite numérique définie pour $n \geq n_0$.
\begin{itemize}
\item $(u_n)$ est \cle{croissante} si, pour tout $n \geq n_0$ : $u_{n+1} \geq u_n$.
\item $(u_n)$ est \cle{strictement croissante} si, pour tout $n \geq n_0$ : $u_{n+1} > u_n$.
\item $(u_n)$ est \cle{décroissante} si, pour tout $n \geq n_0$ : $u_{n+1} \leq u_n$.
\item $(u_n)$ est \cle{strictement décroissante} si, pour tout $n \geq n_0$ : $u_{n+1} < u_n$.
\item $(u_n)$ est \cle{constante} si, pour tout $n \geq n_0$ : $u_{n+1} = u_n$.
\item $(u_n)$ est \cle{monotone} si elle est croissante ou décroissante ; \cle{strictement monotone} si elle est strictement croissante ou strictement décroissante.
\end{itemize}
\end{definition}

\begin{attention}[Une suite n'est pas forcément monotone]
Certaines suites ne sont ni croissantes ni décroissantes. Par exemple $u_n = (-1)^n$ vaut alternativement $-1$ et $1$ : elle \emph{oscille}. Dire « la suite n'est pas croissante » ne signifie donc pas « elle est décroissante ». Autre piège de rédaction : vérifier $u_1 > u_0$ sur les deux premiers termes ne prouve \textbf{rien} sur l'ensemble de la suite ; il faut une comparaison valable pour \textbf{tout} entier $n$.
\end{attention}

\courssub{Méthode 1 : le signe de la différence $u_{n+1} - u_n$}

\textbf{Intuition.} Pour savoir si l'on monte ou descend entre deux marches consécutives, on mesure l'écart : $u_{n+1} - u_n$. Si cet écart est toujours positif, on monte à chaque pas ; s'il est toujours négatif, on descend.

\begin{propriete}[Critère de la différence]
\begin{itemize}
\item Si, pour tout $n$, $u_{n+1} - u_n \geq 0$, alors $(u_n)$ est croissante.
\item Si, pour tout $n$, $u_{n+1} - u_n \leq 0$, alors $(u_n)$ est décroissante.
\item Avec des inégalités strictes, on obtient la monotonie stricte.
\end{itemize}
\end{propriete}

C'est la méthode \emph{de base}, à essayer en premier. Elle est particulièrement efficace pour les suites \textbf{polynomiales} et \textbf{rationnelles}, car la différence se calcule et se factorise souvent bien.

\begin{exemplebox}[Étude complète par la différence]
Étudions la monotonie de la suite définie pour $n \geq 1$ par $u_n = \dfrac{n^2+1}{2n^2-3}$.

\textbf{Étape 1 : calculer la différence.}
\begin{align*}
u_{n+1} - u_n &= \frac{(n+1)^2+1}{2(n+1)^2-3} - \frac{n^2+1}{2n^2-3} \\
&= \frac{n^2+2n+2}{2n^2+4n-1} - \frac{n^2+1}{2n^2-3} \\
&= \frac{(n^2+2n+2)(2n^2-3) - (n^2+1)(2n^2+4n-1)}{(2n^2+4n-1)(2n^2-3)}.
\end{align*}

\textbf{Étape 2 : développer le numérateur.}
\begin{align*}
\text{Numérateur} &= (2n^4 + 4n^3 - 3n^2 - 6n + 4n^2 - 6) - (2n^4 + 4n^3 - n^2 + 2n^2 + 4n - 1) \\
&= (2n^4 + 4n^3 + n^2 - 6n - 6) - (2n^4 + 4n^3 + n^2 + 4n - 1) \\
&= -10n - 5.
\end{align*}
Donc
\[ u_{n+1} - u_n = \frac{-10n - 5}{(2n^2+4n-1)(2n^2-3)}. \]

\textbf{Étape 3 : étudier le signe.} Pour $n \geq 1$ :
\begin{itemize}
\item numérateur : $-10n - 5 < 0$ (somme de termes négatifs) ;
\item dénominateur : le facteur $2n^2+4n-1$ est positif pour $n \geq 1$ (valeur $5$ pour $n=1$, croissant ensuite). Le facteur $2n^2-3$ est négatif pour $n=1$ ($-1$), nul pour $n \approx 1,22$, et positif pour $n \geq 2$.
\end{itemize}
\begin{itemize}
\item Pour $n = 1$ : numérateur $< 0$, dénominateur $< 0$ $\Rightarrow$ différence $> 0$ : $u_2 > u_1$.
\item Pour $n \geq 2$ : numérateur $< 0$, dénominateur $> 0$ $\Rightarrow$ différence $< 0$ : $u_{n+1} < u_n$.
\end{itemize}
La suite est donc \textbf{décroissante à partir du rang 2}. (Vérification numérique : $u_1 = -2$, $u_2 = 1$, $u_3 = \tfrac{10}{15} \approx 0{,}67$, $u_4 = \tfrac{17}{29} \approx 0{,}59$ : la suite augmente entre $n=1$ et $n=2$, puis descend régulièrement.)
\end{exemplebox}

\courssub{Méthode 2 : comparer le quotient $\dfrac{u_{n+1}}{u_n}$ à 1}

\textbf{Intuition.} Quand les termes sont tous positifs, au lieu de soustraire, on peut \emph{diviser} : si chaque terme vaut plus de $100\,\%$ du précédent, la suite augmente.

\begin{propriete}[Critère du quotient — suites à termes strictement positifs]
Soit $(u_n)$ une suite telle que $u_n > 0$ pour tout $n$.
\begin{itemize}
\item Si, pour tout $n$, $\dfrac{u_{n+1}}{u_n} \geq 1$, alors $(u_n)$ est croissante.
\item Si, pour tout $n$, $\dfrac{u_{n+1}}{u_n} \leq 1$, alors $(u_n)$ est décroissante.
\end{itemize}
\end{propriete}

Cette méthode est \textbf{indispensable} pour les suites géométriques, les suites contenant des exponentielles, des puissances ou des factorielles, car le quotient se simplifie spectaculairement.

\begin{exemplebox}[Quotient pour une suite avec puissances]
Soit $u_n = \dfrac{3^n}{n!}$ pour $n \geq 1$. Tous les termes sont strictement positifs.
\[ \frac{u_{n+1}}{u_n} = \frac{3^{n+1}}{(n+1)!} \times \frac{n!}{3^n} = \frac{3}{n+1}. \]
Or $\dfrac{3}{n+1} \leq 1 \iff n + 1 \geq 3 \iff n \geq 2$. La suite est donc \textbf{décroissante à partir du rang 2}.
\end{exemplebox}

\begin{attention}[Le quotient exige des termes de signe constant]
Le critère du quotient n'est valable que si $u_n > 0$ pour tout $n$ (ou, avec précautions, $u_n < 0$ pour tout $n$). Si les termes changent de signe, comparer le quotient à 1 ne dit \textbf{rien} sur la monotonie : multiplier une inégalité par un nombre négatif en change le sens ! Avant d'utiliser cette méthode, il faut donc \emph{prouver} que $u_n > 0$, souvent par récurrence.
\end{attention}

\begin{savaistu}[Différence et quotient disent la même chose]
Pour une suite à termes strictement positifs, les deux critères sont équivalents :
\[ \frac{u_{n+1}}{u_n} \geq 1 \iff u_{n+1} \geq u_n \iff u_{n+1} - u_n \geq 0, \]
car multiplier par $u_n > 0$ conserve le sens de l'inégalité. Le choix entre les deux méthodes est donc une simple question de \emph{confort de calcul} : différence pour les sommes et fractions, quotient pour les produits, puissances et factorielles.
\end{savaistu}

\begin{methode}[Étudier la monotonie d'une suite explicite $u_n = f(n)$]
\begin{enumerate}
\item Regarder la \emph{forme} de $u_n$ : sommes, fractions $\rightarrow$ \textbf{différence} ; produits, puissances, factorielles $\rightarrow$ \textbf{quotient}.
\item Si quotient : vérifier (ou prouver) que $u_n > 0$ pour tout $n$.
\item Calculer $u_{n+1} - u_n$ (ou $\frac{u_{n+1}}{u_n}$) et \textbf{simplifier au maximum} (factoriser !).
\item Étudier le signe (ou comparer à 1) pour tout $n$ du domaine.
\item Conclure en rédigeant : « pour tout $n \geq \ldots$, $u_{n+1} - u_n \geq 0$, donc $(u_n)$ est croissante ».
\end{enumerate}
\end{methode}

\courssub{Méthode 3 : suites récurrentes $u_{n+1} = f(u_n)$}

\textbf{Intuition.} Quand la suite est définie par $u_{n+1} = f(u_n)$, chaque terme est l'\emph{image} du précédent par une fonction $f$. Si $f$ « conserve l'ordre » (fonction croissante), alors l'ordre entre deux termes se transmet au rang suivant : la monotonie se propage comme dans un raisonnement par récurrence.

\begin{propriete}[Monotonie et fonction associée]
Soit $(u_n)$ définie par $u_{n+1} = f(u_n)$, où $f$ est définie sur un intervalle $I$ contenant tous les termes de la suite.
\begin{itemize}
\item Si $f$ est \textbf{croissante} sur $I$, alors $(u_n)$ est \textbf{monotone} : croissante si $u_1 \geq u_0$, décroissante si $u_1 \leq u_0$.
\item Si $f$ est \textbf{décroissante} sur $I$, alors $(u_n)$ \textbf{n'est pas monotone} en général : les sous-suites $(u_{2n})$ et $(u_{2n+1})$ sont monotones de sens contraires (la suite « oscille »).
\end{itemize}
\end{propriete}

\begin{experience}[Démonstration guidée : si $f$ est croissante et $u_1 \geq u_0$, alors $(u_n)$ est croissante]
Démontrons par récurrence la propriété $\mathcal{P}(n)$ : « $u_{n+1} \geq u_n$ ».
\begin{enumerate}
\item \textbf{Initialisation.} $\mathcal{P}(0)$ est vraie par hypothèse : $u_1 \geq u_0$.
\item \textbf{Hérédité.} Supposons $u_{n+1} \geq u_n$ pour un rang $n$ fixé. Comme $f$ est croissante sur $I$ et que $u_n, u_{n+1} \in I$, appliquer $f$ conserve l'inégalité :
\[ f(u_{n+1}) \geq f(u_n) \quad \text{c'est-à-dire} \quad u_{n+2} \geq u_{n+1}. \]
Donc $\mathcal{P}(n+1)$ est vraie.
\item \textbf{Conclusion.} Par récurrence, $u_{n+1} \geq u_n$ pour tout $n$ : $(u_n)$ est croissante. $\square$
\end{enumerate}
Remarquez le point crucial : l'hérédité exige que \emph{tous les termes restent dans} $I$, l'intervalle où $f$ est croissante. C'est pourquoi on commence toujours par prouver un encadrement des termes.
\end{experience}

\begin{methode}[Étudier la monotonie d'une suite récurrente $u_{n+1} = f(u_n)$]
\begin{enumerate}
\item \textbf{Étape 1 — Encadrement.} Deviner un intervalle $I$ stable par $f$ et prouver \textbf{par récurrence} que $u_n \in I$ pour tout $n$ (majoration/minoration).
\item \textbf{Étape 2 — Sens de variation de $f$.} Montrer que $f$ est monotone sur $I$ (dérivée, ou composition de fonctions de sens connu).
\item \textbf{Étape 3 — Initialisation.} Calculer $u_0$ et $u_1$ et les comparer.
\item \textbf{Conclusion.} Si $f$ est croissante sur $I$ : $u_1 \geq u_0 \Rightarrow (u_n)$ croissante ; $u_1 \leq u_0 \Rightarrow (u_n)$ décroissante. Si $f$ est décroissante : la suite n'est pas monotone (oscillation).
\end{enumerate}
\end{methode}

\begin{exoresolu}[La suite $u_{n+1} = \sqrt{2 + u_n}$, $u_0 = 0$]
Soit $(u_n)$ définie par $u_0 = 0$ et, pour tout $n \in \mathbb{N}$, $u_{n+1} = \sqrt{2 + u_n}$. Montrer que $(u_n)$ est croissante.
\tcblower
\textbf{Étape 1 : par récurrence, $0 \leq u_n < 2$ pour tout $n$.}

\emph{Initialisation} : $u_0 = 0$, donc $0 \leq u_0 < 2$ : vrai.

\emph{Hérédité} : supposons $0 \leq u_n < 2$. Alors $2 \leq 2 + u_n < 4$, et la fonction racine carrée étant croissante sur $[0;+\infty[$ :
\[ \sqrt{2} \leq \sqrt{2+u_n} < \sqrt{4} \quad \text{soit} \quad 0 < \sqrt{2} \leq u_{n+1} < 2. \]
La propriété est héréditaire. Par récurrence, $0 \leq u_n < 2$ pour tout $n$ : tous les termes vivent dans $I = [0;2]$.

\textbf{Étape 2 : monotonie de $f$.} Posons $f(x) = \sqrt{2+x}$. La fonction $x \mapsto 2+x$ est croissante et positive sur $[0;2]$, et la racine carrée est croissante sur $[0;+\infty[$ : par composition, $f$ est \textbf{croissante sur $[0;2]$}.

\textbf{Étape 3 : comparaison de $u_0$ et $u_1$.} $u_1 = \sqrt{2+0} = \sqrt{2} \approx 1{,}41$, donc $u_1 > u_0$.

\textbf{Conclusion.} $f$ est croissante sur $[0;2]$ qui contient tous les termes, et $u_1 \geq u_0$ : par la propriété précédente (démontrée par récurrence), la suite $(u_n)$ est \textbf{croissante}. De plus elle est majorée par $2$ — ceci sera décisif pour la convergence (section suivante).
\end{exoresolu}

\begin{popfigure}
\begin{center}
\begin{tikzpicture}[scale=1.1]
% ---- f croissante : escalier monotone ----
\begin{scope}
\draw[->,popmuted] (-0.2,0) -- (2.3,0) node[below left]{$x$};
\draw[->,popmuted] (0,-0.2) -- (0,2.3) node[below left]{$y$};
\draw[popmuted,dashed] (0,0) -- (2.2,2.2) node[pos=0.85,below right]{$y=x$};
% f(x) = sqrt(2+x)
\draw[popblue,very thick,domain=0:2.2,samples=100] plot (\x,{sqrt(max(0,2+\x))}) node[right]{$y=f(x)$};
% Escalier (cobweb) depuis u0 = 0
\pgfmathsetmacro{\uA}{0}
\pgfmathsetmacro{\uB}{sqrt(2+\uA)} % 1.414
\pgfmathsetmacro{\uC}{sqrt(2+\uB)} % 1.847
\pgfmathsetmacro{\uD}{sqrt(2+\uC)} % 1.961
\draw[poporange,thick] (\uA,0) node[below]{$u_0$} -- (\uA,\uB);
\draw[poporange,thick] (\uA,\uB) -- (\uB,\uB);
\draw[poporange,thick] (\uB,\uB) -- (\uB,\uC);
\draw[poporange,thick] (\uB,\uC) -- (\uC,\uC);
\draw[poporange,thick] (\uC,\uC) -- (\uC,\uD);
\node[popdark,font=\small] at (1.1,-0.4) {$f$ croissante : montée monotone};
\end{scope}
% ---- f décroissante : toile d'araignée ----
\begin{scope}[xshift=7cm]
\draw[->,popmuted] (-0.2,0) -- (2.3,0) node[below left]{$x$};
\draw[->,popmuted] (0,-0.2) -- (0,2.3) node[below left]{$y$};
\draw[popmuted,dashed] (0,0) -- (2.2,2.2) node[pos=0.85,below right]{$y=x$};
% f(x) = 1/x
\draw[poppurple,very thick,domain=0.45:2.2,samples=100] plot (\x,{1/\x}) node[right]{$y=f(x)$};
% Spirale : u0=0.5
\draw[poporange,thick] (0.5,0) node[below]{$u_0$} -- (0.5,2);
\draw[poporange,thick] (0.5,2) -- (2,2);
\draw[poporange,thick] (2,2) -- (2,0.5);
\draw[poporange,thick] (2,0.5) -- (0.5,0.5);
\draw[poporange,thick] (0.5,0.5) -- (0.5,2);
\node[popdark,font=\small] at (1.1,-0.4) {$f$ décroissante : oscillation};
\end{scope}
\end{tikzpicture}
\end{center}
\legende{Construction géométrique des termes de $u_{n+1}=f(u_n)$ : on lit $u_{n+1}$ sur la courbe puis on le « reporte » sur l'axe grâce à la droite $y=x$. À gauche, $f(x)=\sqrt{2+x}$ croissante : l'escalier monte régulièrement vers le point fixe $2$. À droite, $f(x)=1/x$ décroissante : la « toile d'araignée » tourne autour du point fixe $1$, la suite oscille ($0,5 \to 2 \to 0,5 \dots$).}
\end{popfigure}

\courssub{Synthèse visuelle : les deux tests de signe}

Le tableau ci-dessous résume comment conclure selon le signe obtenu.

\begin{center}
\begin{tabularx}{\linewidth}{|l|X|X|}
\hline
\rowcolor{popblueL} \textbf{Résultat obtenu} & \textbf{Conclusion} & \textbf{Condition} \\
\hline
$u_{n+1} - u_n \geq 0$ pour tout $n$ & $(u_n)$ croissante & aucune \\
\hline
$u_{n+1} - u_n \leq 0$ pour tout $n$ & $(u_n)$ décroissante & aucune \\
\hline
$\dfrac{u_{n+1}}{u_n} \geq 1$ pour tout $n$ & $(u_n)$ croissante & $u_n > 0$ \\
\hline
$\dfrac{u_{n+1}}{u_n} \leq 1$ pour tout $n$ & $(u_n)$ décroissante & $u_n > 0$ \\
\hline
\end{tabularx}
\end{center}

\begin{popfigure}
\begin{center}
\begin{tikzpicture}[scale=1.0]
\node[font=\small] at (-3.4,0.6) {$u_{n+1}-u_n$};
\draw[thick] (-2.4,0.6) -- (2.4,0.6);
\draw[thick] (0,0.35) -- (0,0.85);
\node[popgreen,font=\small] at (-1.2,0.6) {$+$};
\node[popgreen,font=\small] at (1.2,0.6) {$+$};
\node[font=\small] at (-1.2,0.1) {croissante};
\node[font=\small] at (1.2,0.1) {croissante};
\node[font=\small] at (-3.4,-1.0) {$u_{n+1}-u_n$};
\draw[thick] (-2.4,-1.0) -- (2.4,-1.0);
\draw[thick] (0,-1.25) -- (0,-0.75);
\node[popink,font=\small] at (-1.2,-1.0) {$-$};
\node[popink,font=\small] at (1.2,-1.0) {$-$};
\node[font=\small] at (-1.2,-1.5) {décroissante};
\node[font=\small] at (1.2,-1.5) {décroissante};
\node[font=\small] at (-3.4,-2.6) {$\dfrac{u_{n+1}}{u_n}-1$};
\draw[thick] (-2.4,-2.6) -- (2.4,-2.6);
\draw[thick] (0,-2.85) -- (0,-2.35);
\node[popgreen,font=\small] at (1.2,-2.6) {$+$};
\node[popink,font=\small] at (-1.2,-2.6) {$-$};
\node[font=\small] at (1.2,-3.1) {croissante};
\node[font=\small] at (-1.2,-3.1) {décroissante};
\node[font=\small,popmuted] at (0,-3.7) {(troisième ligne valable seulement si $u_n>0$)};
\end{tikzpicture}
\end{center}
\legende{Lecture des signes : le signe constant de $u_{n+1}-u_n$ (ou de $\frac{u_{n+1}}{u_n}-1$ si $u_n>0$) donne directement le sens de variation de la suite.}
\end{popfigure}

\begin{aretenir}[L'essentiel sur la monotonie]
\begin{itemize}
\item $(u_n)$ croissante $\iff u_{n+1} - u_n \geq 0$ pour tout $n$.
\item Si $u_n > 0$ : $(u_n)$ croissante $\iff \dfrac{u_{n+1}}{u_n} \geq 1$.
\item Si $u_{n+1} = f(u_n)$ avec $f$ \textbf{croissante} sur un intervalle contenant tous les termes : la suite est monotone, et son sens est donné par la comparaison de $u_0$ et $u_1$.
\item Si $f$ est \textbf{décroissante} : la suite oscille, elle n'est pas monotone.
\item La rédaction complète d'une étude de suite récurrente comporte toujours trois étapes : encadrement par récurrence, monotonie de $f$, comparaison des deux premiers termes.
\end{itemize}
\end{aretenir}

\begin{exercice}[À vous de jouer]
\begin{enumerate}
\item Étudier la monotonie de la suite $u_n = n^2 - 4n + 3$ pour $n \geq 0$ (par la différence).
\item Soit $v_n = \dfrac{2^n}{n}$ pour $n \geq 1$. Montrer que $v_n > 0$, puis étudier la monotonie par le quotient.
\item Soit $(w_n)$ définie par $w_0 = 1$ et $w_{n+1} = \dfrac{w_n}{2} + 1$. Montrer par récurrence que $0 \leq w_n \leq 2$, étudier la monotonie de $f(x) = \dfrac{x}{2}+1$ sur $[0;2]$, comparer $w_0$ et $w_1$, et conclure.
\end{enumerate}
\end{exercice}

\begin{corrige}[Exercices 1 à 3]
\textbf{1.} $u_{n+1} - u_n = [(n+1)^2 - 4(n+1) + 3] - [n^2 - 4n + 3] = 2n - 3$. Pour $n \geq 2$, $2n-3 > 0$ : la suite est croissante à partir du rang 2 (elle descend d'abord de $u_0=3$ à $u_2 = -1$).

\textbf{2.} $v_n = \dfrac{2^n}{n} > 0$ comme quotient de quantités positives. $\dfrac{v_{n+1}}{v_n} = \dfrac{2^{n+1}}{n+1}\times\dfrac{n}{2^n} = \dfrac{2n}{n+1}$. Or $\dfrac{2n}{n+1} \geq 1 \iff 2n \geq n+1 \iff n \geq 1$ : toujours vrai pour $n \ge 1$. Donc $(v_n)$ est croissante.

\textbf{3.} \emph{Encadrement} : $w_0 = 1 \in [0;2]$ ; si $0 \leq w_n \leq 2$ alors $0 \leq \tfrac{w_n}{2} \leq 1$ donc $1 \leq w_{n+1} \leq 2 \subset [0;2]$ : hérédité acquise, donc $w_n \in [0;2]$ pour tout $n$. \emph{Monotonie de $f$} : $f(x) = \tfrac{x}{2}+1$ est affine de coefficient $\tfrac12 > 0$, donc croissante sur $[0;2]$. \emph{Comparaison} : $w_1 = \tfrac{1}{2}+1 = \tfrac32 > w_0 = 1$. \emph{Conclusion} : $(w_n)$ est croissante.
\end{corrige}

\coursec{Bornitude et Convergence : Théorèmes fondamentaux}

Dans la section précédente, nous avons appris à déterminer si une suite est \cle{croissante} ou \cle{décroissante} : par le signe de $u_{n+1}-u_n$, par le quotient $\dfrac{u_{n+1}}{u_n}$ lorsque les termes sont strictement positifs, ou par récurrence. Mais une question fondamentale reste posée : une suite croissante finit-elle toujours par « s'arrêter » quelque part ? Pensez à un technicien qui remplit un réservoir de \SI{1000}{L} : le niveau monte, monte... mais il ne peut pas dépasser la capacité du réservoir. Le niveau croît et reste \emph{plafonné} : il va nécessairement se stabiliser. C'est exactement l'idée des théorèmes fondamentaux de cette section : la \cle{monotonie} (section 2) combinée à la \cle{bornitude} (notion nouvelle) garantit la \cle{convergence}.

\courssub{Suites majorées, minorées, bornées}

\textbf{Intuition.}\par
Un \cle{majorant} est un « plafond » que les termes de la suite ne dépassent jamais. Un \cle{minorant} est un « plancher » en dessous duquel les termes ne descendent jamais. Attention : un plafond n'a pas besoin d'être atteint ! Si tous les termes sont inférieurs à $5$, alors $5$ est un majorant, mais $6$, $100$ aussi. Le \emph{meilleur} plafond est le plus petit possible.

\begin{definition}[Suite majorée, minorée, bornée]
Soit $(u_n)$ une suite numérique réelle.
\begin{itemize}
\item On dit que $(u_n)$ est \cle{majorée} s'il existe un réel $M$ tel que, pour tout entier naturel $n$ : $u_n \leq M$. Un tel réel $M$ est appelé un \cle{majorant} de la suite.
\item On dit que $(u_n)$ est \cle{minorée} s'il existe un réel $m$ tel que, pour tout entier naturel $n$ : $u_n \geq m$. Un tel réel $m$ est appelé un \cle{minorant} de la suite.
\item On dit que $(u_n)$ est \cle{bornée} si elle est à la fois majorée et minorée, c'est-à-dire s'il existe deux réels $m$ et $M$ tels que, pour tout $n$ : $m \leq u_n \leq M$.
\end{itemize}
\end{definition}

\begin{definition}[Borne supérieure et borne inférieure (notions admises)]
Si $(u_n)$ est majorée, l'ensemble de ses majorants admet un \emph{plus petit élément} : on l'appelle la \cle{borne supérieure} des valeurs de la suite, notée $\sup(u_n)$. C'est le « meilleur plafond » : le plus bas possible.

De même, si $(u_n)$ est minorée, le \emph{plus grand} des minorants est la \cle{borne inférieure}, notée $\inf(u_n)$ : le « meilleur plancher ».

Caractérisation essentielle de la borne supérieure $S$ (admise) : pour tout réel $\varepsilon > 0$, aussi petit soit-il, il existe au moins un terme de la suite tel que $u_n > S - \varepsilon$. Autrement dit, on ne peut pas abaisser le plafond, même d'un cheveu $\varepsilon$, sans qu'un terme de la suite ne le dépasse.
\end{definition}

\begin{exemplebox}[Premiers exemples]
\begin{itemize}
\item La suite $u_n = \dfrac{n}{n+1}$ vérifie $0 \leq u_n < 1$ pour tout $n$. Elle est minorée par $0$ et majorée par $1$ (donc bornée). Sa borne supérieure est $1$ (jamais atteinte, mais approchée d'aussi près qu'on veut).
\item La suite $u_n = n^2$ est minorée par $0$ mais n'est \emph{pas} majorée : ses termes dépassent tout plafond.
\item La suite $u_n = (-1)^n$ est bornée : $-1 \leq u_n \leq 1$.
\end{itemize}
\end{exemplebox}

\courssub{Théorème 1 : toute suite convergente est bornée}

\textbf{Intuition.}\par
Si les termes d'une suite se rapprochent indéfiniment d'une limite $\ell$, alors à partir d'un certain rang ils sont tous « coincés » dans un petit couloir autour de $\ell$. Il ne reste qu'un nombre \emph{fini} de termes avant ce rang, et un ensemble fini possède toujours un plus grand et un plus petit élément. La suite est donc prisonnière entre un plancher et un plafond.

\begin{propriete}[Théorème 1 — admis]
Toute suite \cle{convergente} est \cle{bornée}.
\end{propriete}

\begin{attention}[La réciproque est fausse !]
Une suite bornée n'est pas forcément convergente. Contre-exemple classique : la suite $u_n = (-1)^n$. Elle est bornée ($-1 \leq u_n \leq 1$) pourtant elle \emph{diverge} : ses termes oscillent éternellement entre $-1$ et $1$ sans se stabiliser. Retenir : \textbf{bornée n'implique pas convergente}. Au Probatoire, écrire « la suite est bornée donc elle converge » est une erreur éliminatoire : il manque l'hypothèse de monotonie.
\end{attention}

\courssub{Théorème 2 : le théorème de la limite monotone}

\textbf{Intuition.}\par
Revenons au réservoir : le niveau monte (suite croissante) et ne peut pas dépasser \SI{1000}{L} (suite majorée). Le niveau n'a alors \emph{nulle part où aller} : il est obligé de se stabiliser vers une valeur limite, qui est exactement le meilleur plafond, la borne supérieure. C'est l'un des théorèmes les plus importants de l'année.

\begin{propriete}[Théorème 2 — Théorème de la limite monotone (admis)]
\begin{itemize}
\item Toute suite \cle{croissante} et \cle{majorée} est convergente. Sa limite est la borne supérieure de ses valeurs : $\lim u_n = \sup(u_n)$.
\item Toute suite \cle{décroissante} et \cle{minorée} est convergente. Sa limite est la borne inférieure de ses valeurs : $\lim u_n = \inf(u_n)$.
\end{itemize}
\end{propriete}

\begin{experience}[Démonstration guidée : suite croissante majorée]
Soit $(u_n)$ une suite croissante et majorée. Notons $S = \sup(u_n)$ sa borne supérieure (elle existe car la suite est majorée — ceci est admis). Montrons que $(u_n)$ converge vers $S$, c'est-à-dire que pour tout $\varepsilon > 0$, il existe un rang $N$ tel que pour tout $n \geq N$ : $|u_n - S| < \varepsilon$.

\textbf{Étape 1 : utiliser la caractérisation du supremum.}
Soit $\varepsilon > 0$ fixé. Puisque $S$ est le \emph{plus petit} des majorants, le nombre $S - \varepsilon$ n'est \emph{pas} un majorant (sinon $S$ ne serait pas le plus petit). Il existe donc un rang $N$ tel que :
\[ u_N > S - \varepsilon. \]

\textbf{Étape 2 : utiliser la croissance.}
La suite est croissante, donc pour tout $n \geq N$ :
\[ u_n \geq u_N > S - \varepsilon. \]

\textbf{Étape 3 : utiliser le fait que $S$ est un majorant.}
Pour tout entier $n$, $u_n \leq S$, donc a fortiori pour $n \geq N$ :
\[ S - \varepsilon < u_n \leq S < S + \varepsilon. \]

\textbf{Étape 4 : conclure.}
Pour tout $n \geq N$, on a $|u_n - S| < \varepsilon$. Comme $\varepsilon > 0$ était arbitraire, ceci prouve exactement que $\lim\limits_{n \to +\infty} u_n = S$. \hfill $\square$

Le cas d'une suite décroissante minorée se traite de manière symétrique avec la borne inférieure.
\end{experience}

\begin{popfigure}
\begin{center}
\begin{tikzpicture}[scale=1.05]
% axes
\draw[->,popmuted] (0,0) -- (10.6,0) node[below left] {$n$};
\draw[->,popmuted] (0,0) -- (0,3.4) node[below left] {$u_n$};
% limite S
\draw[popblue,thick,dashed] (0,2.4) -- (10.4,2.4) node[right,popblue] {$S=\sup(u_n)$};
% zone epsilon
\fill[popblueL,opacity=0.55] (0,1.8) rectangle (10.4,2.4);
\draw[<->,popblue] (10.15,1.8) -- (10.15,2.4) node[midway,right,popblue] {$\varepsilon$};
\draw[popblue!60] (0,1.8) -- (10.4,1.8);
% points de la suite croissante vers 2.4
\foreach \x/\y in {0.5/0.35, 1.5/0.9, 2.5/1.35, 3.5/1.68, 4.5/1.9, 5.5/2.05, 6.5/2.16, 7.5/2.24, 8.5/2.29, 9.5/2.33}
  \fill[poporange] (\x,\y) circle (0.09);
% rang N
\draw[popgreen,thick] (4.5,0) -- (4.5,1.9);
\node[popgreen,below] at (4.5,0) {$N$};
\node[popgreen,align=center] at (7.2,0.9) {à partir du rang $N$,\\ tous les termes sont\\ dans la bande $]S-\varepsilon\,;\,S]$};
\end{tikzpicture}
\end{center}
\legende{Suite croissante majorée : les termes montent et finissent piégés dans la bande $]S-\varepsilon\,;\,S]$, aussi étroite soit-elle. La suite converge vers $S$.}
\end{popfigure}

\begin{attention}[Ce que le théorème dit... et ne dit pas]
\begin{itemize}
\item Le théorème garantit l'\emph{existence} de la limite, mais ne donne pas toujours sa \emph{valeur} de façon explicite : il faudra souvent la calculer ensuite (voir la méthode en 4 étapes plus bas).
\item Les deux hypothèses sont indispensables : $u_n = n$ est croissante mais non majorée (elle diverge vers $+\infty$) ; $u_n = (-1)^n$ est bornée mais non monotone (elle diverge).
\item Une suite croissante \emph{non majorée} diverge vers $+\infty$ ; une suite décroissante \emph{non minorée} diverge vers $-\infty$.
\end{itemize}
\end{attention}

\courssub{Théorème 3 : passage à la limite dans les inégalités et théorème des gendarmes}

\begin{propriete}[Théorème 3 — admis]
Soient $(u_n)$ et $(v_n)$ deux suites convergentes, de limites respectives $\ell$ et $\ell'$. Si, à partir d'un certain rang, $u_n \leq v_n$, alors :
\[ \lim_{n \to +\infty} u_n \leq \lim_{n \to +\infty} v_n \quad \text{c'est-à-dire} \quad \ell \leq \ell'. \]
Autrement dit : \emph{le passage à la limite conserve les inégalités larges}.
\end{propriete}

\begin{attention}[Les inégalités strictes deviennent larges !]
Même si $u_n < v_n$ pour \emph{tout} $n$, on ne peut conclure que $\ell \leq \ell'$, jamais $\ell < \ell'$. Exemple : $u_n = 0$ et $v_n = \dfrac{1}{n}$. On a $u_n < v_n$ pour tout $n \geq 1$, pourtant les deux limites sont égales à $0$. Écrire $\ell < \ell'$ serait faux.
\end{attention}

\begin{propriete}[Théorème des gendarmes — admis]
Soient $(u_n)$, $(v_n)$ et $(w_n)$ trois suites telles qu'à partir d'un certain rang :
\[ v_n \leq u_n \leq w_n. \]
Si $(v_n)$ et $(w_n)$ convergent vers la \emph{même} limite $\ell$, alors $(u_n)$ converge aussi vers $\ell$.
\end{propriete}

\textbf{Intuition.}\par
Deux « gendarmes » $(v_n)$ et $(w_n)$ escortent un « prisonnier » $(u_n)$. Si les deux gendarmes se dirigent vers le même poste $\ell$, le prisonnier, coincé entre eux, n'a d'autre choix que d'y aller aussi.

\begin{exemplebox}[Application immédiate des gendarmes]
Étudions la convergence de $u_n = \dfrac{\sin(n)}{n}$. Pour tout $n \geq 1$ :
\[ -1 \leq \sin(n) \leq 1 \quad \Longrightarrow \quad -\frac{1}{n} \leq \frac{\sin(n)}{n} \leq \frac{1}{n}. \]
Or $\lim\limits_{n\to+\infty} -\dfrac{1}{n} = \lim\limits_{n\to+\infty} \dfrac{1}{n} = 0$. Par le théorème des gendarmes, $\lim\limits_{n\to+\infty} u_n = 0$.
\end{exemplebox}

\courssub{Stratégie type pour les suites récurrentes $u_{n+1} = f(u_n)$}

De nombreuses suites du programme (et de la vie technique : remboursement d'un crédit, concentration d'un produit, algorithme de calcul) sont définies par une relation de récurrence $u_{n+1} = f(u_n)$. Voici la démarche complète, à connaître par cœur pour le Probatoire.

\begin{methode}[Prouver la convergence de $u_{n+1} = f(u_n)$ — algorithme en 4 étapes]
\begin{enumerate}
\item \textbf{Bornitude par récurrence.} Conjecturer un intervalle $[m\,;\,M]$ stable (souvent donné par l'énoncé ou deviné sur les premiers termes), puis prouver par récurrence : « pour tout $n$, $m \leq u_n \leq M$ ». L'hérédité utilise le fait que $f([m\,;\,M]) \subset [m\,;\,M]$.
\item \textbf{Monotonie.} Étudier le signe de $u_{n+1} - u_n = f(u_n) - u_n$ en utilisant l'encadrement de l'étape 1 (ou comparer $u_1$ et $u_0$ puis raisonner par récurrence si $f$ est croissante).
\item \textbf{Conclusion par le théorème de la limite monotone.} Suite croissante majorée (ou décroissante minorée) $\Rightarrow$ $(u_n)$ converge vers une limite $\ell \in [m\,;\,M]$.
\item \textbf{Calcul de la limite.} Par continuité de $f$ (admise), passer à la limite dans $u_{n+1} = f(u_n)$ : la limite $\ell$ vérifie l'\cle{équation des points fixes} $f(\ell) = \ell$. Résoudre cette équation et \emph{sélectionner} la solution compatible avec l'intervalle $[m\,;\,M]$ et avec la monotonie.
\end{enumerate}
\end{methode}

\begin{attention}[Erreur fréquente au Probatoire]
L'équation $f(\ell) = \ell$ peut avoir \emph{plusieurs} solutions ! Il faut vérifier que la limite trouvée appartient bien à l'intervalle où vivent les termes de la suite et qu'elle est cohérente avec la monotonie. Par exemple, si la suite est croissante avec $u_0 = 1$, sa limite ne peut pas valoir $0$. Une solution de $f(\ell)=\ell$ hors de $[m\,;\,M]$ doit être rejetée, avec justification. Oublier cette vérification fait perdre des points même quand tout le reste est juste.
\end{attention}

\begin{exoresolu}[Méthode de Héron : calcul de $\sqrt{3}$]
On considère la suite définie par $u_0 = 1$ et, pour tout $n \in \mathbb{N}$ :
\[ u_{n+1} = \frac{1}{2}\left(u_n + \frac{3}{u_n}\right). \]
C'est l'algorithme babylonien (ou méthode de Héron) pour approcher $\sqrt{3}$, encore utilisé dans les calculatrices. Montrer que $(u_n)$ converge et déterminer sa limite.

\tcblower
\textbf{Solution détaillée.}

\textbf{Étape 1 : bornitude.} Montrons par récurrence que pour tout $n \geq 1$ : $\sqrt{3} \leq u_n \leq 3$.

\emph{Initialisation :} $u_1 = \dfrac{1}{2}\left(1 + \dfrac{3}{1}\right) = 2$, et $\sqrt{3} \approx 1{,}732 \leq 2 \leq 3$. La propriété est vraie au rang $1$.

\emph{Hérédité :} supposons $\sqrt{3} \leq u_n \leq 3$ pour un rang $n \geq 1$.
\begin{itemize}
\item \textbf{Majoration :} comme $u_n \geq \sqrt{3}$, on a $\dfrac{3}{u_n} \leq \sqrt{3}$. De plus $u_n \leq 3$. Donc :
      \[ u_{n+1} = \frac{1}{2}\left(u_n + \frac{3}{u_n}\right) \leq \frac{1}{2}\left(3 + \sqrt{3}\right) = \frac{3+\sqrt{3}}{2} \approx 2{,}37 \leq 3. \]
\item \textbf{Minoration :} par l'inégalité entre moyenne arithmétique et moyenne géométrique (ou directement) : pour $x > 0$,
      \[ \frac{1}{2}\left(x + \frac{3}{x}\right) - \sqrt{3} = \frac{(x-\sqrt{3})^2}{2x} \geq 0, \]
      donc $u_{n+1} \geq \sqrt{3}$.
\end{itemize}
La propriété est héréditaire.

\emph{Conclusion :} pour tout $n \geq 1$, $\sqrt{3} \leq u_n \leq 3$ : la suite est bornée.

\textbf{Étape 2 : monotonie.} Pour $n \geq 1$ :
\[ u_{n+1} - u_n = \frac{1}{2}\left(\frac{3}{u_n} - u_n\right) = \frac{3 - u_n^2}{2u_n}. \]
Comme $u_n \geq \sqrt{3}$, on a $u_n^2 \geq 3$, donc $u_{n+1} - u_n \leq 0$ : la suite est \textbf{décroissante} (à partir du rang 1).

\textbf{Étape 3 : convergence.} La suite $(u_n)$ est décroissante et minorée par $\sqrt{3}$ : d'après le théorème de la limite monotone, elle converge vers une limite $\ell \geq \sqrt{3}$.

\textbf{Étape 4 : calcul de $\ell$.} La fonction $f(x) = \dfrac{1}{2}\left(x + \dfrac{3}{x}\right)$ est continue sur $]0\,;\,+\infty[$ (admis). En passant à la limite dans $u_{n+1} = f(u_n)$ :
\[ \ell = \frac{1}{2}\left(\ell + \frac{3}{\ell}\right) \Longrightarrow 2\ell = \ell + \frac{3}{\ell} \Longrightarrow \ell = \frac{3}{\ell} \Longrightarrow \ell^2 = 3. \]
Donc $\ell = \sqrt{3}$ ou $\ell = -\sqrt{3}$. Or $\ell \geq \sqrt{3} > 0$ (étape 3) : on \textbf{rejette} $-\sqrt{3}$. Finalement :
\[ \boxed{\lim_{n \to +\infty} u_n = \sqrt{3}}. \]

\emph{Vérification numérique :} $u_1 = 2$ ; $u_2 = 1{,}75$ ; $u_3 \approx 1{,}7321$ ; $u_4 \approx 1{,}7320508$. En 4 itérations, on a déjà 7 décimales exactes de $\sqrt{3}$ : la convergence est très rapide !
\end{exoresolu}

\begin{savaistu}[Héron d'Alexandrie et les Babyloniens]
La méthode de l'exercice précédent était connue des Babyloniens il y a près de 4000 ans ! Des tablettes d'argile montrent qu'ils calculaient $\sqrt{2}$ avec une précision remarquable. Héron d'Alexandrie (I\textsuperscript{er} siècle apr. J.-C.) l'a formalisée. Aujourd'hui encore, ce type d'algorithme itératif est au cœur des microprocesseurs : votre calculatrice ne « connaît » pas $\sqrt{3}$, elle le \emph{construit} par une suite récurrente, exactement comme vous venez de le faire.
\end{savaistu}

\courssub{Arbre de décision : comment prouver une convergence ?}

Face à une suite dont on cherche la limite, il faut choisir le bon outil. Le schéma suivant résume la stratégie à adopter.

\begin{popfigure}
\begin{center}
\begin{tikzpicture}[
  node distance=0.55cm and 0.9cm,
  quest/.style={rectangle, rounded corners, draw=popblue, fill=popblueL, text=popdark, align=center, font=\small, inner sep=4pt},
  tool/.style={rectangle, rounded corners, draw=popgreen, fill=popgreenL, text=popdark, align=center, font=\small, inner sep=4pt},
  lbl/.style={font=\footnotesize\itshape, popmuted},
  arr/.style={-{Stealth[length=2.5mm]}, thick, popmuted}
]
\node[quest, align=center] (q1){La suite est-elle\\ monotone et bornée ?};
\node[tool, below left=0.7cm and 0.2cm of q1, align=center] (t1){Théorème de la limite\\ monotone (Th. 2)\\ puis $f(\ell)=\ell$ si récurrente};
\node[quest, below right=0.7cm and 0.2cm of q1, align=center] (q2){Peut-on encadrer $u_n$\\ entre deux suites\\ de même limite ?};
\node[tool, below left=0.7cm and 0.2cm of q2, align=center] (t2){Théorème des\\ gendarmes};
\node[quest, below right=0.7cm and 0.2cm of q2, align=center] (q3){Suite de référence\\ ou opérations\\ sur les limites ?};
\node[tool, below=0.7cm of q3, align=center] (t3){Limites usuelles :\\ $q^n$, $\dfrac{1}{n}$, sommes,\\ produits, quotients};
\draw[arr] (q1) -- node[lbl, above left] {oui} (t1);
\draw[arr] (q1) -- node[lbl, above right] {non} (q2);
\draw[arr] (q2) -- node[lbl, above left] {oui} (t2);
\draw[arr] (q2) -- node[lbl, above right] {non} (q3);
\draw[arr] (q3) -- node[lbl, right] {oui} (t3);
\end{tikzpicture}
\end{center}
\legende{Arbre de décision pour l'étude de la convergence d'une suite : on teste d'abord monotonie + bornitude, puis l'encadrement, puis les limites de référence.}
\end{popfigure}

\courssub{Complément : suites contractantes (non exigible)}

Que faire si la suite récurrente $u_{n+1} = f(u_n)$ n'est \emph{pas} monotone ? Cela arrive : les termes peuvent osciller en se rapprochant de la limite, comme un pendule qui s'amortit. Il existe alors un outil puissant (hors programme, simple évocation) : si la fonction $f$ est \cle{contractante} sur un intervalle stable $I$, c'est-à-dire s'il existe un réel $k$ avec $0 \leq k < 1$ tel que $|f'(x)| \leq k$ pour tout $x \in I$, alors la suite converge vers l'unique point fixe de $f$ dans $I$, et l'erreur est contrôlée : $|u_n - \ell| \leq k^n \, |u_0 - \ell|$. Comme $k^n \to 0$ (suite géométrique de raison dans $[0\,;\,1[$), la convergence est garantie même sans monotonie. Cette idée sera approfondie dans l'enseignement supérieur ; au Probatoire, seuls le théorème de la limite monotone et le théorème des gendarmes sont exigibles.

\begin{aretenir}[L'essentiel de la section]
\begin{itemize}
\item $(u_n)$ \cle{bornée} = majorée \emph{et} minorée : $m \leq u_n \leq M$ pour tout $n$.
\item Toute suite \textbf{convergente est bornée}, mais une suite bornée ne converge pas forcément (contre-exemple : $(-1)^n$).
\item \textbf{Théorème de la limite monotone} : croissante + majorée $\Rightarrow$ convergente vers $\sup(u_n)$ ; décroissante + minorée $\Rightarrow$ convergente vers $\inf(u_n)$.
\item Les inégalités larges passent à la limite ; les inégalités strictes deviennent larges.
\item \textbf{Gendarmes} : $v_n \leq u_n \leq w_n$ avec $v_n, w_n \to \ell$ implique $u_n \to \ell$.
\item Pour $u_{n+1} = f(u_n)$ : bornitude (récurrence) $\to$ monotonie $\to$ convergence (Th. 2) $\to$ limite par $f(\ell) = \ell$, en \emph{rejetant} les solutions incompatibles.
\end{itemize}
\end{aretenir}

\begin{exercice}[Pour s'entraîner]
On considère la suite définie par $u_0 = 2$ et $u_{n+1} = \sqrt{u_n + 2}$ pour tout $n \in \mathbb{N}$.
\begin{enumerate}
\item Démontrer par récurrence que pour tout $n$ : $0 \leq u_n \leq 2$.
\item Étudier la monotonie de $(u_n)$.
\item En déduire que $(u_n)$ converge et déterminer sa limite.
\end{enumerate}
\end{exercice}

\begin{corrige}[Exercice : suite $u_{n+1}=\sqrt{u_n+2}$]
\textbf{1.} Initialisation : $u_0 = 2$, donc $0 \leq u_0 \leq 2$. Hérédité : si $0 \leq u_n \leq 2$, alors $2 \leq u_n + 2 \leq 4$, donc $\sqrt{2} \leq u_{n+1} \leq 2$, et a fortiori $0 \leq u_{n+1} \leq 2$. La propriété est héréditaire, donc vraie pour tout $n$.

\textbf{2.} Pour $u_n \in [0\,;\,2]$ : $u_{n+1}^2 - u_n^2 = u_n + 2 - u_n^2 = -(u_n - 2)(u_n + 1) \geq 0$ car $u_n \leq 2$. Comme les termes sont positifs, $u_{n+1} \geq u_n$ : la suite est \textbf{croissante}.

\textbf{3.} La suite est croissante et majorée par $2$ : par le théorème de la limite monotone, elle converge vers $\ell \in [0\,;\,2]$. La fonction $f(x) = \sqrt{x+2}$ étant continue, $\ell$ vérifie $\ell = \sqrt{\ell + 2}$, soit $\ell^2 = \ell + 2$, c'est-à-dire $\ell^2 - \ell - 2 = 0$, d'où $\ell = 2$ ou $\ell = -1$. Or $\ell \geq u_0 = 2$ (suite croissante), donc $\ell = 2$.
\end{corrige}

\coursec{Calcul de limites : Opérations et suites de référence}

Dans la section précédente, nous avons établi des \emph{théorèmes d'existence} : une suite croissante majorée converge, une suite décroissante minorée converge. Mais ces théorèmes ont une limite (sans jeu de mots !) : ils garantissent qu'une limite $\ell$ existe, sans en donner la \emph{valeur}. Cette section répond à la question pratique : \textbf{comment calculer effectivement la limite d'une suite ?} C'est une compétence centrale au Probatoire et au Baccalauréat technique : presque chaque sujet contient une question du type « déterminer la limite de la suite $(u_n)$ ».

L'idée générale est simple : on mémorise un petit stock de \cle{suites de référence} dont les limites sont connues, puis on apprend à \emph{combiner} ces limites grâce aux opérations (somme, produit, quotient). Quand les règles de calcul ne s'appliquent pas directement, on est face à une \cle{forme indéterminée} qu'il faut « lever » par une transformation algébrique.

\courssub{Les suites de référence}

\begin{propriete}[Suites de référence (admises)]
Soit $\alpha > 0$ un réel, $q$ un réel et $\lambda$ un réel quelconque. On a :
\begin{itemize}
\item $\displaystyle\lim_{n\to+\infty} n^{\alpha} = +\infty$ (croissance des puissances) ;
\item $\displaystyle\lim_{n\to+\infty} \frac{1}{n^{\alpha}} = 0$ ;
\item si $|q| < 1$, $\displaystyle\lim_{n\to+\infty} q^n = 0$ ; si $q > 1$, $\displaystyle\lim_{n\to+\infty} q^n = +\infty$ ; si $q = 1$, $q^n = 1$ pour tout $n$ ;
\item $\displaystyle\lim_{n\to+\infty} \frac{\lambda^n}{n!} = 0$ (la factorielle domine toute exponentielle).
\end{itemize}
\end{propriete}

\begin{aretenir}[Hiérarchie des croissances]
Quand $n$ devient grand, les quantités « gagnent » dans l'ordre suivant : la \cle{factorielle} $n!$ domine l'\cle{exponentielle} $q^n$ (avec $q>1$), qui domine la \cle{puissance} $n^{\alpha}$, qui domine le \cle{logarithme} $\ln n$. Dans un quotient de deux quantités qui tendent vers $+\infty$, c'est le terme \emph{dominant} qui impose la limite.
\end{aretenir}

\begin{popfigure}
\begin{tikzpicture}[scale=1, every node/.style={font=\small}]
\fill[popblueL] (0,0) rectangle (3.4,1.1);
\node at (1.7,0.55) {\textbf{Puissances} : $n^{\alpha}\to +\infty$,\ \ $\dfrac{1}{n^{\alpha}}\to 0$};
\fill[poporangeL] (3.9,0) rectangle (7.6,1.1);
\node at (5.75,0.55) {\textbf{Géométriques} : $|q|<1 \Rightarrow q^n\to 0$};
\node at (5.75,0.18) {$q>1 \Rightarrow q^n\to +\infty$};
\fill[popgreenL] (8.1,0) rectangle (11.6,1.1);
\node at (9.85,0.55) {\textbf{Factorielle} : $\dfrac{\lambda^n}{n!}\to 0$};
\draw[-{Stealth},very thick,popdark] (0.3,-0.5) -- (11.3,-0.5);
\node[below,popdark] at (5.8,-0.55) {\textbf{Priorité croissante} : $\ln n \ \prec\ n^{\alpha}\ \prec\ q^n\ (q>1)\ \prec\ n!$};
\end{tikzpicture}
\legende{Tableau récapitulatif des limites de référence et règle de priorité des croissances.}
\end{popfigure}

\begin{exemplebox}[Limites immédiates]
\begin{itemize}
\item $\displaystyle\lim_{n\to+\infty}\frac{1}{n^2}=0$ ; \quad $\displaystyle\lim_{n\to+\infty}\left(\frac{1}{3}\right)^n = 0$ car $\left|\tfrac13\right|<1$ ;
\item $\displaystyle\lim_{n\to+\infty} 2^n = +\infty$ car $2>1$ ; \quad $\displaystyle\lim_{n\to+\infty}\frac{5^n}{n!}=0$.
\end{itemize}
\end{exemplebox}

\courssub{Opérations sur les limites (théorèmes admis)}

\begin{propriete}[Règles de calcul sur les limites (admises)]
Soient $(u_n)$ et $(v_n)$ deux suites convergentes avec $\lim u_n = \ell$ et $\lim v_n = \ell'$ ($\ell, \ell' \in \mathbb{R}$). Alors :
\begin{itemize}
\item \textbf{Somme :} $\lim (u_n + v_n) = \ell + \ell'$ ;
\item \textbf{Produit :} $\lim (u_n \times v_n) = \ell \times \ell'$ ;
\item \textbf{Quotient :} si $\ell' \neq 0$, alors $\displaystyle\lim \frac{u_n}{v_n} = \frac{\ell}{\ell'}$ ;
\item \textbf{Composition :} si $f$ est une fonction \emph{continue en $\ell$}, alors $\lim f(u_n) = f(\ell)$.
\end{itemize}
Ces règles s'étendent aux limites infinies selon les conventions usuelles : $+\infty + \ell = +\infty$, $+\infty \times \ell = +\infty$ si $\ell>0$, $\dfrac{\ell}{+\infty}=0$, etc.
\end{propriete}

\begin{attention}[Le piège du quotient]
L'erreur la plus fréquente au Probatoire : appliquer la règle du quotient \emph{sans vérifier que la limite du dénominateur est non nulle}. Si $\lim v_n = 0$, la règle du quotient \textbf{ne s'applique pas} : on ne peut rien conclure directement (c'est une forme indéterminée si le numérateur tend aussi vers $0$). Attention à ne pas confondre :
\begin{itemize}
\item « la limite du quotient est $0$ » (conclusion possible : $\dfrac{\ell}{+\infty} = 0$) ;
\item « le dénominateur tend vers $0$ » (situation interdite pour la règle du quotient : il faut transformer l'expression).
\end{itemize}
De même, on n'écrit \emph{jamais} $\dfrac{1}{0}$ : on étudie le signe du dénominateur et on conclut $+\infty$ ou $-\infty$ le cas échéant.
\end{attention}

\begin{exemplebox}[Application directe des opérations]
Soit $u_n = 3 + \dfrac{2}{n}$ et $v_n = 5 - \dfrac{1}{n^2}$. On a $\lim u_n = 3$ et $\lim v_n = 5 \neq 0$. Donc :
\[ \lim_{n\to+\infty} \frac{u_n}{v_n} = \frac{3}{5}, \qquad \lim_{n\to+\infty} (u_n \cdot v_n) = 15. \]
\end{exemplebox}

\courssub{Les formes indéterminées et comment les lever}

\begin{definition}[Forme indéterminée]
On dit qu'on est en présence d'une \cle{forme indéterminée} (FI) lorsque les règles de calcul ne permettent pas de conclure. Les cinq formes indéterminées principales sont :
\[ \infty - \infty, \qquad 0 \times \infty, \qquad \frac{0}{0}, \qquad \frac{\infty}{\infty}, \qquad 1^{\infty}. \]
« Indéterminée » ne signifie pas « sans limite » : cela signifie que \emph{le résultat dépend des expressions précises} et qu'il faut transformer l'écriture pour conclure.
\end{definition}

\begin{methode}[Lever une forme indéterminée]
\begin{enumerate}
\item \textbf{Quotient de polynômes ou de puissances ($\infty/\infty$ ou $0/0$)} : factoriser le \cle{terme de plus haut degré} au numérateur et au dénominateur, simplifier, puis appliquer les règles de calcul.
\item \textbf{Différence avec des racines carrées ($\infty - \infty$)} : multiplier par la \cle{quantité conjuguée} : $a - b = \dfrac{a^2 - b^2}{a+b}$, ce qui fait disparaître les racines au numérateur.
\item \textbf{Produit $0 \times \infty$} : réécrire en quotient ($\frac{0}{1/\infty}$ ou $\frac{\infty}{1/0}$) pour retomber sur $0/0$ ou $\infty/\infty$.
\item \textbf{Forme $1^{\infty}$} : écrire $u_n^{v_n} = e^{v_n \ln u_n}$ et étudier la limite de l'exposant $v_n \ln u_n$ (nécessite les outils du chapitre fonctions logarithmes).
\item \textbf{Quotient de quantités de natures différentes} : utiliser la hiérarchie des croissances (le terme dominant impose la limite).
\end{enumerate}
\end{methode}

\begin{exoresolu}[Limite d'une fraction rationnelle]
Énoncé : déterminer $\displaystyle\lim_{n\to+\infty} u_n$ où $u_n = \dfrac{2n^2 + 3n}{n^2 - 5}$.
\tcblower
\textbf{Solution détaillée.} Le numérateur et le dénominateur tendent vers $+\infty$ : c'est une FI du type $\infty/\infty$. On factorise par le terme dominant $n^2$ :
\[ u_n = \frac{n^2\left(2 + \dfrac{3}{n}\right)}{n^2\left(1 - \dfrac{5}{n^2}\right)} = \frac{2 + \dfrac{3}{n}}{1 - \dfrac{5}{n^2}}. \]
Or $\displaystyle\lim_{n\to+\infty}\frac{3}{n} = 0$ et $\displaystyle\lim_{n\to+\infty}\frac{5}{n^2}=0$. Par somme, le numérateur tend vers $2$ et le dénominateur vers $1 \neq 0$. La règle du quotient s'applique :
\[ \lim_{n\to+\infty} u_n = \frac{2}{1} = 2. \]
\textbf{Remarque} : la limite est le quotient des coefficients dominants, $\frac{2}{1}$. Ce raccourci est valable pour toute fraction rationnelle en $n$.
\end{exoresolu}

\begin{exoresolu}[Limite par la quantité conjuguée]
Énoncé : déterminer $\displaystyle\lim_{n\to+\infty} n\left(\sqrt{n+1} - \sqrt{n}\right)$.
\tcblower
\textbf{Solution détaillée.} On a $\sqrt{n+1} - \sqrt{n} \to 0$ (forme $\infty - \infty$ à l'intérieur) et $n \to +\infty$ : FI du type $0 \times \infty$. On multiplie par la quantité conjuguée :
\[ \sqrt{n+1} - \sqrt{n} = \frac{(\sqrt{n+1}-\sqrt{n})(\sqrt{n+1}+\sqrt{n})}{\sqrt{n+1}+\sqrt{n}} = \frac{(n+1)-n}{\sqrt{n+1}+\sqrt{n}} = \frac{1}{\sqrt{n+1}+\sqrt{n}}. \]
Donc :
\[ n\left(\sqrt{n+1} - \sqrt{n}\right) = \frac{n}{\sqrt{n+1}+\sqrt{n}} = \frac{n}{\sqrt{n}\left(\sqrt{1+\frac1n}+1\right)} = \frac{\sqrt{n}}{\sqrt{1+\frac1n}+1}. \]
Le numérateur tend vers $+\infty$ et le dénominateur tend vers $2$ : par quotient, $\displaystyle\lim_{n\to+\infty} n\left(\sqrt{n+1}-\sqrt{n}\right) = +\infty$.
\textbf{Remarque} : au passage, $\sqrt{n+1}-\sqrt{n} = \dfrac{1}{\sqrt{n+1}+\sqrt{n}} \to 0$ : la suite $v_n = \sqrt{n+1}-\sqrt{n}$ converge vers $0$.
\end{exoresolu}

\courssub{Comparaison, encadrement et équivalents}

\begin{propriete}[Théorème des gendarmes (admis)]
Soient $(u_n)$, $(v_n)$ et $(w_n)$ trois suites telles qu'à partir d'un certain rang : $v_n \le u_n \le w_n$. Si $(v_n)$ et $(w_n)$ convergent vers la \emph{même} limite $\ell$, alors $(u_n)$ converge aussi vers $\ell$.
\end{propriete}

\begin{popfigure}
\begin{tikzpicture}
\begin{axis}[width=\linewidth, height=6.2cm, xmin=0, xmax=20, ymin=-1.4, ymax=1.6,
axis lines=middle, xlabel={$n$}, ylabel={}, samples=60,
xtick={5,10,15,20}, ytick={0}, tick label style={font=\small}]
\addplot[popcurve, domain=1:20] {1.2/sqrt(max(0,x))} node[pos=0.9, above left, popblue, font=\small]{$w_n$};
\addplot[popcurve, domain=1:20] {-1.2/sqrt(max(0,x))} node[pos=0.9, below left, poporange, font=\small]{$v_n$};
\addplot[popcurve, domain=1:20] {1.2*sin(deg(3*x))/sqrt(max(0,x))} node[pos=0.55, above, popgreen!60!black, font=\small]{$u_n$};
\addplot[popline, dashed, domain=0:20] {0};
\end{axis}
\end{tikzpicture}
\legende{Théorème des gendarmes : $u_n$ est « coincée » entre $v_n$ et $w_n$ qui tendent toutes deux vers $0$ ; elle est forcée de tendre vers $0$ aussi.}
\end{popfigure}

\begin{exemplebox}[Application des gendarmes]
Soit $u_n = \dfrac{\sin n}{n}$. Pour tout $n \ge 1$ : $-\dfrac{1}{n} \le \dfrac{\sin n}{n} \le \dfrac{1}{n}$. Comme $\lim \dfrac{1}{n} = \lim \left(-\dfrac{1}{n}\right) = 0$, le théorème des gendarmes donne $\lim u_n = 0$.
\end{exemplebox}

\begin{propriete}[Suites équivalentes (admis)]
On dit que $(u_n)$ et $(v_n)$ sont \cle{équivalentes}, et on note $u_n \sim v_n$, lorsque $\displaystyle\lim_{n\to+\infty}\frac{u_n}{v_n} = 1$ (ce qui suppose $v_n \neq 0$ pour $n$ assez grand). Si $u_n \ge 0$, $v_n \ge 0$ et $u_n \sim v_n$, alors $(u_n)$ et $(v_n)$ ont le \emph{même comportement limite} : $\lim u_n = \lim v_n$ (finie ou infinie).
\end{propriete}

\begin{exemplebox}[Équivalent d'une fraction rationnelle]
Pour $u_n = \dfrac{2n^2+3n}{n^2-5}$, on a $u_n \sim \dfrac{2n^2}{n^2} = 2$, ce qui redonne $\lim u_n = 2$. Plus généralement, une fraction rationnelle est équivalente au quotient de ses termes dominants.
\end{exemplebox}

\courssub{Application aux suites récurrentes : l'équation $\ell = f(\ell)$}

\begin{propriete}[Point fixe d'une suite récurrente (admis)]
Soit $(u_n)$ définie par $u_{n+1} = f(u_n)$, où $f$ est \emph{continue}. Si $(u_n)$ converge vers une limite $\ell$, alors $\ell$ vérifie nécessairement :
\[ \ell = f(\ell). \]
On dit que $\ell$ est un \cle{point fixe} de $f$.
\end{propriete}

\begin{methode}[Trouver la limite d'une suite récurrente convergente]
\begin{enumerate}
\item \textbf{Prouver d'abord la convergence} (par exemple : suite croissante et majorée, voir section 3). Sans cette étape, la méthode est invalide !
\item Poser $\ell = \lim u_n$ et passer à la limite dans $u_{n+1} = f(u_n)$ : par continuité de $f$, on obtient $\ell = f(\ell)$.
\item Résoudre l'équation $\ell = f(\ell)$. Si elle a plusieurs solutions, \textbf{choisir la bonne racine} à l'aide de la monotonie et de la bornitude de la suite (par exemple : si $u_n \ge 2$ pour tout $n$, alors $\ell \ge 2$).
\end{enumerate}
\end{methode}

\begin{exoresolu}[Limite d'une suite récurrente]
Énoncé : soit $(u_n)$ définie par $u_0 = 1$ et $u_{n+1} = \sqrt{u_n + 2}$. On admet que $(u_n)$ est croissante et majorée par $2$. Déterminer sa limite.
\tcblower
\textbf{Solution détaillée.} La suite est croissante et majorée par $2$ : elle converge vers un réel $\ell \in [1\,;\,2]$ (car $u_n \ge u_0 = 1$). La fonction $f(x) = \sqrt{x+2}$ est continue sur $[0\,;\,+\infty[$. En passant à la limite :
\[ \ell = \sqrt{\ell + 2} \iff \ell^2 = \ell + 2 \iff \ell^2 - \ell - 2 = 0 \iff (\ell - 2)(\ell + 1) = 0. \]
Les solutions sont $\ell = 2$ et $\ell = -1$. Comme $u_n \ge 1$ pour tout $n$, on a $\ell \ge 1$, donc $\ell = -1$ est impossible. \textbf{Conclusion :} $\lim u_n = 2$.
\end{exoresolu}

\begin{attention}[L'équation $\ell = f(\ell)$ ne suffit pas à prouver la convergence]
Résoudre $\ell = f(\ell)$ ne fait que donner les limites \emph{possibles}. Si la suite diverge, l'équation peut quand même avoir des solutions ! Exemple : $u_{n+1} = 2u_n$ avec $u_0 = 1$ donne $\ell = 2\ell$, soit $\ell = 0$... alors que $u_n = 2^n \to +\infty$. Il faut \textbf{toujours justifier la convergence avant} de passer à la limite.
\end{attention}

\begin{savaistu}[Interprétation graphique des points fixes]
L'équation $\ell = f(\ell)$ signifie géométriquement que le point $(\ell\,;\,\ell)$ appartient à la courbe de $f$ : les points fixes de $f$ sont les abscisses des \textbf{intersections de la courbe $y = f(x)$ avec la droite $y = x$} (la première bissectrice). C'est le principe de la construction « en escalier » ou « en toile d'araignée » que l'on trace pour visualiser le comportement d'une suite récurrente : on part de $u_0$ sur l'axe des abscisses, on monte à la courbe, on rejoint la bissectrice, et on recommence. Selon la pente de $f$ au point fixe, la suite converge en spirale, en escalier... ou s'éloigne ! Ce point de vue graphique est à la base des méthodes numériques de résolution d'équations utilisées en bureautique technique et en calcul scientifique.
\end{savaistu}

\begin{exercice}[Entraînement type Probatoire]
\begin{enumerate}
\item Calculer $\displaystyle\lim_{n\to+\infty} \frac{3n^2 - n + 1}{6n^2 + 5}$.
\item Calculer $\displaystyle\lim_{n\to+\infty} \left(\sqrt{n^2 + n} - n\right)$ (utiliser la quantité conjuguée).
\item Soit $(u_n)$ définie par $u_0 = 3$ et $u_{n+1} = \dfrac{u_n + 4}{2}$. On admet que $(u_n)$ est croissante et majorée par $4$. Vérifier que cette information est cohérente avec l'équation $\ell = f(\ell)$ et déterminer la limite.
\end{enumerate}
\end{exercice}

\begin{corrige}[Exercice n°1]
\begin{enumerate}
\item FI $\infty/\infty$. On factorise par $n^2$ : $\dfrac{3 - \frac1n + \frac{1}{n^2}}{6 + \frac{5}{n^2}} \to \dfrac{3}{6} = \dfrac12$.
\item $\sqrt{n^2+n} - n = \dfrac{n}{\sqrt{n^2+n}+n} = \dfrac{1}{\sqrt{1+\frac1n}+1} \to \dfrac12$.
\item $\ell = \dfrac{\ell+4}{2} \iff 2\ell = \ell + 4 \iff \ell = 4$. Cohérent avec « majorée par $4$ » et $u_0=3 < 4$ (suite croissante) : $\lim u_n = 4$.
\end{enumerate}
\end{corrige}

\coursec{Modélisation et applications concrètes (Vie courante \& Technique)}

Après avoir maîtrisé les outils théoriques --- monotonie, bornitude, théorèmes de convergence et calcul de limites --- nous voici au cœur de la puissance des suites : la \textbf{modélisation}. Une suite n'est pas seulement un objet abstrait ; c'est le langage naturel pour décrire tout phénomène qui évolue par \emph{états successifs} : un capital qui grossit mois après mois, une concentration de médicament qui varie prise après prise, une population qui fluctue génération après génération.

La démarche est toujours la même : \textbf{traduire le réel en mathématiques} (identifier $U_0$, la relation de récurrence $U_{n+1}=f(U_n)$, la question posée), \textbf{résoudre le problème mathématique} (monotonie, convergence, formule explicite, recherche de rang), puis \textbf{retraduire la réponse en langage courant} avec son unité.

\courssub{Le modèle universel : la suite arithmético-géométrique $U_{n+1} = a U_n + b$}

C'est le « couteau suisse » de la modélisation en Terminale. Il apparaît dès qu'une quantité évolue par un \textbf{pourcentage} (facteur multiplicatif $a$) auquel s'ajoute ou se retire une \textbf{valeur fixe} (terme additif $b$) à chaque étape.

\begin{methode}[Modéliser un problème d'épargne, de prêt ou de population par une suite AG]
\textbf{Étape 1 : Identifier les paramètres.}
\begin{itemize}
    \item $U_0$ : valeur initiale (capital de départ, dose initiale, population initiale).
    \item $t$ : taux d'évolution par période (ex: 5\% $\rightarrow t=0,05$).
    \item $a = 1+t$ : coefficient multiplicateur (attention : $a=1{,}05$ pour 5\%, pas $0,05$ !).
    \item $b$ : apport constant (versement positif) ou retrait constant (remboursement, prélèvement, négatif).
\end{itemize}
\textbf{Étape 2 : Écrire la récurrence.} $U_{n+1} = a U_n + b$.
\textbf{Étape 3 : Déterminer la formule explicite (si $a \neq 1$).}
\[ U_n = a^n U_0 + b\,\frac{a^n - 1}{a - 1} \]
\textbf{Étape 4 : Étudier la convergence.}
\begin{itemize}
    \item Si $|a| < 1$ : convergence vers la limite $L = \frac{b}{1-a}$ (équilibre).
    \item Si $a > 1$ : divergence vers $+\infty$ (si $U_0 > L$) ou $-\infty$ (si $U_0 < L$). Pas de limite finie.
    \item Si $a = 1$ : suite arithmétique $U_n = U_0 + nb$ (divergente sauf $b=0$).
\end{itemize}
\textbf{Étape 5 : Répondre à la question concrète} (rang pour dépasser un seuil, capital final, durée de remboursement) \textbf{en redonnant l'unité} (FCFA, kg, individus, années).
\end{methode}

\begin{propriete}[Formule explicite de la suite arithmético-géométrique (admise)]
Soit $(U_n)$ définie par $U_{n+1} = a U_n + b$ avec $a \neq 1$ et $U_0$ donné.
On pose $V_n = U_n - \frac{b}{1-a}$ (la suite centrée sur l'éventuelle limite).
Alors $(V_n)$ est géométrique de raison $a$ et de premier terme $V_0 = U_0 - \frac{b}{1-a}$.
D'où l'explicite : $\quad U_n = a^n \left(U_0 - \frac{b}{1-a}\right) + \frac{b}{1-a} = a^n U_0 + b\,\frac{a^n - 1}{a - 1}$.
\emph{Au Probatoire : on peut utiliser cette formule directement après l'avoir citée. La démonstration par récurrence sur $V_n$ est exigible si le sujet le demande explicitement.}
\end{propriete}

\begin{exemplebox}[Monsieur Kamga : Épargne avec retraits]
Monsieur Kamga place \num{100 000} FCFA dans un compte rapportant 5\% par an. Chaque fin d'année, il retire \num{2 000} FCFA pour les frais de scolarité.
On note $U_n$ le capital (en milliers de FCFA) après $n$ ans. $U_0 = 100$.
\begin{enumerate}
    \item Modéliser : $a = 1{,}05$, $b = -2$. Récurrence : $U_{n+1} = 1{,}05 U_n - 2$.
    \item Explicite : $U_n = 1{,}05^n \cdot 100 - 2 \cdot \frac{1{,}05^n - 1}{0{,}05} = 1{,}05^n \cdot 100 - 40(1{,}05^n - 1) = 60 \cdot 1{,}05^n + 40$.
    \item Convergence : $a=1{,}05 > 1 \Rightarrow$ la suite diverge vers $+\infty$. Le capital croît indéfiniment car les intérêts ($5\% \times U_n$) finissent par dépasser le retrait fixe de 2.
    \item \textbf{Question :} À partir de quel rang le capital dépasse-t-il 200 (soit \num{200 000} FCFA) ?
    \[ 60 \cdot 1{,}05^n + 40 > 200 \iff 1{,}05^n > \frac{160}{60} = \frac{8}{3} \approx 2{,}667 \]
    \[ n \ln(1{,}05) > \ln(2{,}667) \iff n > \frac{\ln(2{,}667)}{\ln(1{,}05)} \approx \frac{0{,}981}{0{,}0488} \approx 20{,}1 \]
    \textbf{Réponse :} Au bout de \textbf{21 ans} (fin de la 21ème année), le capital dépassera \num{200 000} FCFA.
\end{enumerate}
\end{exemplebox}

\begin{popfigure}
\begin{tikzpicture}
\begin{axis}[
    width=\linewidth, height=0.55\linewidth,
    axis lines=middle,
    xlabel={Années $n$}, ylabel={Capital $U_n$ (kFCFA)},
    xmin=0, xmax=25, ymin=0, ymax=350,
    xtick={0,5,10,15,20,25}, ytick={0,100,200,300},
    grid=major, grid style={popline},
    domain=0:25, samples=26,
]
% Continuous curve for visual guide
\addplot[popmuted, dashed, thick, samples=200, domain=0:25] {60*1.05^x + 40};
% Discrete points (the actual sequence)
\addplot[popblue, only marks, mark=*, mark size=2.5, thick] table {
0 100
1 103
2 106.15
3 109.46
4 112.93
5 116.58
6 120.41
7 124.43
8 128.65
9 133.08
10 137.73
11 142.62
12 147.75
13 153.14
14 158.80
15 164.74
16 170.97
17 177.52
18 184.40
19 191.62
20 199.20
21 207.16
22 215.51
23 224.29
24 233.50
25 243.18
};
% Threshold line
\addplot[poporange, dashed, thick] coordinates {(0,200) (25,200)};
\node[poporange, anchor=south west] at (axis cs:0,200) {Seuil 200 kFCFA};
\end{axis}
\end{tikzpicture}
\legende{Évolution du capital de M. Kamga : croissance exponentielle (suite AG avec $a>1$). Les points discrets sont la réalité ; la courbe continue guide l'œil.}
\end{popfigure}

\begin{attention}[Pièges classiques sur les suites AG]
\begin{itemize}
    \item \textbf{Taux vs Coefficient :} « 5\% d'intérêts » $\Rightarrow$ $a = 1{,}05$ (et non $0{,}05$). « Dépréciation de 10\% » $\Rightarrow$ $a = 0{,}90$.
    \item \textbf{Signe de $b$ :} Apport/Versement $\Rightarrow b > 0$. Retrait/Remboursement/Prélèvement $\Rightarrow b < 0$.
    \item \textbf{Unité oubliée :} Écrire « La limite est 40 » au lieu de « Le capital tend vers \num{40 000} FCFA ». \textbf{Zéro point} à la conclusion au Probatoire.
    \item \textbf{Formule explicite :} Vérifier $a \neq 1$. Si $a=1$, c'est une suite arithmétique : $U_n = U_0 + nb$.
    \item \textbf{Rang $n$ :} $U_n$ est la valeur \emph{après} $n$ périodes. $U_0$ est l'état initial. Ne pas décaler d'un cran.
\end{itemize}
\end{attention}

\courssub{Modèle à convergence assurée : Pharmacocinétique et Pollution ($0 < a < 1$)}

Ici, le coefficient $a$ (souvent noté $p$) est \textbf{strictement compris entre 0 et 1}. Il modélise une \emph{élimination proportionnelle} (le corps élimine 30\% du médicament, le lac évacue 10\% de son polluant) suivie d'un \emph{apport constant} (nouvelle prise, rejet usine). Le système tend vers un \textbf{équilibre} (dose résiduelle, concentration limite).

\begin{exemplebox}[Traitement médical : Dose résiduelle]
Un patient prend \textbf{50 mg} d'un anti-inflammatoire toutes les 6 heures. Son organisme élimine \textbf{30\%} de la substance présente toutes les 6 heures.
On note $U_n$ la quantité de médicament (en mg) dans le sang \emph{juste après} la $n$-ième prise. $U_0 = 50$ (première prise).
\begin{enumerate}
    \item \textbf{Modélisation :} Entre deux prises, 30\% disparaissent $\rightarrow$ il reste 70\% ($p=0{,}7$). Puis on ajoute 50 mg.
    \[ U_{n+1} = 0{,}7 U_n + 50 \quad (U_0 = 50) \]
    \item \textbf{Convergence :} $0 < 0{,}7 < 1 \Rightarrow$ convergence vers $L = \frac{b}{1-a} = \frac{50}{1-0{,}7} = \frac{50}{0{,}3} \approx 166{,}67$ mg.
    \item \textbf{Interprétation :} La quantité de médicament dans le sang \textbf{tend vers 166,7 mg} (dose résiduelle). Elle n'excède jamais cette valeur de manière significative.
    \item \textbf{Explicite :} $U_n = 0{,}7^n \cdot 50 + 50 \cdot \frac{0{,}7^n - 1}{-0{,}3} = \frac{500}{3} - \frac{350}{3} \cdot 0{,}7^n$.
    \item \textbf{Question pratique :} Au bout de combien de prises la dose est-elle à 95\% de la limite ? ($U_n > 0{,}95 \times 166{,}67 \approx 158{,}3$)
    \[ \frac{500}{3} - \frac{350}{3} \cdot 0{,}7^n > 158{,}3 \iff 0{,}7^n < \frac{500 - 475}{350} \approx 0{,}0714 \]
    \[ n \ln(0{,}7) < \ln(0{,}0714) \iff n > \frac{\ln(0{,}0714)}{\ln(0{,}7)} \approx \frac{-2{,}64}{-0{,}357} \approx 7{,}4 \]
    \textbf{Réponse :} Dès le \textbf{rang $n=8$} (soit après la 9ème prise, 48 heures), la concentration est stabilisée à 95\%.
\end{enumerate}
\end{exemplebox}

\begin{popfigure}
\begin{tikzpicture}[
    node distance=1.8cm and 2.5cm,
    block/.style={draw=popblue, thick, fill=popblueL, rounded corners, minimum width=2.8cm, minimum height=1.2cm, align=center},
    arrow/.style={-Stealth, thick, popdark},
    labelstyle/.style={font=\small, text=popdark}
]
\node[block, align=center] (entree){Apport constant\\$q$ (mg, kg, FCFA)};
\node[block, right=of entree, align=center] (systeme){Système\\Élimination proportionnelle $p$\%\\Stock $U_n$};
\node[block, right=of systeme, align=center] (sortie){Sortie\\$(1-p)U_n$ éliminé};

\draw[arrow] (entree) -- node[above, labelstyle] {$+ q$} (systeme);
\draw[arrow] (systeme) -- node[above, labelstyle] {$U_{n+1} = p U_n + q$} (sortie);

% Feedback loop visual
\draw[arrow, poporange] (systeme.north) -- ++(0,1) -| node[above, pos=0.25, labelstyle] {Rétroaction : $p U_n$} (systeme.west);

% Parameters box
\node[draw=poppurple, fill=poppurpleL, rounded corners, align=left, font=\small, below=1.5cm of systeme] (params) {
\textbf{Paramètres :}\\
$p \in ]0,1[$ : fraction conservée (ex: 0,7)\\
$q > 0$ : apport constant par période\\
$U_0$ : stock initial\\
\textbf{Limite (équilibre) :} $L = \frac{q}{1-p}$
};
\end{tikzpicture}
\legende{Schéma bloc de la modélisation pharmacocinétique / pollution : entrée constante, élimination proportionnelle, convergence vers l'équilibre $q/(1-p)$.}
\end{popfigure}

\courssub{Algorithme de Héron : Calcul de $\sqrt{A}$ par suite récurrente}

C'est un exemple historique et algorithmique majeur : construire une suite qui converge \textbf{très vite} (convergence quadratique) vers $\sqrt{A}$ sans utiliser de racine carrée, seulement les 4 opérations.

\begin{propriete}[Suite de Héron (ou de Babylone) pour $\sqrt{A}$]
Soit $A > 0$. La suite définie par :
\[ U_0 > 0 \quad \text{(arbitraire, ex: } U_0 = A \text{ ou } 1\text{)} \qquad U_{n+1} = \frac{1}{2}\left( U_n + \frac{A}{U_n} \right) \]
converge vers $\sqrt{A}$.
\emph{Preuve (idée) :} $U_{n+1}$ est la moyenne arithmétique de $U_n$ et $A/U_n$. Par inégalité AM-GM, $U_{n+1} \ge \sqrt{A}$. La suite est décroissante (pour $n\ge 1$) et minorée par $\sqrt{A}$, donc convergente. La limite $L$ vérifie $L = \frac{1}{2}(L + A/L) \Rightarrow L^2 = A$.
\end{propriete}

\begin{exoresolu}[Calcul de $\sqrt{7}$ par l'algorithme de Héron]
Calculer $\sqrt{7}$ à $10^{-6}$ près.
\emph{Choix :} $U_0 = 7$ (ou $3$, peu importe $>0$).
\[ U_{n+1} = \frac{1}{2}\left( U_n + \frac{7}{U_n} \right) \]
\begin{align*}
U_0 &= 7 \\
U_1 &= \frac{1}{2}(7 + 1) = 4 \\
U_2 &= \frac{1}{2}(4 + 1{,}75) = 2{,}875 \\
U_3 &= \frac{1}{2}(2{,}875 + 2{,}43478\ldots) \approx 2{,}65489 \\
U_4 &\approx \frac{1}{2}(2{,}65489 + 2{,}63664) \approx 2{,}645765 \\
U_5 &\approx \frac{1}{2}(2{,}645765 + 2{,}645751) \approx 2{,}64575131
\end{align*}
$\sqrt{7} \approx 2{,}64575131$. Déjà à $U_4$, on a 5 décimales exactes. \textbf{Convergence quadratique :} le nombre de décimales exactes double à chaque étape.
\end{exoresolu}

\begin{savaistu}[Héron d'Alexandrie (Ier siècle ap. J.-C.)]
Cet ingénieur grec a décrit cet algorithme dans son ouvrage \emph{Metrica}. C'est l'ancêtre direct de la méthode de Newton-Raphson pour résoudre $f(x)=0$ (ici $f(x)=x^2-A$). Les tablettes babyloniennes (1800 av. J.-C.) montrent déjà ce calcul pour $\sqrt{2}$. C'est l'algorithme utilisé encore aujourd'hui dans les processeurs pour l'instruction \texttt{SQRT} (souvent combiné avec une approximation initiale par table de lookup).
\end{savaistu}

\courssub{Complément : La suite logistique $U_{n+1} = r U_n (1 - U_n)$ — Quand le déterminisme devient chaos}

Ce modèle (Pierre-François Verhulst, 1845) décrit une population avec ressources limitées. $U_n \in [0,1]$ représente la population normalisée par la capacité maximale. $r \in [0,4]$ est le taux de croissance intrinsèque. C'est un \textbf{complément de culture scientifique} : pas d'étude analytique exigée au Probatoire, mais une illustration graphique puissante.

\begin{itemize}
    \item $0 < r \le 1$ : Extinction ($U_n \to 0$).
    \item $1 < r \le 3$ : Convergence vers un équilibre stable $L = 1 - 1/r$.
    \item $3 < r < 3,45\ldots$ : Oscillation entre 2 valeurs (cycle d'ordre 2).
    \item $r \approx 3,57\ldots$ : \textbf{Cascade de bifurcations} (cycles 4, 8, 16\ldots).
    \item $r > 3,57\ldots$ : \textbf{Chaos déterministe}. Sensibilité extrême aux conditions initiales (effet papillon). Le comportement est imprévisible à long terme, bien que l'équation soit parfaitement déterministe.
\end{itemize}

\begin{popfigure}
\begin{tikzpicture}
\begin{axis}[
    width=\linewidth, height=0.6\linewidth,
    axis lines=middle,
    xlabel={$U_n$}, ylabel={$U_{n+1}$},
    xmin=0, xmax=1, ymin=0, ymax=1,
    domain=0:1, samples=200,
    grid=major, grid style={popline},
    clip=false
]
% Function curve f(x) = r x (1-x) with r=3.5
\def\r{3.5}
\addplot[popblue, thick] {\r*x*(1-x)} node[right, pos=0.9, popblue] {$y = f(x)$};
% Bisector y=x
\addplot[popdark, dashed] {x} node[above right, pos=0.9, popdark] {$y=x$};
% Cobweb starting from x0=0.3
\def\x{0.3}
\pgfmathsetmacro{\y}{\r*\x*(1-\x)}
\draw[poporange, thick] (\x,0) -- (\x,\y);
\foreach \i in {1,...,30} {
    \pgfmathsetmacro{\y}{\r*\x*(1-\x)}
    \draw[poporange, thick] (\x,\y) -- (\y,\y);
    \pgfmathsetmacro{\x}{\y}
    \draw[poporange, thick] (\x,\x) -- (\x,\y);
}
\node[poporange, anchor=south west] at (axis cs:0.1,0.9) {Toile d'araignée ($r=3,5$, $U_0=0,3$) : Cycle d'ordre 4};
\end{axis}
\end{tikzpicture}
\legende{Toile d'araignée (cobweb plot) pour $r=3,5$. La suite ne converge pas vers un point fixe, mais oscille entre 4 valeurs (cycle d'ordre 4). Illustration du chaos naissant.}
\end{popfigure}

\begin{savaistu}[Le chaos déterministe et la météo]
La suite logistique est le modèle le plus simple exhibant le \textbf{chaos déterministe}. Edward Lorenz (1963) a découvert le même phénomène en simplifiant des équations de convection atmosphérique : « Le battement d'ailes d'un papillon au Brésil peut déclencher une tornade au Texas ». C'est la fin du rêve laplacien de prévisibilité parfaite : même sans hasard, l'avenir peut être imprévisible si on ne connaît pas l'état initial avec une précision infinie.
\end{savaistu}

\courssub{Synthèse : La checklist de modélisation pour le Probatoire}

Face à un énoncé concret (« Une usine rejette... », « Un crédit de... », « Une population de... »), adoptez ce réflexe en 5 étapes :

\begin{enumerate}
    \item \textbf{Nommer la suite :} « Soit $(U_n)$ la quantité ... (unité) à l'instant $n$ (période : an, mois, jour, prise). »
    \item \textbf{Donner $U_0$ :} Lire la valeur initiale dans le texte.
    \item \textbf{Établir la récurrence :} Exprimer $U_{n+1}$ en fonction de $U_n$ en justifiant \emph{chaque terme} (intérêts $\to \times 1{,}05$ ; retrait $\to -200$ ; élimination $\to \times 0{,}7$ ; apport $\to +50$).
    \item \textbf{Analyser mathématiquement :} Monotonie ? Bornitude ? Convergence ? Formule explicite ? Calcul de limite ? Recherche de rang (inéquation + logarithmes) ?
    \item \textbf{Conclure en français :} « Le capital dépassera \num{5 000 000} FCFA au bout de 12 ans. » « La concentration en polluant se stabilise vers 12,5 mg/L. » \textbf{Jamais de conclusion nue sans unité ni phrase.}
\end{enumerate}

\begin{exercice}[Entraînement type Probatoire — Modélisation mixte]
Une commune rurale compte \num{5 000} habitants en 2020 ($n=0$). Chaque année, la population augmente de 2\% (naissance - mortalité) et \num{150} nouveaux habitants s'installent (exode rural).
\begin{enumerate}
    \item Modéliser l'évolution de la population $P_n$ (en habitants) par une suite. Préciser $P_0$ et la relation de récurrence.
    \item Donner l'expression explicite de $P_n$.
    \item La commune prévoit de construire un nouveau lycée quand la population dépassera \num{10 000} habitants. En quelle année cela se produira-t-il ?
    \item Même question si, au lieu de s'installer, \num{150} habitants partent chaque année ($b = -150$). La population tend-elle vers une limite ? Laquelle ?
\end{enumerate}
\end{exercice}

\begin{corrige}[Exercice n°1 — Modélisation mixte]
\begin{enumerate}
    \item $P_0 = 5000$. Taux 2\% $\Rightarrow a = 1{,}02$. Apport $+150 \Rightarrow b = 150$.
    Récurrence : $\boxed{P_{n+1} = 1{,}02 P_n + 150}$.
    \item Formule explicite ($a \neq 1$) :
    $P_n = 1{,}02^n \cdot 5000 + 150 \cdot \frac{1{,}02^n - 1}{0{,}02} = 1{,}02^n \cdot 5000 + 7500(1{,}02^n - 1) = 12500 \cdot 1{,}02^n - 7500$.
    \item On cherche $n$ tel que $P_n > 10000$.
    $12500 \cdot 1{,}02^n - 7500 > 10000 \iff 1{,}02^n > \frac{17500}{12500} = 1{,}4$.
    $n \ln(1{,}02) > \ln(1{,}4) \iff n > \frac{\ln(1{,}4)}{\ln(1{,}02)} \approx \frac{0{,}3365}{0{,}0198} \approx 16{,}99$.
    $n = 17$. Année : $2020 + 17 = \boxed{2037}$.
    \item Si $b = -150$ : $P_{n+1} = 1{,}02 P_n - 150$.
    $a=1{,}02 > 1 \Rightarrow$ divergence. Pas de limite finie.
    Explicite : $P_n = 1{,}02^n \cdot 5000 - 7500(1{,}02^n - 1) = -2500 \cdot 1{,}02^n + 7500$.
    Comme $1{,}02^n \to +\infty$, $P_n \to -\infty$ (mathématiquement). \textbf{Réalistement :} la population devient nulle (extinction) bien avant. Le modèle n'est plus valide quand $P_n$ approche de 0.
\end{enumerate}
\end{corrige}

\begin{aretenir}[Maîtriser la modélisation par suites pour le Probatoire]
\begin{itemize}
    \item \textbf{AG ($a>1$) :} Épargne, crédit, population croissante $\to$ Divergence (sauf cas rares). Chercher un \textbf{rang} (seuil). Utiliser logarithmes.
    \item \textbf{AG ($0<a<1$) :} Pharmacocinétique, pollution, amortissement $\to$ \textbf{Convergence} vers $L = \frac{b}{1-a}$. Calculer la \textbf{limite} (équilibre) et/ou le \textbf{rang} de proximité (ex: 95\% de la limite).
    \item \textbf{Héron :} Algorithme de calcul de racine. Convergence ultra-rapide (quadratique). Savoir écrire la récurrence et itérer 2-3 fois.
    \item \textbf{Logistique :} Culture G. Comprendre que non-linéarité ($U_n^2$) $\Rightarrow$ comportements complexes (bifurcations, chaos).
    \item \textbf{Rédiger :} Phrase d'intro $\to$ Notations $\to$ Récurrence justifiée $\to$ Étude maths $\to$ \textbf{Conclusion phrase + unité}.
\end{itemize}
\end{aretenir}

\coursec{Synthèse et entraînement type Probatoire}

Après avoir modélisé des situations concrètes dans la section précédente, nous abordons maintenant l'épreuve décisive : l'examen du Probatoire / Baccalauréat technique. Cette section n'apporte pas de nouvelles notions, mais vous donne les \cle{clés de la réussite} : comprendre la structure type des sujets, gérer votre temps, rédiger avec la rigueur attendue par les correcteurs et éviter les pièges classiques qui font perdre des points faciles.

\courssub{Structure type d'un sujet « Suites » au Probatoire}

Depuis plusieurs sessions, les sujets de l'enseignement technique suivent une architecture stable en \textbf{trois parties indépendantes mais séquentielles}. Les identifier dès la lecture du sujet vous fait gagner un temps précieux.

\begin{popfigure}
\centering
\begin{tikzpicture}[
    node distance=1.2cm and 1.5cm,
    box/.style={draw=popblue, thick, rounded corners, fill=popblueL, text width=5.5cm, align=center, minimum height=2.8cm, font=\small},
    arrow/.style={-Stealth, thick, popdark},
    title/.style={font=\bfseries\color{popblue}, align=center}
]
% Nodes
\node[box] (part1) {
    \textbf{PARTIE A : Suite explicite}\par\medskip
    $U_n = f(n)$ (rationnelle, racine, etc.)\par
    \textit{Questions types :}\par
    $\bullet$ Calculer $U_0, U_1, U_2$\par
    $\bullet$ Étudier la monotonie ($U_{n+1}-U_n$ ou $U_{n+1}/U_n$)\par
    $\bullet$ Calculer la limite ($+\infty$, $-\infty$, ou réel)\par
    $\bullet$ Résoudre $U_n > A$ (rang seuil)
};
\node[box, right=of part1] (part2) {
    \textbf{PARTIE B : Suite récurrente}\par\medskip
    $V_{n+1} = f(V_n), \quad V_0$ donné\par
    \textit{Questions types :}\par
    $\bullet$ Calculer premiers termes\par
    $\bullet$ Conjecturer monotonie / convergence\par
    $\bullet$ \textbf{Récurrence} : Majoration / Minoration\par
    $\bullet$ \textbf{Récurrence} : Monotonie\par
    $\bullet$ Convergence (Th. croissante majorée)\par
    $\bullet$ Calcul limite $l = f(l)$\par
    $\bullet$ Encadrement / Rang seuil
};
\node[box, below=of part1] (part3) {
    \textbf{PARTIE C : Modélisation}\par\medskip
    Situation concrète (épargne, population, production, chimie...)\par
    \textit{Questions types :}\par
    $\bullet$ Traduire en suite $(W_n)$\par
    $\bullet$ Identifier type (arithmético-géométrique le plus souvent)\par
    $\bullet$ Expliciter $W_n$ (somme géométrique)\par
    $\bullet$ Interpréter : « Au bout de combien d'années... »\par
    $\bullet$ Limite / Seuil / Somme
};
\draw[arrow] (part1.east) -- (part2.west) node[midway, above, font=\tiny\bfseries\color{poporange}] {Indépendantes};
\draw[arrow] (part1.south) -- (part3.north) node[midway, left, font=\tiny\bfseries\color{poporange}] {Réinvestissement possible};
\end{tikzpicture}
\legende{Architecture type d'un sujet « Suites » au Probatoire Technique (3 parties, ~12-15 points au total).}
\end{popfigure}

\courssub{Gestion du temps et stratégie de l'épreuve}

L'exercice « Suites » représente généralement \textbf{30 à 40 minutes} sur les 3h ou 4h de l'épreuve. Voici la répartition idéale :

\begin{center}
\begin{tabularx}{\linewidth}{|l|X|c|}\hline
\textbf{Phase} & \textbf{Actions} & \textbf{Durée} \\ \hline
\rowcolor{popblueL} \textbf{Lecture \& Repérage} & Identifier les 3 parties, noter les valeurs initiales, la fonction $f$. & 5 min \\ \hline
\textbf{Partie A (Explicite)} & Calculs directs, monotonie par différence/quotient, limite (règles usuelles), inéquation. & 10-12 min \\ \hline
\textbf{Partie B (Récurrente)} & \textbf{Cœur du sujet.} Récurrences (majoration $\rightarrow$ monotonie $\rightarrow$ convergence $\rightarrow$ limite). Soigner la rédaction. & 15-18 min \\ \hline
\textbf{Partie C (Modélisation)} & Mise en équation, explicitation (souvent $W_n = a r^n + b$), questions de sens. & 8-10 min \\ \hline
\textbf{Relecture} & Vérifier les conclusions, les encadrements, les unités, le français mathématique. & 3-5 min \\ \hline
\end{tabularx}
\end{center}

\begin{attention}[Priorité absolue : La rédaction de la Récurrence]
Au Probatoire, \textbf{la forme vaut autant que le fond}. Une récurrence mal rédigée (initialisation oubliée, hérédité mal formulée, conclusion absente) perd \textbf{3 à 4 points} sur 5, même si l'idée est juste. Apprenez par cœur le canevas de la \begin{methode}Réussir l'exercice 'Suites' au Probatoire\end{methode} ci-dessous.
\end{attention}

\courssub{La Check-list de rédaction : Le « Français Mathématique » attendu}

Les correcteurs MINESEC sanctionnent l'absence de connecteurs logiques et de phrases complètes. Voici les \cle{mots-clés obligatoires} à insérer dans vos copies.

\begin{popfigure}
\centering
\begin{tikzpicture}[
    node distance=0.8cm and 1cm,
    stepbox/.style={draw=popgreen, thick, rounded corners, fill=popgreenL, text width=7cm, align=left, minimum height=1.8cm, font=\small},
    kwbox/.style={draw=poporange, thick, rounded corners, fill=poporangeL, text width=5cm, align=center, minimum height=1.8cm, font=\small\bfseries},
    arrow/.style={-Stealth, thick, popdark}
]
% Left column: Steps
\node[stepbox, align=center] (s1){
    \textbf{1. Initialisation}\par
    \textit{« Pour $n=0$, on a $U_0 = \dots$ donc $P(0)$ est vraie. »}
};
\node[stepbox, below=of s1] (s2) {
    \textbf{2. Hérédité}\par
    « Soit $k \in \mathbb{N}$. Supposons $P(k)$ vraie (Hérédité).\par
    Montrons $P(k+1)$.\par
    $U_{k+1} = \dots = \dots \ge \dots$ (car $P(k)$)\par
    Donc $P(k+1)$ est vraie. »
};
\node[stepbox, below=of s2] (s3) {
    \textbf{3. Conclusion}\par
    \textit{« Par le principe de récurrence, $P(n)$ est vraie pour tout $n \in \mathbb{N}$. »}
};
% Right column: Keywords
\node[kwbox, right=of s1] (k1) {Mots-clés :\par\color{popdark}{\bfseries Soit $P(n)$ : ...}\par{\bfseries Initialisation}\par{\bfseries Hérédité}\par{\bfseries Conclusion}};
\node[kwbox, right=of s2] (k2) {Connecteurs :\par\color{popdark}{\bfseries Donc, Car, Or,\par Par conséquent,\par Puisque, Comme}};
\node[kwbox, right=of s3] (k3) {Finalisation :\par\color{popdark}{\bfseries Donc la suite est...\par Par le théorème de la suite croissante majorée...\par Passage à la limite : $l = f(l)$\par Donc $\lim U_n = \dots$}};
\draw[arrow] (s1) -- (k1);
\draw[arrow] (s2) -- (k2);
\draw[arrow] (s3) -- (k3);
\end{tikzpicture}
\legende{Infographie : La rédaction idéale au Probatoire. Encadrez vos résultats finaux (\texttt{\textbackslash boxed\{\}} ou trait rouge).}
\end{popfigure}

\courssub{Pièges classiques et stratégies de contournement}

\begin{attention}[Piège n°1 : Conclure la convergence sans l'avoir prouvée]
\textbf{Erreur fréquente :} On pose $l = f(l)$, on résout $l = \sqrt{3+l}$, on trouve $l = \frac{1+\sqrt{13}}{2}$, et on écrit : « \emph{Donc la suite converge vers $\frac{1+\sqrt{13}}{2}$} ».
\medskip

\textbf{Pourquoi c'est faux :} L'équation $l=f(l)$ donne les \emph{seules valeurs possibles} de la limite \textbf{SI la suite converge}. Elle ne prouve pas la convergence.
\medskip

\textbf{La parade :} Toujours écrire : « La suite $(V_n)$ est croissante et majorée (prouvé par récurrence), \textbf{donc elle converge} (Théorème de la suite monotone bornée). \textbf{Passons à la limite} dans $V_{n+1} = f(V_n)$... »
\end{attention}

\begin{attention}[Piège n°2 : Suite définie par $U_{n+1} = f(U_n)$ avec $f$ \textbf{décroissante}]
Si $f$ est décroissante, la suite $(U_n)$ n'est \textbf{pas monotone} (elle oscille). La méthode standard (récurrence simple sur la monotonie) \textbf{échoue}.
\medskip

\textbf{Stratégie (Monotonie de rang 2) :} Étudier les deux sous-suites paires et impaires :
\begin{itemize}
    \item $U_{2n+2} = f(f(U_{2n})) = g(U_{2n})$ avec $g = f \circ f$ \textbf{croissante}.
    \item $U_{2n+3} = g(U_{2n+1})$.
\end{itemize}
On étudie la monotonie de $(U_{2n})$ et $(U_{2n+1})$ séparément (elles sont monotones en sens inverse). Si elles sont bornées, elles convergent vers la \textbf{même limite} $l$ (racine de $l=f(l)$), donc $(U_n)$ converge vers $l$.
\medskip

\textit{Au Probatoire Technique :} Ce cas est rare mais possible. S'il apparaît, la question guide souvent : « Étudier la monotonie de $(U_{2n})$ et $(U_{2n+1)}$ ». Suivez le guide !
\end{attention}

\begin{attention}[Piège n°3 : Confondre « Divergente » et « Tendance vers $+\infty$ »]
Une suite non majorée n'est pas nécessairement de limite $+\infty$ (ex: $U_n = (-1)^n n$).
\textbf{Règle :} Pour écrire $\lim U_n = +\infty$, il faut prouver : $\forall A > 0, \exists N, \forall n \ge N, U_n \ge A$. En pratique : montrer que la suite est \textbf{croissante} et \textbf{non majorée} (ou minorée et décroissante pour $-\infty$).
\end{attention}

\courssub{La calculatrice : Outil de conjecture, pas de preuve}

En examen, le mode \texttt{TABLE} (Casio) ou \texttt{Sequence} (NumWorks) est votre allié pour \textbf{conjecturer} avant de démontrer.

\begin{popfigure}
\centering
\begin{tikzpicture}[
    font=\ttfamily\small,
    screen/.style={draw=popdark, very thick, fill=popmuted!20, rounded corners=3pt, minimum width=12cm, minimum height=7cm},
    header/.style={fill=popblue, text=white, font=\bfseries\small, minimum height=0.7cm},
    cell/.style={minimum height=0.6cm, minimum width=2.5cm, draw=popline, anchor=north west},
    diffcell/.style={cell, fill=poporangeL, font=\bfseries\color{popdark}},
    cursor/.style={draw=poppink, very thick, rounded corners=1pt}
]
\node[screen] (scr) {};
% Header
\node[header, minimum width=2.5cm] at (scr.north west) {n};
\node[header, minimum width=3cm] at ([xshift=2.5cm]scr.north west) {U\_n};
\node[header, minimum width=3.5cm] at ([xshift=5.5cm]scr.north west) {U\_{n+1}-U\_n};
\node[header, minimum width=3cm] at ([xshift=9cm]scr.north west) {U\_{n+1}/U\_n};
% Data rows
\foreach \n/\un/\diff/\quot in {
    0/0.0000/--/--,
    1/1.7321/1.7321/--,
    2/2.1753/0.4432/1.2558,
    3/2.2749/0.0996/1.0458,
    4/2.2956/0.0207/1.0091,
    5/2.3007/0.0051/1.0022,
    6/2.3022/0.0015/1.0007} {
    \pgfmathsetmacro{\y}{-0.7 - \n*0.65}
    \node[cell] at ([xshift=0.1cm, yshift=\y cm]scr.north west) {\n};
    \node[cell] at ([xshift=2.6cm, yshift=\y cm]scr.north west) {\un};
    \node[diffcell] at ([xshift=5.6cm, yshift=\y cm]scr.north west) {\diff};
    \node[cell] at ([xshift=9.1cm, yshift=\y cm]scr.north west) {\quot};
}
% Cursor on row 3
\draw[cursor] ([xshift=0.05cm, yshift=-2.65cm]scr.north west) rectangle ([xshift=11.95cm, yshift=-3.25cm]scr.north west);
\node[font=\tiny\color{poppink}, anchor=north west] at ([xshift=0.1cm, yshift=-3.35cm]scr.north west) {Curseur sur n=3 : diff > 0 $\rightarrow$ Croissante ; quot > 1 $\rightarrow$ Croissante};
\node[font=\bfseries\color{popblue}, anchor=south west] at ([xshift=0.1cm, yshift=0.2cm]scr.south west) {Mode TABLE : $Y_1 = \sqrt{3+X}$ ; $Y_2 = Y_1(X+1)-Y_1(X)$ ; $Y_3 = Y_1(X+1)/Y_1(X)$};
\end{tikzpicture}
\legende{Simulation écran calculatrice (Mode TABLE) pour $V_{n+1}=\sqrt{3+V_n}, V_0=0$. La colonne différence positive conjecture la croissance.}
\end{popfigure}

\begin{methode}[Utilisation stratégique de la calculatrice]
\begin{enumerate}
    \item \textbf{Programmez la suite} en mode \texttt{SEQ} ou \texttt{TABLE} ($Y_1 = f(X)$).
    \item \textbf{Affichez 15-20 termes} : Observez la tendance (croît ? décroît ? oscille ? plafonne ?).
    \item \textbf{Calculez $U_{n+1}-U_n$ et $U_{n+1}/U_n$} en colonnes supplémentaires (comme sur la figure).
    \item \textbf{Repérez le rang seuil} pour les inéquations ($U_n > 2,3$ $\rightarrow$ $n \ge 4$).
    \item \textbf{Notez vos conjectures} sur brouillon : « Suite croissante, majorée par 2,31, limite $\approx 2,3028$ ».
    \item \textbf{Passez à la démonstration rigoureuse} sur la copie (récurrence, théorèmes). \textbf{Jamais} n'écrivez « D'après ma calculatrice... » sur la copie d'examen.
\end{enumerate}
\end{methode}

\courssub{Sujet type complet corrigé (Entraînement Probatoire)}

Voici un sujet représentatif, mélangeant explicite, récurrente et modélisation légère. Traitez-le en conditions réelles (35 min, sans cours) avant de lire le corrigé.

\begin{exoresolu}[Sujet Type Probatoire — Suites Numériques]
\textbf{Partie A : Suite explicite}\par
Soit la suite $(U_n)$ définie pour tout $n \in \mathbb{N}$ par $U_n = \dfrac{n^2 + 3n}{2n^2 - 1}$.
\begin{enumerate}
    \item Calculer $U_0$, $U_1$, $U_2$.
    \item Étudier la monotonie de la suite $(U_n)$.
    \item Calculer $\displaystyle\lim_{n\to+\infty} U_n$.
    \item Déterminer le plus petit entier $N$ tel que pour tout $n \ge N$, $U_n > 0,49$.
\end{enumerate}

\textbf{Partie B : Suite récurrente}\par
Soit la suite $(V_n)$ définie par $V_0 = 0$ et pour tout $n \in \mathbb{N}$, $V_{n+1} = \sqrt{3 + V_n}$.
\begin{enumerate}
    \item Calculer $V_1$, $V_2$, $V_3$. Conjecturer la monotonie et une majoration.
    \item Démontrer par récurrence que pour tout $n \in \mathbb{N}$, $0 \le V_n \le \frac{1+\sqrt{13}}{2}$.
    \item Démontrer par récurrence que la suite $(V_n)$ est croissante.
    \item En déduire que la suite $(V_n)$ converge. Calculer sa limite.
    \item Déterminer le plus petit entier $n$ tel que $V_n > 2,3$.
\end{enumerate}

\textbf{Partie C : Modélisation — Épargne}\par
Un élève verse 50\,000\,FCFA sur un livret au 1er janvier de chaque année. Le taux d'intérêt annuel est de 4\%, capitalisés en fin d'année. On note $C_n$ le capital disponible juste après le versement de l'année $n$ ($n=0$ correspond à la 1ère année).
\begin{enumerate}
    \item Exprimer $C_{n+1}$ en fonction de $C_n$. Préciser $C_0$.
    \item En déduire que la suite $(C_n - 1\,250\,000)$ est géométrique. En déduire $C_n$ en fonction de $n$.
    \item Au bout de combien d'années le capital dépasse-t-il 300\,000\,FCFA ?
\end{enumerate}
\tcblower
\textbf{Solution détaillée}\par

\underline{\textbf{Partie A}}\par
\textbf{1.} $U_0 = \frac{0}{-1} = 0$. \quad $U_1 = \frac{1+3}{2-1} = 4$. \quad $U_2 = \frac{4+6}{8-1} = \frac{10}{7} \approx 1,43$.\par
\textbf{2.} Pour $n \ge 1$, $2n^2-1 > 0$. Étudions le signe de $U_{n+1} - U_n$ :
\[
U_{n+1} - U_n = \frac{(n+1)^2+3(n+1)}{2(n+1)^2-1} - \frac{n^2+3n}{2n^2-1}
\]
Mise au même dénominateur (positif pour $n\ge 1$) :
\[
\text{Numérateur} = [(n^2+2n+1+3n+3)(2n^2-1)] - [(n^2+3n)(2n^2+4n+1)]
\]
Développement :
$= (n^2+5n+4)(2n^2-1) - (n^2+3n)(2n^2+4n+1)$
$= (2n^4 + 10n^3 + 8n^2 - n^2 -5n -4) - (2n^4 + 4n^3 + n^2 + 6n^3 + 12n^2 + 3n)$
$= 2n^4 + 10n^3 + 7n^2 -5n -4 - 2n^4 - 10n^3 - 13n^2 - 3n$
$= -6n^2 -8n -4 < 0$ pour tout $n \ge 1$.
Donc $U_{n+1} - U_n < 0$ pour $n \ge 1$. La suite est \textbf{décroissante} à partir du rang 1 (et $U_0=0 < U_1=4$, donc pas monotone sur $\mathbb{N}$ entier, mais décroissante sur $\mathbb{N}^*$).\par
\textbf{3.} $\lim_{n\to+\infty} U_n = \lim_{n\to+\infty} \frac{n^2(1+\frac{3}{n})}{n^2(2-\frac{1}{n^2})} = \frac{1}{2}$.\par
\textbf{4.} On cherche $n$ tel que $U_n > 0,49 = \frac{49}{100}$.
$\frac{n^2+3n}{2n^2-1} > \frac{49}{100} \iff 100(n^2+3n) > 49(2n^2-1) \iff 100n^2+300n > 98n^2-49 \iff 2n^2+300n+49 > 0$.
Toujours vrai pour $n \ge 1$. Mais attention : $U_1=4 > 0,49$, $U_2=10/7>0,49$... La suite décroît vers $0,5$. Elle reste toujours $> 0,5 > 0,49$.
Donc $N=1$ (ou $N=0$ si on considère $U_0=0 \not> 0,49$). Le plus petit $N$ tel que $\forall n\ge N, U_n>0,49$ est \textbf{$N=1$}.

\underline{\textbf{Partie B}}\par
\textbf{1.} $V_1 = \sqrt{3} \approx 1,732$. $V_2 = \sqrt{3+\sqrt{3}} \approx 2,175$. $V_3 = \sqrt{3+\sqrt{3+\sqrt{3}}} \approx 2,275$.
Conjecture : Suite croissante, majorée par $2,31$ (environ).\par
\textbf{2.} Notons $L = \frac{1+\sqrt{13}}{2} \approx 2,3028$. C'est la racine positive de $L^2 = 3+L$ ($L^2-L-3=0$).
\emph{Propriété $P(n)$ :} $0 \le V_n \le L$.
\textit{Initialisation :} $n=0$, $V_0=0$. $0 \le 0 \le L$ (vrai car $L>0$).
\textit{Hérédité :} Supposons $P(k)$ vraie : $0 \le V_k \le L$.
Alors $3 \le 3+V_k \le 3+L = L^2$ (car $L^2=3+L$).
La fonction racine carrée est croissante sur $\mathbb{R}^+$, donc :
$\sqrt{3} \le \sqrt{3+V_k} \le \sqrt{L^2} = L$.
Or $V_{k+1} = \sqrt{3+V_k}$. Donc $0 < \sqrt{3} \le V_{k+1} \le L$.
$P(k+1)$ est vraie.
\textit{Conclusion :} Par récurrence, $\forall n \in \mathbb{N}, 0 \le V_n \le L$. La suite est majorée (par $L$) et minorée (par $0$).\par
\textbf{3.} \emph{Propriété $Q(n)$ :} $V_{n+1} \ge V_n$ (croissante).
\textit{Initialisation :} $n=0$, $V_1 = \sqrt{3} \approx 1,73 \ge 0 = V_0$. Vrai.
\textit{Hérédité :} Supposons $Q(k)$ vraie : $V_{k+1} \ge V_k$.
Comme $f(x)=\sqrt{3+x}$ est croissante sur $\mathbb{R}^+$ :
$V_{k+2} = f(V_{k+1}) \ge f(V_k) = V_{k+1}$.
Donc $Q(k+1)$ vraie.
\textit{Conclusion :} Par récurrence, $(V_n)$ est croissante.\par
\textbf{4.} La suite $(V_n)$ est croissante et majorée (par $L$). \textbf{Elle converge} (Théorème de la suite croissante majorée).
Soit $\ell$ sa limite. En passant à la limite dans $V_{n+1} = \sqrt{3+V_n}$ (fonction continue) :
$\ell = \sqrt{3+\ell} \implies \ell^2 = 3+\ell \implies \ell^2 - \ell - 3 = 0$.
$\ell = \frac{1 \pm \sqrt{13}}{2}$. Comme $V_n \ge 0$, $\ell \ge 0$. Donc $\ell = \frac{1+\sqrt{13}}{2}$.
\textbf{5.} D'après le tableau de la calculatrice (ou calcul manuel) :
$V_4 \approx 2,2956 < 2,3$ et $V_5 \approx 2,3008 > 2,3$.
Donc le plus petit $n$ est \textbf{$n=5$}.

\underline{\textbf{Partie C}}\par
\textbf{1.} Intérêts : $0,04 C_n$. Versement : $50\,000$.
$C_{n+1} = C_n + 0,04 C_n + 50\,000 = 1,04 C_n + 50\,000$.
$C_0 = 50\,000$ (juste après le 1er versement).\par
\textbf{2.} On cherche une suite géométrique. Résolvons $x = 1,04x + 50\,000 \iff -0,04x = 50\,000 \iff x = -1\,250\,000$.
Posons $W_n = C_n - (-1\,250\,000) = C_n + 1\,250\,000$.
$W_{n+1} = C_{n+1} + 1\,250\,000 = 1,04 C_n + 50\,000 + 1\,250\,000 = 1,04 C_n + 1\,300\,000$.
Mais $1\,300\,000 = 1,04 \times 1\,250\,000$.
Donc $W_{n+1} = 1,04(C_n + 1\,250\,000) = 1,04 W_n$.
$(W_n)$ est géométrique de raison $1,04$ et premier terme $W_0 = C_0 + 1\,250\,000 = 1\,300\,000$.
$W_n = 1\,300\,000 \times 1,04^n$.
$C_n = 1\,300\,000 \times 1,04^n - 1\,250\,000$.\par
\textbf{3.} $C_n > 300\,000 \iff 1\,300\,000 \times 1,04^n - 1\,250\,000 > 300\,000$
$\iff 1,04^n > \frac{1\,550\,000}{1\,300\,000} = \frac{31}{26} \approx 1,1923$.
$n \ln(1,04) > \ln(31/26) \iff n > \frac{\ln(1,1923)}{\ln(1,04)} \approx \frac{0,176}{0,0392} \approx 4,49$.
Le plus petit entier $n$ est \textbf{5}.
$n=5$ correspond au capital au 1er janvier de l'année 6 (après le 6ème versement), soit \textbf{au bout de 5 ans}.
\end{exoresolu}

\courssub{Auto-évaluation : Grille de critères pour l'examen}

Avant de rendre votre copie, cochez mentalement ces cases. C'est la grille utilisée par les correcteurs.

\begin{center}
\begin{tabularx}{\linewidth}{|>{\bfseries}l|X|c|}\hline
\textbf{Critère} & \textbf{Indicateurs de réussite} & \textbf{Points} \\ \hline
\rowcolor{popblueL} \textbf{Rigueur logique} & Récurrence complète (Init + Hérédité + Conclusion). Théorèmes cités nommément. Pas de saut logique. & /5 \\ \hline
\textbf{Vocabulaire précis} & « Croissante », « Majorée », « Converge », « Limite », « Rang seuil ». Pas de « Monte », « Plafonne », « Tendance ». & /3 \\ \hline
\rowcolor{popgreenL} \textbf{Présentation} & Résultats encadrés. Étapes séparées. Calculs alignés. Unités (FCFA, ans, \%). & /2 \\ \hline
\textbf{Calcul exact} & Pas d'arrondi prématuré. Formes exactes ($\frac{1+\sqrt{13}}{2}$, $1,04^n$) privilégiées. & /3 \\ \hline
\textbf{Modélisation} & Traduction correcte. Interprétation de la réponse finale (phrase). & /2 \\ \hline
\rowcolor{poporangeL} \textbf{TOTAL} & \textbf{Objectif : $\ge 12/15$ pour viser l'excellence.} & /15 \\ \hline
\end{tabularx}
\end{center}

\courssub{Boîte à outils : Fiche Mémo « Monotonie / Bornitude / Limite »}

\begin{center}

    \begin{center}
    \textbf{\large FICHE MEMO — SUITES (TERMINALE TECH)} \\[4pt]
    \textit{À glisser dans le cartable / À mémoriser pour le Probatoire}
    \end{center}
    \vspace{4pt}
    \begin{tabularx}{\linewidth}{|l|X|X|}\hline
    \rowcolor{poppurple!40} \multicolumn{3}{|c|}{\textbf{1. MONOTONIE — Comment prouver ?}} \\ \hline
    \textbf{Suite explicite $U_n=f(n)$} & \textbf{Différence} : $U_{n+1}-U_n$ (signe constant $\rightarrow$ mono). & \textbf{Quotient} : $U_{n+1}/U_n$ (si $U_n>0$, $\ge 1 \rightarrow$ croissante). \\ \hline
    \textbf{Suite récurrente $U_{n+1}=f(U_n)$} & \textbf{Récurrence sur $P(n): U_{n+1} \ge U_n$}. & \textbf{Fonction $f$} : Si $f$ croissante et $U_1 \ge U_0 \rightarrow$ croissante. \\ \hline
    \textbf{Cas $f$ décroissante} & \textbf{Rang 2} : Étudier $U_{2n}$ et $U_{2n+1}$ via $g=f\circ f$ (croissante). & \textit{Attention :} Vérifier que les deux sous-suites ont même limite. \\ \hline
    \rowcolor{poppurple!40} \multicolumn{3}{|c|}{\textbf{2. BORNITUDE — Majorer / Minorer}} \\ \hline
    \textbf{Par récurrence} & Chercher $M$ tel que $f(M) \le M$ (majorant) ou $f(m) \ge m$ (minorant). & Prouver $U_n \le M$ (Init + Hérédité : $U_k \le M \Rightarrow f(U_k) \le f(M) \le M$). \\ \hline
    \textbf{Par fonction} & Si $f$ croissante sur $[m,M]$ et $U_0 \in [m,M]$, alors $\forall n, U_n \in [m,M]$. & Utiliser les bornes de $f$ (ex: $\sqrt{\cdot} \ge 0$, $\frac{1}{1+x^2} \le 1$). \\ \hline
    \rowcolor{poppurple!40} \multicolumn{3}{|c|}{\textbf{3. CONVERGENCE — Les 3 piliers}} \\ \hline
    \textbf{Th 1 : Monotone + Bornée} & Croissante + Majorée $\rightarrow$ Converge. \quad Décroissante + Minorée $\rightarrow$ Converge. & \textbf{Le plus utilisé au Probatoire pour les suites récurrentes.} \\ \hline
    \textbf{Th 2 : Gendarmes} & $g_n \le u_n \le h_n$ et $\lim g_n = \lim h_n = \ell \Rightarrow \lim u_n = \ell$. & Utile pour suites explicites complexes ($|\sin| \le 1$, etc.). \\ \hline
    \textbf{Th 3 : Cauchy (Culture)} & $\forall \varepsilon>0, \exists N, \forall p,q\ge N, |u_p-u_q|<\varepsilon$. & \textit{Pas au programme Terminale Tech, mais base de l'analyse supérieure.} \\ \hline
    \rowcolor{poppurple!40} \multicolumn{3}{|c|}{\textbf{4. CALCUL DE LIMITES — Opérations}} \\ \hline
    \textbf{Formes déterminées} & $l+L$, $l\times L$, $l/L (L\ne 0)$, $\sqrt{l} (l\ge 0)$, $f(l)$ ($f$ continue). & Remplacer $U_n$ par sa limite \textbf{après avoir prouvé la convergence}. \\ \hline
    \textbf{Formes indéterminées} & $\frac{\infty}{\infty}$, $\infty-\infty$, $0\times\infty$, $1^\infty$, $0^0$. & \textbf{Ne JAMAIS écrire « = ? »}. Factoriser, conjuguer, changement de variable, encadrer. \\ \hline
    \textbf{Suites de référence} & $q^n \ (|q|<1 \to 0)$, $n^\alpha \to +\infty$, $\frac{\ln n}{n} \to 0$, $\frac{n^\alpha}{b^n} \to 0 \ (b>1)$. & Comparer les vitesses : $\ln n \ll n^\alpha \ll b^n$. \\ \hline
    \rowcolor{poppurple!40} \multicolumn{3}{|c|}{\textbf{5. RÉDACTION TYPE RÉCURRENCE (Copier-coller mental)}} \\ \hline
    \multicolumn{3}{|l|}{\parbox{0.9\linewidth}{
    \textbf{Soit $P(n)$ la propriété : ``$U_n \le M$'' (ou ``$U_{n+1} \ge U_n$'').}\\
    \textbf{Initialisation :} Pour $n=0$, $U_0 = \dots \le M$, donc $P(0)$ vraie.\\
    \textbf{Hérédité :} Soit $k\in\mathbb{N}$. Supposons $P(k)$ vraie : $U_k \le M$.\\
    Alors $U_{k+1} = f(U_k) \le f(M) \le M$ (car $f$ croissante et $f(M)\le M$).\\
    Donc $P(k+1)$ est vraie.\\
    \textbf{Conclusion :} Par le principe de récurrence, $\forall n\in\mathbb{N}, P(n)$ est vraie. $\blacksquare$
    }} \\ \hline
    \end{tabularx}
\end{center}

\begin{savaistu}[Le correcteur lit votre français]
L'examinateur du Probatoire / Baccalauréat Technique est sensible à la \textbf{qualité de l'expression écrite}. Une démonstration juste mais mal rédigée (phrases nominales, absence de « Donc », « Car », « Or », « Par conséquent ») perd 1 à 2 points sur la note de « Rigueur / Communication ».
\begin{itemize}
    \item Remplacez « $\Rightarrow$ » par \textbf{« Donc »} ou \textbf{« Par conséquent »} dans le texte.
    \item Remplacez « $\because$ » ou « car » (seul) par \textbf{« Car »} ou \textbf{« Puisque »} en début de phrase.
    \item Utilisez \textbf{« Or »} pour introduire un fait décisif (ex: « Or $f$ est croissante sur $\mathbb{R}^+$ »).
    \item Terminez toujours par une \textbf{phrase de conclusion} : « La suite $(V_n)$ converge vers $\frac{1+\sqrt{13}}{2}$ ».
\end{itemize}
\textit{Astuce :} Relisez votre copie en ne lisant que les mots en gras (connecteurs + conclusions). Le fil logique doit tenir debout tout seul.
\end{savaistu}

\courssub{Méthode synthétique : Plan d'attaque en 5 étapes (Révision finale)}

\begin{methode}[Réussir l'exercice « Suites » au Probatoire — Plan d'attaque 5 étapes]
\begin{enumerate}
    \item \textbf{SCANNER} (3 min) : Lire tout le sujet. Identifier : Type (Explicite / Récurrente / Modélisation), Valeurs initiales, Fonction $f$, Questions posées (Monotonie ? Limite ? Rang ? Somme ?).
    \item \textbf{CONJECTURER} (3 min) : Calculatrice en mode TABLE. 10-15 termes. Noter sur brouillon : Croissante/Décroissante ? Majorée/Minorée ? Limite approximative ? Rang seuil visible ?
    \item \textbf{DÉMONTRER PARTIE A (Explicite)} (10 min) : Différence ou Quotient $\rightarrow$ Signe $\rightarrow$ Monotonie. Limite (règles, factorisation). Inéquation $\rightarrow$ Rang $N$. \textbf{Encadrer résultats.}
    \item \textbf{DÉMONTRER PARTIE B (Récurrente)} (15 min) — \textbf{Ordre immuable :}
        \begin{enumerate}
            \item Majoration / Minoration (Récurrence).
            \item Monotonie (Récurrence ou $f$ croissante + $U_1 \ge U_0$).
            \item \textbf{Convergence} (Citer le théorème : « Croissante et majorée $\Rightarrow$ Converge »).
            \item \textbf{Passage à la limite} : $l = f(l)$ $\rightarrow$ Résoudre $\rightarrow$ Choisir la racine compatible avec les bornes.
            \item Rang seuil (Calculatrice ou inéquation).
        \end{enumerate}
    \item \textbf{RÉSOUDRE PARTIE C (Modélisation)} (8 min) : Relation $C_{n+1} = a C_n + b$. Chercher $x_0$ point fixe. Poser $W_n = C_n - x_0$ (Géométrique). Expliciter $C_n$. Question de sens (années, somme, seuil).
\end{enumerate}
\end{methode}

\begin{aretenir}[Derniers conseils pour le jour J]
\begin{itemize}
    \item \textbf{Gérez votre temps :} Si vous bloquez sur une question (ex: monotonie explicite complexe), \textbf{passez à la suite}. La partie récurrente et la modélisation rapportent souvent plus de points et sont plus standardisées.
    \item \textbf{Ne laissez jamais blanc :} Même si vous ne trouvez pas la limite, écrire « La suite est croissante et majorée, donc elle converge » rapporte des points. Écrire « $l = f(l) \Rightarrow l = \dots$ » rapporte des points (même sans convergence prouvée, cela montre la méthode).
    \item \textbf{Unités :} En modélisation, la réponse finale \textbf{doit} avoir une unité (FCFA, ans, mètres, \%) et une phrase : « Le capital dépassera 300\,000 FCFA au bout de 5 ans. »
    \item \textbf{Forme exacte :} Préférez $\frac{1+\sqrt{13}}{2}$ à $2,302...$ sauf si l'énoncé demande un arrondi.
\end{itemize}
\end{aretenir}

\textbf{Fin de la leçon 06 — Suites numériques.} Vous disposez maintenant de l'ensemble des outils théoriques, techniques et stratégiques pour aborder sereinement cette chapitre au Probatoire. L'entraînement régulier sur des sujets d'annales (disponibles sur OnBuch+) reste la meilleure préparation. Bon courage !

\newpage

\coursec{Activité d'intégration}

\coursec{Leçon 06 : Suites numériques — Activité d'intégration : Gestion hydraulique du barrage de Lom Pangar}

\courssub{Objectifs de l'activité}
Cette activité d'intégration vise à \cle{mobiliser l'ensemble des savoirs et savoir-faire} de la leçon sur les suites numériques dans un contexte technique authentique : la gestion hydraulique du barrage de Lom Pangar. Elle permet de valider les compétences suivantes :
\begin{itemize}
    \item Modéliser une situation concrète par deux suites imbriquées : une suite arithmético-géométrique (volume d'eau) et une suite arithmétique (capacité utile).
    \item Étudier la \cle{monotonie}, la \cle{bornitude} et la \cle{convergence} de ces suites.
    \item Calculer les \cle{limites} et les interpréter physiquement (volume maximal théorique, colmatage définitif).
    \item Résoudre une inéquation $V_n > C_n$ par l'outil numérique (mode \texttt{TABLE} de la calculatrice) puis justifier rigoureusement par \cle{encadrement} et formules explicites.
    \item Proposer une \cle{mesure technique corrective} (curage, rehausse) et en simuler l'impact sur la durée de vie de l'ouvrage.
    \item Rédiger un \cle{rapport technique structuré} et assurer une \cle{restitution orale argumentée}.
\end{itemize}

\courssub{Situation problème}
Le barrage hydroélectrique de \textbf{Lom Pangar} (région de l'Est, Cameroun) est un ouvrage stratégique pour la régulation du fleuve Sanaga et la production d'électricité. Sa mise en service a constitué une avancée majeure pour le secteur énergétique national. Cependant, tout réservoir artificiel subit deux phénomènes naturels antagonistes : l'apport en eau (pluies, affluents) et la perte de capacité de stockage (sédimentation, évaporation).

\medskip
\noindent \textbf{Données techniques initiales (année 0 = mise en service) :}
\begin{itemize}
    \item Capacité utile initiale du réservoir : $C_0 = \num{6}$ milliards de $\si{\meter\cubed}$.
    \item Volume d'eau stocké initial : $V_0 = \num{4}$ milliards de $\si{\meter\cubed}$.
\end{itemize}

\noindent \textbf{Évolution annuelle (hypothèses moyennes pluriannuelles) :}
\begin{enumerate}
    \item \textbf{Apports pluviométriques et affluents :} $+\num{0,5}$ milliard de $\si{\meter\cubed}$ par an (volume constant ajouté).
    \item \textbf{Pertes par évaporation et fuites :} $\num{2}\%$ du volume d'eau présent en début d'année (proportionnel à $V_n$).
    \item \textbf{Sédimentation (charriage alluvionnaire) :} Réduction de la capacité utile de $\num{0,05}$ milliard de $\si{\meter\cubed}$ par an (quantité constante soustraite).
\end{enumerate}

\noindent On note $V_n$ le volume d'eau stocké (en milliards de $\si{\meter\cubed}$) et $C_n$ la capacité utile du réservoir (en milliards de $\si{\meter\cubed}$) à la \textbf{fin de l'année $n$} ($n \in \mathbb{N}$).

\begin{popfigure}
\begin{tikzpicture}
\begin{axis}[
    width=\linewidth, height=7cm,
    xlabel={Année $n$}, ylabel={Volume / Capacité (milliards $\si{\meter\cubed}$)},
    xmin=0, xmax=25, ymin=0, ymax=7,
    grid=major, grid style={popline},
    legend pos=north west,
    legend style={fill=popcream, draw=popdark},
    samples=100, domain=0:25,
    tick label style={font=\small},
    label style={font=\small}
]
% V_n explicit: 25 - 21*0.98^x
\addplot[popcurve, popblue, thick, samples=26, mark=*] coordinates {
    (0,4) (1,4.42) (2,4.8316) (3,5.235) (4,5.630) (5,6.017) (6,6.397) (7,6.769) (8,7.134)
};
\addlegendentry{$V_n$ (Volume d'eau)}
% C_n explicit: 6 - 0.05*x
\addplot[popcurve, poporange, thick, samples=26, mark=square*] coordinates {
    (0,6) (1,5.95) (2,5.9) (3,5.85) (4,5.8) (5,5.75) (6,5.7) (7,5.65) (8,5.6)
};
\addlegendentry{$C_n$ (Capacité utile)}
% Horizontal line for V limit
\addplot[dashed, poppurple, thick] coordinates {(0,25) (25,25)};
\addlegendentry{Limite théorique de $V_n$}
% Vertical line at overflow n=5
\addplot[popgreen, thick, dashed, forget plot] coordinates {(5,0) (5,6.5)};
\node[anchor=south, font=\tiny, popgreen] at (axis cs:5,6.5) {Débordement $n=5$};
\end{axis}
\end{tikzpicture}
\legende{Évolution prévisionnelle du volume $V_n$ et de la capacité $C_n$ — Croisement vers $n=5$.}
\end{popfigure}

\courssub{Consignes}
Travailler en binôme. Produire un rapport écrit structuré (Introduction, Modélisation, Analyse, Simulation corrective, Conclusion) et préparer une restitution orale de 5 minutes.

\begin{enumerate}
    \item \textbf{Modélisation :} Exprimer $V_{n+1}$ en fonction de $V_n$ et $C_{n+1}$ en fonction de $C_n$. Identifier la nature (type) de chaque suite. Justifier.
    \item \textbf{Étude de convergence :}
    \begin{enumerate}
        \item Étudier la monotonie et la bornitude de $(V_n)$ et $(C_n)$.
        \item Démontrer la convergence de $(V_n)$ et calculer sa limite $\ell_V$. Interpréter $\ell_V$ physiquement.
        \item Étudier la convergence de $(C_n)$. Quelle est la conséquence à long terme sur la capacité du barrage ?
    \end{enumerate}
    \item \textbf{Risque de débordement :} Déterminer le plus petit entier $n$ tel que $V_n > C_n$ (le barrage déborde).
    \begin{enumerate}
        \item Utiliser le mode \texttt{TABLE} de la calculatrice pour conjecturer la valeur de $n$.
        \item Justifier rigoureusement ce résultat par un \cle{encadrement} (utilisation des formules explicites de $V_n$ et $C_n$).
    \end{enumerate}
    \item \textbf{Mesure corrective et simulation :} Le génie civil propose deux options pour prolonger la durée de vie du barrage :
    \begin{itemize}
        \item \textbf{Option A (Curage annuel) :} Réduire la sédimentation à $\num{0,01}$ milliard $\si{\meter\cubed}$/an (coût annuel élevé).
        \item \textbf{Option B (Rehausse du barrage) :} Augmenter la capacité initiale de $\num{1}$ milliard $\si{\meter\cubed}$ ($C_0 = \num{7}$) sans changer la sédimentation ($\num{0,05}$).
    \end{itemize}
    Pour chaque option, modifier le modèle, recalculer l'année de débordement et comparer l'efficacité. Quelle option recommandez-vous pour une gestion durable sur 20 ans ? Argumenter.
\end{enumerate}

\courssub{Ressources théoriques mobilisables}

\begin{propriete}[Suite arithmético-géométrique]
Une suite $(u_n)$ définie par $u_{n+1} = a u_n + b$ avec $a \neq 1$ est arithmético-géométrique.
Formule explicite : $u_n = a^n (u_0 - \ell) + \ell$ où $\ell = \frac{b}{1-a}$ est le point fixe (limite si $|a|<1$).
Convergence : Si $|a| < 1$, $\lim\limits_{n\to+\infty} u_n = \ell$. Monotonie : même sens que $u_1 - u_0$ si $a>0$.
\end{propriete}

\begin{propriete}[Suite arithmétique]
Une suite $(v_n)$ définie par $v_{n+1} = v_n + r$ est arithmétique de raison $r$.
Explicite : $v_n = v_0 + nr$.
Convergence : Convergente si $r=0$ (constante), divergente vers $+\infty$ si $r>0$, vers $-\infty$ si $r<0$.
\end{propriete}

\begin{propriete}[Théorème de convergence monotone (Probatoire)]
Toute suite croissante et majorée converge vers sa borne supérieure (son majorant le plus petit).
Toute suite décroissante et minorée converge vers sa borne inférieure (son minorant le plus grand).
\end{propriete}

\begin{methode}[Résolution $V_n > C_n$ par encadrement]
\begin{enumerate}
    \item Établir les formules explicites $V_n = \dots$ et $C_n = \dots$.
    \item Poser $f(n) = V_n - C_n$. Étudier le signe de $f(n)$.
    \item Comme $V_n$ croît et $C_n$ décroît, $f(n)$ est strictement croissante.
    \item Trouver $n_0$ tel que $f(n_0) \le 0$ et $f(n_0+1) > 0$. Alors $n_0+1$ est l'année de débordement.
\end{enumerate}
\end{methode}

\begin{attention}[Pièges fréquents]
\begin{itemize}
    \item Confondre "fin d'année $n$" et "début d'année $n+1$" : ici $V_0, C_0$ sont les états initiaux (fin année 0). $V_1$ est l'état fin année 1.
    \item Oublier que la sédimentation réduit la \textbf{capacité} $C_n$, pas le volume $V_n$ directement.
    \item Ne pas vérifier la condition $|a|<1$ avant d'affirmer la convergence d'une suite arithmético-géométrique.
    \item Utiliser la virgule décimale dans le code TikZ/pgfplots (utiliser le point : \texttt{0.98} et non \texttt{0,98}).
    \item Oublier d'énoncer explicitement le théorème utilisé pour conclure à la convergence (suite croissante majorée).
\end{itemize}
\end{attention}

\courssub{Démarche de résolution et corrigé commenté}
\begin{exoresolu}[Corrigé détaillé de l'activité d'intégration — Barrage de Lom Pangar]

\textbf{1. Modélisation et identification des suites}\par
\noindent \textbf{Volume $V_n$ :} Chaque année, on ajoute $\num{0,5}$ (apports) et on retire $\num{2}\%$ de $V_n$ (évaporation/fuites).
\[ V_{n+1} = V_n + \num{0,5} - \num{0,02} V_n = \num{0,98} V_n + \num{0,5}. \]
C'est une \cle{suite arithmético-géométrique} de raison $a = \num{0,98}$ et terme constant $b = \num{0,5}$. Comme $|a| = \num{0,98} < 1$, elle converge.

\noindent \textbf{Capacité $C_n$ :} Chaque année, la sédimentation réduit la capacité de $\num{0,05}$.
\[ C_{n+1} = C_n - \num{0,05}. \]
C'est une \cle{suite arithmétique} de raison $r = -\num{0,05} < 0$. Elle est strictement décroissante et divergente vers $-\infty$ (physiquement : le barrage se colmate totalement).

\textbf{2. Étude de convergence}\par
\paragraph{Suite $(V_n)$ :}
Point fixe (limite candidate) : $\ell_V = \frac{b}{1-a} = \frac{\num{0,5}}{1-\num{0,98}} = \frac{\num{0,5}}{\num{0,02}} = \num{25}$.
Formule explicite : $V_n = \num{0,98}^n (V_0 - \num{25}) + \num{25} = \num{25} - \num{21} \times \num{0,98}^n$.
\begin{itemize}
    \item \textbf{Monotonie :} $V_1 - V_0 = \num{4,42} - \num{4} = \num{0,42} > 0$. Comme $a=\num{0,98} > 0$, la suite est de même variation que le premier terme : \cle{strictement croissante}.
    \item \textbf{Bornitude :} Majorée par sa limite $\num{25}$ (car $\num{0,98}^n > 0 \Rightarrow V_n < \num{25}$). Minorée par $V_0 = \num{4}$.
    \item \textbf{Convergence :} Suite croissante majorée $\Rightarrow$ convergente (théorème de la suite monotone bornée). $\lim\limits_{n\to+\infty} V_n = \ell_V = \num{25}$ milliards $\si{\meter\cubed}$.
    \item \textbf{Interprétation :} Sans sédimentation, le volume d'eau tendrait vers un équilibre de $\num{25}$ milliards $\si{\meter\cubed}$ (apports = pertes). Mais la capacité initiale n'est que de $\num{6}$ !
\end{itemize}

\paragraph{Suite $(C_n)$ :}
Formule explicite : $C_n = C_0 + n r = \num{6} - \num{0,05} n$.
\begin{itemize}
    \item \textbf{Monotonie :} $r = -\num{0,05} < 0 \Rightarrow$ \cle{strictement décroissante}.
    \item \textbf{Bornitude :} Non minorée (tend vers $-\infty$). Majorée par $C_0 = \num{6}$.
    \item \textbf{Convergence :} Divergente vers $-\infty$. Physiquement, la capacité utile s'annule vers $n = \num{120}$ ans ($\num{6} / \num{0,05}$), puis le modèle n'a plus de sens (capacité nulle).
\end{itemize}

\textbf{3. Détermination de l'année de débordement ($V_n > C_n$)}\par
\paragraph{Approche numérique (Calculatrice — Mode TABLE) :}
On entre $u(n) = \num{25} - \num{21} \times \num{0,98}^n$ et $v(n) = \num{6} - \num{0,05} n$.
\begin{center}
\begin{tabular}{|c|c|c|c|}\hline
$n$ & $V_n$ (milliards $\si{\meter\cubed}$) & $C_n$ (milliards $\si{\meter\cubed}$) & $V_n - C_n$ \\ \hline
0 & \num{4,000} & \num{6,000} & \num{-2,000} \\
1 & \num{4,420} & \num{5,950} & \num{-1,530} \\
2 & \num{4,832} & \num{5,900} & \num{-1,068} \\
3 & \num{5,235} & \num{5,850} & \num{-0,615} \\
4 & \num{5,630} & \num{5,800} & \num{-0,170} \\
\textbf{5} & \textbf{\num{6,017}} & \textbf{\num{5,750}} & \textbf{+\num{0,267}} \\ \hline
\end{tabular}
\end{center}
\noindent \textbf{Conjecture :} Le débordement a lieu à la fin de l'année \textbf{$n = 5$} (premier entier où $V_n > C_n$).

\paragraph{Justification rigoureuse par encadrement :}
On pose $D_n = V_n - C_n = (\num{25} - \num{21} \times \num{0,98}^n) - (\num{6} - \num{0,05} n) = \num{19} + \num{0,05} n - \num{21} \times \num{0,98}^n$.
On montre que $(D_n)$ est strictement croissante :
\begin{align*}
D_{n+1} - D_n &= \num{0,05} - \num{21} \times \num{0,98}^n (\num{0,98} - 1) \\
&= \num{0,05} + \num{21} \times \num{0,02} \times \num{0,98}^n \\
&= \num{0,05} + \num{0,42} \times \num{0,98}^n > 0.
\end{align*}
Donc $D_n$ croît strictement.
On a $D_4 = \num{19} + \num{0,2} - \num{21} \times \num{0,98}^4 \approx \num{19,2} - \num{19,37} = -\num{0,17} < 0$.
On a $D_5 = \num{19} + \num{0,25} - \num{21} \times \num{0,98}^5 \approx \num{19,25} - \num{18,98} = +\num{0,27} > 0$.
Par monotonie stricte, $D_n \le 0$ pour $n \le 4$ et $D_n > 0$ pour $n \ge 5$.
\cle{Conclusion :} Le barrage déborde à la fin de la \textbf{5\textsuperscript{e} année} ($n=5$) si aucune mesure n'est prise.

\textbf{4. Simulation des mesures correctives}\par
\paragraph{Option A : Curage annuel (sédimentation réduite à $\num{0,01}$).}
Nouvelle suite capacité : $C_n^{(A)} = \num{6} - \num{0,01} n$.
Volume inchangé : $V_n = \num{25} - \num{21} \times \num{0,98}^n$.
Inéquation : $\num{25} - \num{21} \times \num{0,98}^n > \num{6} - \num{0,01} n \iff \num{19} + \num{0,01} n > \num{21} \times \num{0,98}^n$.
Test $n=4$ : $\num{19,04} > \num{21} \times \num{0,922} \approx \num{19,37}$ (Faux).
Test $n=5$ : $\num{19,05} > \num{21} \times \num{0,904} \approx \num{18,98}$ (Vrai).
\cle{Résultat Option A :} Débordement à $n=5$ \textbf{inchangé} à court terme. Mais à long terme ($n=20$), $C_{20}^{(A)} = \num{5,8}$ vs $C_{20} = \num{5,0}$. Le curage préserve la capacité pour l'avenir.

\paragraph{Option B : Rehausse du barrage ($C_0 = \num{7}$, sédimentation $\num{0,05}$).}
Nouvelle capacité : $C_n^{(B)} = \num{7} - \num{0,05} n$.
Inéquation : $\num{25} - \num{21} \times \num{0,98}^n > \num{7} - \num{0,05} n \iff \num{18} + \num{0,05} n > \num{21} \times \num{0,98}^n$.
Test $n=6$ : $\num{18,3} > \num{21} \times \num{0,886} \approx \num{18,60}$ (Faux).
Test $n=7$ : $\num{18,35} > \num{21} \times \num{0,868} \approx \num{18,23}$ (Vrai).
\cle{Résultat Option B :} Débordement repoussé à $n=7$ (gain de 2 ans).

\paragraph{Analyse comparative sur 20 ans :}
\begin{center}
\begin{tabularx}{\linewidth}{|l|X|X|}\hline
\textbf{Critère} & \textbf{Option A (Curage)} & \textbf{Option B (Rehausse)} \\ \hline
Année débordement & $n=5$ (identique) & $n=7$ (+2 ans) \\ \hline
Capacité à 20 ans & $C_{20}=\num{5,8}$ & $C_{20}=\num{6,0}$ \\ \hline
Coût & Annuel récurrent (exploitation) & Investissement unique (CAPEX) \\ \hline
Durabilité & \textbf{Excellente} (ralentit colmatage) & Limitée (sédimentation continue) \\ \hline
\end{tabularx}
\end{center}

\noindent \textbf{Recommandation technique :} L'\textbf{Option A (Curage)} est la seule durable. Bien qu'elle ne repousse pas le premier débordement immédiat (car le volume $V_n$ croît trop vite vers 25), elle garantit une capacité résiduelle de $\num{5,8}$ à 20 ans contre $\num{5,0}$ sans action, et $\num{6,0}$ pour la rehausse mais qui retombera à $\num{5,0}$ à 40 ans. Une stratégie combinée (Rehausse initiale + Curage régulier) serait optimale pour viser une exploitation sur 50 ans.

\end{exoresolu}

\courssub{Bilan et compétences validées}
Cette activité a permis de \cle{valider l'ensemble des compétences} du chapitre « Suites numériques » dans un cadre professionnel réaliste :
\begin{itemize}
    \item \textbf{Modélisation :} Traduction d'un énoncé technique en relations de récurrence ($V_{n+1}=\num{0,98}V_n+\num{0,5}$, $C_{n+1}=C_n-\num{0,05}$) et identification correcte des types de suites (arithmético-géométrique et arithmétique).
    \item \textbf{Analyse mathématique rigoureuse :} Démonstration de la monotonie (par différence ou quotient), bornitude, convergence via le théorème de la suite croissante majorée / décroissante minorée, calcul des limites ($\ell_V=\num{25}$, divergence de $C_n$).
    \item \textbf{Outils numériques et preuve :} Utilisation pertinente du mode \texttt{TABLE} pour conjecturer, suivie d'une justification par \cle{encadrement} utilisant les formes explicites ($V_n = \num{25} - \num{21} \times \num{0,98}^n$, $C_n = \num{6} - \num{0,05}n$) et la monotonie de la différence $D_n$.
    \item \textbf{Esprit critique et technique :} Comparaison de solutions d'ingénierie (curage vs rehausse) par simulation mathématique, analyse coûts/bénéfices sur l'horizon temporel, rédaction d'une recommandation argumentée.
    \item \textbf{Communication :} Structuration d'un rapport technique (hypothèses, modèles, résultats, conclusion) et préparation d'une soutenance orale.
\end{itemize}

\begin{savaistu}[Le barrage de Lom Pangar : enjeux réels]
Mis en service en 2021, le barrage de Lom Pangar (capacité $\num{6}$ milliards $\si{\meter\cubed}$, hauteur $\num{46}$ m) régule le débit de la Sanaga pour optimiser les centrales d'\textbf{Edéa} et de \textbf{Song Loulou} (production d'aluminium \textsc{Alucam} et électricité nationale). La sédimentation est un défi majeur pour tous les barrages tropicaux : le bassin versant de la Lom charrie des quantités massives d'alluvions lors des crues. Le curage (dragage) est techniquement complexe et coûteux en exploitation ; la rehausse (hausse de la digue) est un investissement lourd (génie civil) mais unique. La modélisation par suites permet aux ingénieurs d'EDC (Electricity Development Corporation) de planifier la maintenance sur 30-50 ans.
\end{savaistu}

\courssub{Entraînement type Probatoire}
\begin{exercice}[Exercice 1 — Suite arithmético-géométrique classique]
On considère la suite $(u_n)$ définie par $u_0 = 3$ et $u_{n+1} = \num{0,8} u_n + 4$ pour tout $n \in \mathbb{N}$.
\begin{enumerate}
    \item Démontrer que la suite $(u_n)$ converge et calculer sa limite.
    \item Déterminer le sens de variation de $(u_n)$.
    \item Trouver le plus petit entier $n$ tel que $u_n > 19,5$.
\end{enumerate}
\end{exercice}

\begin{corrige}[Exercice 1]
\begin{enumerate}
    \item La suite est arithmético-géométrique avec $a=\num{0,8}$ et $b=4$. $|a|<1$, elle converge vers $\ell = \frac{4}{1-\num{0,8}} = 20$.
    \item $u_1 = \num{0,8}\times 3 + 4 = 6,4 > u_0$. Comme $a>0$, la suite est strictement croissante.
    \item Explicite : $u_n = 20 - 17 \times \num{0,8}^n$. Inéquation : $20 - 17 \times \num{0,8}^n > 19,5 \iff 17 \times \num{0,8}^n < 0,5 \iff \num{0,8}^n < \frac{0,5}{17} \approx 0,0294$.
    Par calculatrice (TABLE) ou logarithmes : $n=16$ donne $\approx 0,028$, $n=15$ donne $\approx 0,035$. Le plus petit $n$ est \textbf{16}.
\end{enumerate}
\end{corrige}

\begin{exercice}[Exercice 2 — Modélisation population / épidémie]
Une usine traite des eaux usées. Le volume de boue $B_n$ (en $\si{\meter\cubed}$) évolue selon $B_{n+1} = \num{0,9} B_n + 50$ avec $B_0 = 200$. La capacité du bassin est de $C_n = 600 - 10n$.
\begin{enumerate}
    \item Étudier la convergence de $(B_n)$ et donner sa limite.
    \item Déterminer l'année de saturation ($B_n > C_n$) par encadrement.
\end{enumerate}
\end{exercice}

\begin{corrige}[Exercice 2]
\begin{enumerate}
    \item $a=\num{0,9}, b=50$. $\ell = \frac{50}{0,1} = 500$. Suite croissante ($B_1=230>200$) majorée par 500 $\Rightarrow$ converge vers 500.
    \item $B_n = 500 - 300 \times \num{0,9}^n$. $C_n = 600 - 10n$.
    Inéquation : $500 - 300 \times \num{0,9}^n > 600 - 10n \iff 10n - 100 > 300 \times \num{0,9}^n$.
    $n=10$ : $0 > 300 \times 0,348 \approx 104$ (Faux).
    $n=15$ : $50 > 300 \times 0,205 \approx 61,5$ (Faux).
    $n=16$ : $60 > 300 \times 0,185 \approx 55,5$ (Vrai).
    Saturation à l'année \textbf{16}.
\end{enumerate}
\end{corrige}

\newpage

\coursec{Méthodes et exercices résolus}

\courssub{Raisonnement par récurrence : principe et rédaction}

\begin{methode}[Preuve par récurrence : structure et rédaction]
    \textbf{Objectif :} Démontrer une propriété $\mathcal{P}(n)$ pour tout entier naturel $n \ge n_0$.
    \begin{enumerate}
        \item \textbf{Initialisation :} Vérifier que $\mathcal{P}(n_0)$ est vraie (souvent $n_0 = 0$ ou $1$). Écrire explicitement le calcul.
        \item \textbf{Hérédité :} Supposer $\mathcal{P}(k)$ vraie pour un $k \ge n_0$ arbitraire (hypothèse de récurrence). En déduire rigoureusement que $\mathcal{P}(k+1)$ est vraie.
        \item \textbf{Conclusion :} Écrire la phrase type : « Par le principe de récurrence, la propriété $\mathcal{P}(n)$ est vraie pour tout entier $n \ge n_0$. »
    \end{enumerate}
    \textbf{Réflexes Probatoire :} Ne jamais oublier l'initialisation. Bien distinguer l'hypothèse (supposition) de la thèse (ce qu'on doit prouver). Utiliser $\Rightarrow$ pour l'enchaînement logique dans l'hérédité.
\end{methode}

\begin{exoresolu}[Récurrence : inégalité et somme]
    Soit la suite $(U_n)$ définie par $U_0 = 2$ et $U_{n+1} = 3U_n - 2$ pour tout $n \in \mathbb{N}$.
    \begin{enumerate}
        \item Calculer $U_1$, $U_2$, $U_3$. Conjecturer une expression de $U_n$ en fonction de $n$.
        \item Démontrer par récurrence que pour tout $n \in \mathbb{N}$, $U_n = 3^n + 1$.
        \item En déduire la nature de la suite $(U_n)$ et sa limite.
    \end{enumerate}
    \tcblower
    \textbf{Solution détaillée}
    \begin{enumerate}
        \item \textbf{Calculs :}
        $U_1 = 3 \times 2 - 2 = 4$,
        $U_2 = 3 \times 4 - 2 = 10$,
        $U_3 = 3 \times 10 - 2 = 28$.
        On observe : $2 = 3^0+1$, $4=3^1+1$, $10=3^2+1$, $28=3^3+1$.
        \textbf{Conjecture :} $U_n = 3^n + 1$.
        \item \textbf{Démonstration par récurrence.}
        \textbf{Propriété $\mathcal{P}(n)$ :} « $U_n = 3^n + 1$ ».
        \begin{itemize}
            \item \textbf{Initialisation ($n=0$) :} $U_0 = 2$ et $3^0 + 1 = 1 + 1 = 2$. Donc $\mathcal{P}(0)$ est vraie.
            \item \textbf{Hérédité :} Soit $k \in \mathbb{N}$. Supposons $\mathcal{P}(k)$ vraie, c'est-à-dire $U_k = 3^k + 1$.
            Calculons $U_{k+1}$ en utilisant la relation de récurrence :
            \begin{align*}
                U_{k+1} &= 3U_k - 2 \\
                &= 3(3^k + 1) - 2 \quad \text{(par hypothèse de récurrence)} \\
                &= 3^{k+1} + 3 - 2 \\
                &= 3^{k+1} + 1.
            \end{align*}
            Ainsi $\mathcal{P}(k+1)$ est vraie.
        \end{itemize}
        \textbf{Conclusion :} Par le principe de récurrence, pour tout $n \in \mathbb{N}$, $U_n = 3^n + 1$.
        \item \textbf{Nature et limite :}
        La suite $(3^n)$ est strictement croissante et tend vers $+\infty$.
        Donc $(U_n)$ est strictement croissante et $\lim_{n\to+\infty} U_n = +\infty$.
        La suite diverge vers $+\infty$.
    \end{enumerate}
\end{exoresolu}

\courssub{Monotonie d'une suite : méthodes d'étude}

\begin{methode}[Étude de la monotonie d'une suite]
    \textbf{Trois méthodes selon la forme de la suite :}
    \begin{enumerate}
        \item \textbf{Suite explicite $U_n = f(n)$ :} Étudier le signe de $U_{n+1} - U_n$ (différence) ou le rapport $\frac{U_{n+1}}{U_n}$ (si termes de signe constant).
        \item \textbf{Suite récurrente $U_{n+1} = f(U_n)$ :} Étudier la fonction $f$ sur un intervalle $I$ contenant les termes.
            \begin{itemize}
                \item Si $f$ est croissante sur $I$ : la suite est monotone (même sens que $U_1 - U_0$).
                \item Si $f$ est décroissante sur $I$ : la suite n'est pas monotone (mais les sous-suites pairs/impairs le sont).
            \end{itemize}
        \item \textbf{Preuve par récurrence :} Souvent nécessaire pour valider le sens de variation si on utilise la fonction $f$ (montrer que $U_n \in I$ pour tout $n$).
    \end{enumerate}
    \textbf{Rédaction :} « On étudie le signe de $U_{n+1} - U_n$... » ou « La fonction $f$ est croissante sur $I$ et $U_0 \in I$, par récurrence $U_n \in I$... »
\end{methode}

\begin{exoresolu}[Monotonie : suite récurrente et fonction associée]
    Soit la suite $(V_n)$ définie par $V_0 = \frac{1}{2}$ et $V_{n+1} = \dfrac{2V_n + 1}{V_n + 2}$.
    \begin{enumerate}
        \item Montrer que pour tout $n \in \mathbb{N}$, $0 < V_n < 1$.
        \item Étudier la monotonie de $(V_n)$.
        \item En déduire la convergence de $(V_n)$ et calculer sa limite.
    \end{enumerate}
    \tcblower
    \textbf{Solution détaillée}
    \begin{enumerate}
        \item \textbf{Bornitude par récurrence.}
        Propriété $\mathcal{P}(n)$ : « $0 < V_n < 1$ ».
        \begin{itemize}
            \item $n=0$ : $V_0 = 0,5$, donc $0 < V_0 < 1$. $\mathcal{P}(0)$ vraie.
            \item Hérédité : Supposons $0 < V_k < 1$. La fonction $f(x) = \frac{2x+1}{x+2}$ est définie et croissante sur $]-2, +\infty[$ (car $f'(x) = \frac{3}{(x+2)^2} > 0$).
            Comme $0 < V_k < 1$, on a $f(0) < f(V_k) < f(1)$.
            $f(0) = \frac{1}{2} > 0$ et $f(1) = \frac{3}{3} = 1$.
            Donc $0 < \frac{1}{2} \le V_{k+1} < 1$. $\mathcal{P}(k+1)$ vraie.
        \end{itemize}
        Par récurrence, $\forall n \in \mathbb{N},\ 0 < V_n < 1$.
        \item \textbf{Monotonie.}
        Pour $n \in \mathbb{N}$, $V_{n+1} - V_n = \frac{2V_n+1}{V_n+2} - V_n = \frac{2V_n+1 - V_n(V_n+2)}{V_n+2} = \frac{1 - V_n^2}{V_n+2}$.
        Comme $0 < V_n < 1$, le numérateur $1-V_n^2 > 0$ et le dénominateur $V_n+2 > 0$.
        Donc $V_{n+1} - V_n > 0$. La suite $(V_n)$ est \textbf{strictement croissante}.
        \item \textbf{Convergence et limite.}
        La suite est croissante et majorée par $1$ (d'après 1.). D'après le théorème de la suite monotone bornée, elle converge vers une limite $\ell \in \mathbb{R}$.
        Passage à la limite dans $V_{n+1} = f(V_n)$ : $\ell = f(\ell) = \frac{2\ell+1}{\ell+2}$.
        $\ell(\ell+2) = 2\ell+1 \iff \ell^2 + 2\ell = 2\ell + 1 \iff \ell^2 = 1 \iff \ell = 1 \text{ ou } -1$.
        Comme $V_n > 0$, $\ell \ge 0$. Donc $\ell = 1$.
        \textbf{Conclusion :} $\lim_{n\to+\infty} V_n = 1$.
    \end{enumerate}
\end{exoresolu}

\courssub{Bornitude et Convergence : Théorèmes fondamentaux}

\begin{methode}[Théorème de la suite monotone bornée : application]
    \textbf{Théorème (Exigible au Probatoire) :} Toute suite monotone et bornée converge vers une limite finie.
    \begin{itemize}
        \item Suite croissante + majorée $\Rightarrow$ Converge vers son \textbf{majorant le plus petit} (supremum).
        \item Suite décroissante + minorée $\Rightarrow$ Converge vers son \textbf{minorant le plus grand} (infimum).
    \end{itemize}
    \textbf{Démarche type pour une suite récurrente $U_{n+1}=f(U_n)$ :}
    \begin{enumerate}
        \item Déterminer un intervalle $I$ stable par $f$ ($f(I) \subset I$) contenant $U_0$ (souvent par récurrence).
        \item Étudier la monotonie de $f$ sur $I$ pour en déduire celle de $(U_n)$ (comparer $f(x)$ et $x$).
        \item Vérifier la bornitude (conséquence de $U_n \in I$).
        \item Conclure à la convergence et résoudre $\ell = f(\ell)$ en gardant la racine dans $I$.
    \end{enumerate}
\end{methode}

\begin{exoresolu}[Convergence : suite définie par récurrence (type Probatoire)]
    On considère la suite $(W_n)$ définie par $W_0 = 0$ et $W_{n+1} = \sqrt{W_n + 6}$.
    \begin{enumerate}
        \item Montrer que pour tout $n \in \mathbb{N}$, $0 \le W_n \le 3$.
        \item Montrer que la suite $(W_n)$ est croissante.
        \item En déduire que $(W_n)$ converge et calculer sa limite.
    \end{enumerate}
    \tcblower
    \textbf{Solution détaillée}
    \begin{enumerate}
        \item \textbf{Bornitude par récurrence sur $I=[0,3]$.}
        $\mathcal{P}(n)$ : « $0 \le W_n \le 3$ ».
        \begin{itemize}
            \item $n=0$ : $W_0=0 \in [0,3]$. Vrai.
            \item Hérédité : Supposons $0 \le W_k \le 3$. La fonction $f(x)=\sqrt{x+6}$ est croissante sur $[-6,+\infty[$.
            $f(0) = \sqrt{6} \approx 2,45 \ge 0$ et $f(3) = \sqrt{9} = 3$.
            Donc $0 \le f(0) \le f(W_k) \le f(3) = 3$, soit $0 \le W_{k+1} \le 3$.
        \end{itemize}
        Par récurrence, $\forall n,\ W_n \in [0,3]$.
        \item \textbf{Monotonie.}
        Pour $x \in [0,3]$, comparons $f(x)$ et $x$ : $f(x) - x = \sqrt{x+6} - x$.
        Posons $g(x) = \sqrt{x+6} - x$. $g'(x) = \frac{1}{2\sqrt{x+6}} - 1$.
        Pour $x \in [0,3]$, $\sqrt{x+6} \in [\sqrt{6}, 3]$, donc $\frac{1}{2\sqrt{x+6}} \le \frac{1}{2\sqrt{6}} < 1$, ainsi $g'(x) < 0$.
        $g$ est décroissante sur $[0,3]$. $g(3) = 0$, donc pour tout $x \in [0,3]$, $g(x) \ge 0$, c'est-à-dire $f(x) \ge x$.
        Comme $W_n \in [0,3]$, $W_{n+1} = f(W_n) \ge W_n$.
        La suite $(W_n)$ est \textbf{croissante}.
        \item \textbf{Convergence et limite.}
        La suite est croissante et majorée par $3$. D'après le théorème de la suite monotone bornée, elle converge vers une limite $\ell \in [0,3]$.
        Passage à la limite : $\ell = \sqrt{\ell + 6} \iff \ell^2 = \ell + 6 \iff \ell^2 - \ell - 6 = 0 \iff (\ell-3)(\ell+2)=0$.
        $\ell = 3$ ou $\ell = -2$. Comme $\ell \ge 0$, $\ell = 3$.
        \textbf{Conclusion :} $\lim_{n\to+\infty} W_n = 3$.
    \end{enumerate}
\end{exoresolu}

\courssub{Calcul de limites : Opérations et suites de référence}

\begin{methode}[Calcul de limites : opérations, suites de référence, théorème des gendarmes]
    \textbf{Règles de calcul (formes indéterminées interdites) :}
    \begin{itemize}
        \item Somme/Produit/Quotient : Si $\lim U_n = \ell_1$, $\lim V_n = \ell_2$ (finis), alors $\lim(U_n+V_n)=\ell_1+\ell_2$, $\lim(U_nV_n)=\ell_1\ell_2$, $\lim(U_n/V_n)=\ell_1/\ell_2$ ($\ell_2 \neq 0$).
        \item Puissance : $\lim U_n^{\alpha} = \ell^{\alpha}$.
    \end{itemize}
    \textbf{Suites de référence (à connaître par cœur) :}
    \begin{itemize}
        \item $q^n \to 0$ si $|q|<1$ ; $1$ si $q=1$ ; diverge sinon.
        \item $n^\alpha \to +\infty$ si $\alpha>0$ ; $1$ si $\alpha=0$ ; $0$ si $\alpha<0$.
        \item $\frac{n^\alpha}{b^n} \to 0$ pour tout $\alpha>0, b>1$ (l'exponentielle gagne).
        \item $\frac{\ln n}{n^\alpha} \to 0$ pour $\alpha>0$.
    \end{itemize}
    \textbf{Théorème des gendarmes :} Si $U_n \le V_n \le W_n$ et $\lim U_n = \lim W_n = \ell$, alors $\lim V_n = \ell$.
    \textbf{Technique pour limites de fractions rationnelles :} Factoriser par la plus grande puissance du numérateur et dénominateur.
\end{methode}

\begin{exoresolu}[Calcul de limites : cas classiques et pièges]
    Calculer les limites suivantes (justifier chaque étape) :
    \begin{enumerate}
        \item $\displaystyle \lim_{n\to+\infty} \frac{3n^2 - 5n + 2}{2n^2 + 7}$
        \item $\displaystyle \lim_{n\to+\infty} \left( \sqrt{n^2+n} - n \right)$
        \item $\displaystyle \lim_{n\to+\infty} \frac{2^n + 3^n}{4^n - 5^n}$
        \item $\displaystyle \lim_{n\to+\infty} \frac{\sin(n)}{n}$
    \end{enumerate}
    \tcblower
    \textbf{Solution détaillée}
    \begin{enumerate}
        \item \textbf{Fraction rationnelle.} Plus forte puissance : $n^2$.
        $\frac{3n^2 - 5n + 2}{2n^2 + 7} = \frac{n^2(3 - \frac{5}{n} + \frac{2}{n^2})}{n^2(2 + \frac{7}{n^2})} = \frac{3 - \frac{5}{n} + \frac{2}{n^2}}{2 + \frac{7}{n^2}}$.
        $\lim \frac{5}{n} = 0$, $\lim \frac{2}{n^2}=0$, $\lim \frac{7}{n^2}=0$.
        Par opérations sur les limites : $\frac{3 - 0 + 0}{2 + 0} = \frac{3}{2}$.
        \item \textbf{Forme indéterminée $\infty - \infty$.} Multiplier par le conjugué.
        $\sqrt{n^2+n} - n = \frac{(\sqrt{n^2+n} - n)(\sqrt{n^2+n} + n)}{\sqrt{n^2+n} + n} = \frac{n^2+n - n^2}{\sqrt{n^2+n} + n} = \frac{n}{\sqrt{n^2+n} + n}$.
        Factoriser $n$ au dénominateur : $\frac{n}{n(\sqrt{1+\frac{1}{n}} + 1)} = \frac{1}{\sqrt{1+\frac{1}{n}} + 1}$.
        $\lim \frac{1}{n} = 0 \Rightarrow \lim \sqrt{1+\frac{1}{n}} = 1$.
        Limite = $\frac{1}{1+1} = \frac{1}{2}$.
        \item \textbf{Exponentielles.} Plus fort terme : $5^n$ au dénominateur (base la plus grande en valeur absolue).
        $\frac{2^n + 3^n}{4^n - 5^n} = \frac{5^n\left( (\frac{2}{5})^n + (\frac{3}{5})^n \right)}{5^n\left( (\frac{4}{5})^n - 1 \right)} = \frac{ (\frac{2}{5})^n + (\frac{3}{5})^n }{ (\frac{4}{5})^n - 1 }$.
        Comme $\frac{2}{5}, \frac{3}{5}, \frac{4}{5} \in ]0,1[$, leurs puissances tendent vers 0.
        Limite = $\frac{0+0}{0-1} = 0$.
        \item \textbf{Théorème des gendarmes.} $-1 \le \sin(n) \le 1 \Rightarrow -\frac{1}{n} \le \frac{\sin(n)}{n} \le \frac{1}{n}$ (pour $n>0$).
        $\lim -\frac{1}{n} = 0$ et $\lim \frac{1}{n} = 0$.
        Par encadrement, $\lim \frac{\sin(n)}{n} = 0$.
    \end{enumerate}
\end{exoresolu}

\courssub{Modélisation et applications concrètes (Vie courante \& Technique)}

\begin{methode}[Modélisation technique : suites arithmético-géométriques et intérêts]
    \textbf{Situations types :}
    \begin{itemize}
        \item \textbf{Intérêts composés / Capitalisation :} $C_{n+1} = C_n(1+t) + V$ (versement périodique) ou $C_{n+1} = C_n(1+t)$ (simple). Solution : $C_n = C_0(1+t)^n$ ou suite arithmético-géométrique.
        \item \textbf{Amortissement / Dépréciation :} $V_{n+1} = V_n(1-t)$ ($t$ taux d'amortissement). Suite géométrique de raison $<1$.
        \item \textbf{Population / Chimie (dilution) :} $U_{n+1} = a U_n + b$. Résolution par recherche de forme explicite : $U_n - \ell = a^n(U_0 - \ell)$ où $\ell = \frac{b}{1-a}$ (point fixe).
        \item \textbf{Électricité (condensateur RC) :} $u_{n+1} = u_n + \frac{E - u_n}{RC}$ (discrétisation).
    \end{itemize}
    \textbf{Démarche de résolution :}
    1. Traduire le problème en suite $(U_n)$ et relation de récurrence.
    2. Identifier le type (géométrique, arithmético-géométrique).
    3. Trouver l'expression explicite $U_n = f(n)$.
    4. Répondre à la question (seuil, durée, limite).
\end{methode}

\begin{exoresolu}[Application technique : Amortissement d'une machine-outil]
    Une entreprise achète une machine-outil pour $\num{50000}$~\texteuro. Elle s'amortit de $15\%$ par an. On note $V_n$ sa valeur (en euros) après $n$ années ($V_0 = 50000$).
    \begin{enumerate}
        \item Exprimer $V_{n+1}$ en fonction de $V_n$. Préciser la nature de la suite.
        \item Donner l'expression de $V_n$ en fonction de $n$.
        \item Au bout de combien d'années la valeur descend-elle sous $\num{10000}$~\texteuro ?
        \item L'entreprise décide de vendre la machine quand sa valeur atteint $5\%$ de sa valeur initiale. Déterminer l'année de vente.
    \end{enumerate}
    \tcblower
    \textbf{Solution détaillée}
    \begin{enumerate}
        \item Perte de $15\%$ $\Rightarrow$ conservation de $85\%$.
        $V_{n+1} = V_n \times 0,85$.
        C'est une \textbf{suite géométrique} de raison $q = 0,85$ et premier terme $V_0 = 50000$.
        \item $V_n = V_0 \times q^n = 50000 \times 0,85^n$.
        \item On cherche $n$ tel que $V_n < 10000$.
        $50000 \times 0,85^n < 10000 \iff 0,85^n < 0,2$.
        Comme $0,85 < 1$, la fonction $\ln$ est croissante mais $\ln(0,85) < 0$, l'inégalité s'inverse :
        $n \ln(0,85) < \ln(0,2) \iff n > \frac{\ln(0,2)}{\ln(0,85)}$.
        $\frac{\ln(0,2)}{\ln(0,85)} \approx \frac{-1,609}{-0,1625} \approx 9,9$.
        Le plus petit entier $n$ est $10$.
        \textbf{Réponse :} Au bout de 10 ans (à la fin de la 10ème année).
        \item $5\%$ de $50000 = 2500$~\texteuro.
        $50000 \times 0,85^n = 2500 \iff 0,85^n = 0,05$.
        $n = \frac{\ln(0,05)}{\ln(0,85)} \approx \frac{-2,9957}{-0,1625} \approx 18,43$.
        La valeur passe sous $2500$ au cours de la 19ème année. « Quand sa valeur atteint... » $\Rightarrow$ premier $n$ tel que $V_n \le 2500$. $n=19$.
        \textbf{Réponse :} Lors de la 19ème année.
    \end{enumerate}
\end{exoresolu}

\courssub{Synthèse et entraînement type Probatoire}

\begin{methode}[Étude complète d'une suite récurrente $U_{n+1}=f(U_n)$ : Algorithmique Probatoire]
    \textbf{Plan type « Copier-Coller » pour l'examen :}
    \begin{enumerate}
        \item \textbf{Domaine / Fonction :} Définir $f$, son ensemble de définition $D_f$, sa dérivée, son sens de variation, ses points fixes ($\ell = f(\ell)$).
        \item \textbf{Bornitude (Récurrence) :} Trouver un intervalle $I = [m, M]$ tel que $U_0 \in I$ et $f(I) \subset I$. Prouver $\forall n, U_n \in I$.
        \item \textbf{Monotonie :} Comparer $f(x)$ et $x$ sur $I$ (signe de $f(x)-x$). En déduire le sens de variation de $(U_n)$.
        \item \textbf{Convergence :} Invoquer le théorème de la suite monotone bornée $\Rightarrow$ Convergence vers $\ell$.
        \item \textbf{Limite :} Résoudre $\ell = f(\ell)$ et sélectionner la racine dans $I$ (cohérente avec la monotonie : $\ell \ge U_0$ si croissante, $\ell \le U_0$ si décroissante).
    \end{enumerate}
\end{methode}

\begin{exoresolu}[Synthèse : Étude complète type Probatoire (Bac 2023 adapté)]
    Soit la suite $(S_n)$ définie par $S_0 = 1$ et $S_{n+1} = \frac{1}{2}\left( S_n + \frac{5}{S_n} \right)$.
    \begin{enumerate}
        \item On pose $f(x) = \frac{1}{2}(x + \frac{5}{x})$ sur $]0, +\infty[$. Étudier les variations de $f$. Montrer que $f(x) \ge \sqrt{5}$.
        \item Montrer que pour tout $n \in \mathbb{N}^*$, $S_n \ge \sqrt{5}$.
        \item Montrer que la suite $(S_n)_{n \ge 1}$ est décroissante.
        \item En déduire que $(S_n)$ converge et calculer sa limite.
        \item Interprétation : À quoi sert cette suite ? (Méthode de Héron / Babylone).
    \end{enumerate}
    \tcblower
    \textbf{Solution détaillée}
    \begin{enumerate}
        \item \textbf{Étude de $f$.} $f(x) = \frac{1}{2}(x + 5x^{-1})$.
        $f'(x) = \frac{1}{2}(1 - \frac{5}{x^2}) = \frac{x^2 - 5}{2x^2}$.
        $f'(x) = 0 \Leftrightarrow x = \sqrt{5}$ (car $x>0$).
        \begin{center}
        \begin{tabular}{|c|ccccc|}\hline
        $x$ & $0$ & & $\sqrt{5}$ & & $+\infty$ \\ \hline
        $f'(x)$ & & $-$ & $0$ & $+$ & \\ \hline
        $f$ & $+\infty$ & $\searrow$ & $\sqrt{5}$ & $\nearrow$ & $+\infty$ \\ \hline
        \end{tabular}
        \end{center}
        Minimum en $\sqrt{5}$ : $f(\sqrt{5}) = \frac{1}{2}(\sqrt{5}+\sqrt{5}) = \sqrt{5}$.
        Donc $\forall x > 0, f(x) \ge \sqrt{5}$.
        \item \textbf{Bornitude (Récurrence sur $n \ge 1$).}
        $\mathcal{P}(n)$ : « $S_n \ge \sqrt{5}$ ».
        \textit{Initialisation ($n=1$) :} $S_1 = f(S_0) = f(1) = \frac{1}{2}(1+5) = 3 \ge \sqrt{5} \approx 2,23$. Vrai.
        \textit{Hérédité :} Supposons $S_k \ge \sqrt{5}$ ($k \ge 1$). Alors $S_{k+1} = f(S_k)$. Comme $S_k \ge \sqrt{5} > 0$, $S_k \in [\sqrt{5}, +\infty[$. Sur cet intervalle, $f$ est croissante et $f(x) \ge \sqrt{5}$. Donc $S_{k+1} = f(S_k) \ge f(\sqrt{5}) = \sqrt{5}$.
        Par récurrence, $\forall n \ge 1, S_n \ge \sqrt{5}$.
        \item \textbf{Monotonie pour $n \ge 1$.}
        Pour $n \ge 1$, $S_n \ge \sqrt{5} > 0$.
        $S_{n+1} - S_n = \frac{1}{2}(S_n + \frac{5}{S_n}) - S_n = \frac{1}{2}(\frac{5}{S_n} - S_n) = \frac{5 - S_n^2}{2S_n}$.
        Comme $S_n \ge \sqrt{5}$, $S_n^2 \ge 5 \Rightarrow 5 - S_n^2 \le 0$. Dénominateur $>0$.
        Donc $S_{n+1} - S_n \le 0$. La suite $(S_n)_{n \ge 1}$ est \textbf{décroissante}.
        (Note : $S_0=1 < S_1=3$, la suite n'est pas monotone sur $\mathbb{N}$ tout entier, mais l'est dès le rang 1).
        \item \textbf{Convergence et limite.}
        Pour $n \ge 1$, $(S_n)$ est décroissante et minorée par $\sqrt{5}$.
        D'après le théorème de la suite monotone bornée, elle converge vers une limite $\ell \in [\sqrt{5}, +\infty[$.
        Passage à la limite : $\ell = \frac{1}{2}(\ell + \frac{5}{\ell}) \iff 2\ell = \ell + \frac{5}{\ell} \iff \ell = \frac{5}{\ell} \iff \ell^2 = 5 \iff \ell = \sqrt{5}$ (car $\ell > 0$).
        \item \textbf{Interprétation :} Cette suite permet de calculer $\sqrt{5}$ par approximation successive (Méthode de Héron / Newton pour $x^2-5=0$). C'est un algorithme de racine carrée utilisé en informatique embarquée (calculatrices, CNC).
    \end{enumerate}
\end{exoresolu}

\begin{popfigure}
\begin{tikzpicture}[scale=0.8, >=Stealth]
    % Axes
    \draw[->] (-0.5,0) -- (5,0) node[right] {$x$};
    \draw[->] (0,-0.5) -- (0,5) node[above] {$y$};
    % y=x
    \draw[popmuted, dashed] (0,0) -- (4.5,4.5) node[above right] {$y=x$};
    % y=f(x) = 0.5(x + 5/x)
    \draw[popblue, thick, domain=0.5:4.5, samples=200] plot (\x, {0.5*(\x + 5/\x)}) node[above right] {$y=f(x)$};
    % Cobweb for S0=1 -> S1=3 -> S2~2.33 -> S3~2.238
    \draw[poporange, thick] (1,0) -- (1,3); % vertical to f(1)
    \draw[poporange, thick] (1,3) -- (3,3); % horizontal to y=x
    \draw[poporange, thick] (3,3) -- (3,2.333); % vertical to f(3)
    \draw[poporange, thick] (3,2.333) -- (2.333,2.333); % horizontal
    \draw[poporange, thick] (2.333,2.333) -- (2.333,2.238); % vertical
    \draw[poporange, thick] (2.333,2.238) -- (2.238,2.238); % horizontal
    % Points
    \fill[poporange] (1,0) circle (2pt) node[below] {$S_0=1$};
    \fill[poporange] (3,3) circle (2pt) node[above right] {$S_1=3$};
    \fill[poporange] (2.333,2.333) circle (2pt) node[above left] {$S_2$};
    \fill[poporange] (2.238,2.238) circle (2pt) node[below right] {$S_3$};
    \fill[popgreen] (2.236,2.236) circle (3pt) node[below left] {$\ell=\sqrt{5}$};
    % sqrt(5) line
    \draw[popgreen, dashed] (2.236,0) node[below] {$\sqrt{5}$} -- (2.236,2.236);
    \draw[popgreen, dashed] (0,2.236) node[left] {$\sqrt{5}$} -- (2.236,2.236);
\end{tikzpicture}
\legende{Représentation graphique (toile d'araignée) de la méthode de Héron pour $\sqrt{5}$. La suite « descend l'escalier » vers le point fixe.}
\end{popfigure}

\begin{attention}[Les 5 erreurs qui coûtent des points au Probatoire]
    \begin{enumerate}
        \item \textbf{Oublier l'initialisation dans une preuve par récurrence.} « Supposons vrai pour $k$... » sans avoir vérifié $n=0$ (ou $n=1$) = \textbf{0 point} sur la démonstration.
        \item \textbf{Confondre monotonie de $f$ et monotonie de $(U_n)$.} « $f$ est croissante donc $(U_n)$ est croissante » \textbf{FAUX}. Il faut comparer $f(x)$ et $x$ (ou $U_1$ et $U_0$) ET vérifier que $U_n$ reste dans l'intervalle où $f$ est croissante.
        \item \textbf{Calcul de limite : écrire $\frac{\infty}{\infty} = 1$ ou $\infty - \infty = 0$.} Les formes indéterminées ne se calculent pas « au feeling ». Il faut factoriser, conjuguer, ou utiliser les suites de référence. Sanction : \textbf{0 point} sur la limite.
        \item \textbf{Oublier de sélectionner la bonne racine pour la limite.} Résoudre $\ell = f(\ell)$ donne souvent 2 solutions. Il faut \textbf{impérativement} dire : « Comme la suite est croissante et $U_0=2$, $\ell \ge 2$, donc $\ell = 3$ (et pas $-1$) ». Sans justification = \textbf{point en moins}.
        \item \textbf{Mauvaise manipulation d'inéquation avec logarithme.} $0,85^n < 0,2 \Rightarrow n \ln 0,85 < \ln 0,2$. Comme $\ln 0,85 < 0$, diviser change le sens : $n > \frac{\ln 0,2}{\ln 0,85}$. Oublier le changement de sens = \textbf{réponse fausse} (ex: $n=9$ au lieu de $10$).
    \end{enumerate}
\end{attention}

\newpage

\coursec{Exercices}

\coursec{Suites numériques}

\courssub{Raisonnement par récurrence : principe et rédaction}

\begin{definition}[Principe de récurrence sur $\mathbb{N}$]
Soit $\mathcal{P}(n)$ une propriété dépendant d'un entier naturel $n$. Si l'on vérifie :
\begin{enumerate}
    \item \textbf{Initialisation :} $\mathcal{P}(0)$ est vraie (ou $\mathcal{P}(n_0)$ pour un certain $n_0 \in \mathbb{N}$).
    \item \textbf{Hérédité :} Pour tout $k \in \mathbb{N}$ (avec $k \ge n_0$), $\mathcal{P}(k) \implies \mathcal{P}(k+1)$.
\end{enumerate}
Alors $\mathcal{P}(n)$ est vraie pour tout entier $n \ge n_0$.
\end{definition}

\begin{methode}[Rédiger une démonstration par récurrence]
\begin{enumerate}
    \item \textbf{Annoncer} la propriété $\mathcal{P}(n)$ à démontrer.
    \item \textbf{Initialisation :} Vérifier $\mathcal{P}(n_0)$ par un calcul explicite.
    \item \textbf{Hérédité :}
    \begin{itemize}
        \item Supposer $\mathcal{P}(k)$ vraie pour un $k \ge n_0$ arbitraire (hypothèse de récurrence).
        \item En déduire $\mathcal{P}(k+1)$ par un raisonnement logique (calcul algébrique, inégalité, etc.).
    \end{itemize}
    \item \textbf{Conclusion :} « Par le principe de récurrence, $\mathcal{P}(n)$ est vraie pour tout $n \ge n_0$. »
\end{enumerate}
\end{methode}

\begin{exemplebox}[Somme des $n$ premiers entiers]
Démontrer que $\forall n \in \mathbb{N}^*,\ 1+2+\dots+n = \frac{n(n+1)}{2}$.
\begin{itemize}
    \item $\mathcal{P}(n) : \sum_{k=1}^n k = \frac{n(n+1)}{2}$.
    \item \textbf{Initialisation ($n=1$) :} $1 = \frac{1\times 2}{2} = 1$. Vrai.
    \item \textbf{Hérédité :} Supposons $\mathcal{P}(k)$ vraie. $\sum_{i=1}^{k+1} i = \frac{k(k+1)}{2} + (k+1) = (k+1)(\frac{k}{2}+1) = \frac{(k+1)(k+2)}{2}$. Donc $\mathcal{P}(k+1)$ vraie.
    \item \textbf{Conclusion :} Propriété vraie pour tout $n \in \mathbb{N}^*$.
\end{itemize}
\end{exemplebox}

\begin{attention}[Erreurs fréquentes]
\begin{itemize}
    \item Oublier l'initialisation (ou la faire pour un mauvais indice).
    \item Dire « Supposons la propriété vraie pour $n$ » au lieu de « pour un $k$ arbitraire ».
    \item Prouver $\mathcal{P}(k) \implies \mathcal{P}(k)$ (tautologie) au lieu de $\mathcal{P}(k) \implies \mathcal{P}(k+1)$.
    \item Utiliser la conclusion comme hypothèse (raisonnement circulaire).
\end{itemize}
\end{attention}

\courssub{Monotonie d'une suite : méthodes d'étude}

\begin{definition}[Monotonie]
Soit $(u_n)$ une suite numérique.
\begin{itemize}
    \item $(u_n)$ est \textbf{croissante} si $\forall n \in \mathbb{N},\ u_{n+1} \ge u_n$.
    \item $(u_n)$ est \textbf{strictement croissante} si $\forall n \in \mathbb{N},\ u_{n+1} > u_n$.
    \item $(u_n)$ est \textbf{décroissante} si $\forall n \in \mathbb{N},\ u_{n+1} \le u_n$.
    \item $(u_n)$ est \textbf{constante} si $\forall n \in \mathbb{N},\ u_{n+1} = u_n$.
\end{itemize}
\end{definition}

\begin{propriete}[Critères de monotonie]
\begin{enumerate}
    \item \textbf{Différence :} Étudier le signe de $u_{n+1} - u_n$.
    \item \textbf{Rapport (si $u_n > 0$) :} Étudier le signe de $\frac{u_{n+1}}{u_n} - 1$ (ou comparer $\frac{u_{n+1}}{u_n}$ à $1$).
    \item \textbf{Fonction associée (suite explicite $u_n = f(n)$) :} Si $f$ est dérivable sur $[1, +\infty[$, le signe de $f'$ donne la monotonie.
    \item \textbf{Suite récurrente $u_{n+1} = f(u_n)$ :} Si $f$ est croissante sur un intervalle $I$ contenant tous les $u_n$, la suite est monotone et son sens de variation est celui de $u_1 - u_0$ (ou signe de $f(x)-x$ sur $I$).
\end{enumerate}
\end{propriete}

\begin{methode}[Étude de la monotonie d'une suite récurrente $u_{n+1}=f(u_n)$]
\begin{enumerate}
    \item Déterminer un intervalle $I$ tel que $\forall n,\ u_n \in I$ (souvent par récurrence via bornitude).
    \item Étudier la monotonie de $f$ sur $I$ (signe de $f'$).
    \item Si $f$ est croissante sur $I$ : comparer $u_1$ et $u_0$. Si $u_1 \ge u_0$, la suite est croissante ; sinon décroissante.
    \item Si $f$ est décroissante sur $I$ : la suite n'est pas monotone en général ; étudier les sous-suites de rang pair et impair.
\end{enumerate}
\end{methode}

\begin{exemplebox}[Suite explicite : $u_n = \frac{4n-1}{2n+3}$]
$u_{n+1}-u_n = \frac{4n+3}{2n+5} - \frac{4n-1}{2n+3} = \frac{(4n+3)(2n+3)-(4n-1)(2n+5)}{(2n+5)(2n+3)} = \frac{14}{(2n+5)(2n+3)} > 0$.
La suite est strictement croissante.
\end{exemplebox}

\begin{exemplebox}[Suite récurrente : $v_{n+1} = 0,6 v_n + 4$ avec $v_0=0$]
$f(x)=0,6x+4$ est croissante sur $\mathbb{R}$. $v_1=4 > v_0=0$. Donc $(v_n)$ est croissante.
\end{exemplebox}

\courssub{Bornitude et Convergence : Théorèmes fondamentaux}

\begin{definition}[Bornitude]
\begin{itemize}
    \item $(u_n)$ est \textbf{majorée} si $\exists M \in \mathbb{R},\ \forall n,\ u_n \le M$.
    \item $(u_n)$ est \textbf{minorée} si $\exists m \in \mathbb{R},\ \forall n,\ u_n \ge m$.
    \item $(u_n)$ est \textbf{bornée} si elle est majorée et minorée ($\exists M>0,\ \forall n,\ |u_n| \le M$).
\end{itemize}
\end{definition}

\begin{definition}[Limite finie et infinie]
\begin{itemize}
    \item $\lim_{n\to+\infty} u_n = \ell \in \mathbb{R}$ si $\forall \varepsilon > 0,\ \exists N \in \mathbb{N},\ \forall n \ge N,\ |u_n - \ell| < \varepsilon$.
    \item $\lim_{n\to+\infty} u_n = +\infty$ si $\forall A > 0,\ \exists N \in \mathbb{N},\ \forall n \ge N,\ u_n > A$.
    \item $\lim_{n\to+\infty} u_n = -\infty$ si $\forall A < 0,\ \exists N \in \mathbb{N},\ \forall n \ge N,\ u_n < A$.
\end{itemize}
\end{definition}

\begin{propriete}[Suite convergente $\Rightarrow$ bornée]
Toute suite convergente vers une limite finie est bornée. La réciproque est fausse (ex: $u_n = (-1)^n$).
\end{propriete}

\begin{propriete}[Théorème de la suite monotone bornée -- Exigible au Probatoire]
\begin{itemize}
    \item Une suite \textbf{croissante et majorée} converge vers son \textbf{majorant le plus petit} (supremum).
    \item Une suite \textbf{décroissante et minorée} converge vers son \textbf{minorant le plus grand} (infimum).
\end{itemize}
\end{propriete}

\begin{aretenir}[Opérations sur les limites (formes indéterminées)]
Si $\lim u_n = \ell$, $\lim v_n = \ell'$ (finies) :
$\lim(u_n+v_n)=\ell+\ell'$, $\lim(u_n v_n)=\ell\ell'$, $\lim(u_n/v_n)=\ell/\ell'$ (si $\ell'\neq 0$).
\textbf{Formes indéterminées :} $\infty - \infty$, $0 \times \infty$, $\frac{\infty}{\infty}$, $\frac{0}{0}$, $1^\infty$, $0^0$, $\infty^0$.
\end{aretenir}

\courssub{Calcul de limites : Opérations et suites de référence}

\begin{propriete}[Suites de référence]
\begin{itemize}
    \item $\forall \alpha > 0,\ \lim_{n\to+\infty} \frac{1}{n^\alpha} = 0$.
    \item $\forall a \in \mathbb{R},\ \lim_{n\to+\infty} \frac{a^n}{n!} = 0$.
    \item Si $|q| < 1$, $\lim q^n = 0$ ; si $q > 1$, $\lim q^n = +\infty$.
    \item $\lim_{n\to+\infty} \left(1+\frac{1}{n}\right)^n = e$.
\end{itemize}
\end{propriete}

\begin{methode}[Calcul de limites usuelles]
\begin{enumerate}
    \item \textbf{Suite rationnelle :} Factoriser par la plus grande puissance du numérateur et dénominateur (terme de plus haut degré).
    \item \textbf{Différence de racines :} Multiplier par le conjugué.
    \item \textbf{Suite exponentielle / polynôme :} Factoriser par l'exponentielle de plus grande base.
    \item \textbf{Forme $1^\infty$ :} Ramener à $\left(1+\frac{1}{u_n}\right)^{u_n}$ ou utiliser $\ln$.
\end{enumerate}
\end{methode}

\begin{exemplebox}[Limites types]
\begin{enumerate}
    \item $\frac{3n^2-5n+2}{2n^2+7} \sim \frac{3n^2}{2n^2} \to \frac{3}{2}$.
    \item $\sqrt{n+1}-\sqrt{n} = \frac{1}{\sqrt{n+1}+\sqrt{n}} \to 0$.
    \item $\frac{2^n+3^n}{4^n} = \left(\frac{1}{2}\right)^n + \left(\frac{3}{4}\right)^n \to 0$.
    \item $\frac{n^3+2n}{5n^2-1} \sim \frac{n}{5} \to +\infty$.
    \item $\left(1+\frac{1}{n}\right)^{2n} = \left[\left(1+\frac{1}{n}\right)^n\right]^2 \to e^2$.
\end{enumerate}
\end{exemplebox}

\courssub{Je vérifie mes connaissances}

\begin{exercice}[QCM : Vocabulaire et propriétés de base]
Pour chacune des affirmations suivantes, cocher la bonne réponse (une seule par question). Aucune justification n'est demandée.
\begin{enumerate}
    \item Soit $(u_n)$ une suite numérique. L'affirmation « $(u_n)$ est croissante » signifie :
    \begin{itemize}
        \item[A.] $\forall n \in \mathbb{N},\ u_{n+1} > u_n$
        \item[B.] $\forall n \in \mathbb{N},\ u_{n+1} \ge u_n$
        \item[C.] $\exists n \in \mathbb{N},\ u_{n+1} \ge u_n$
        \item[D.] $\forall n \in \mathbb{N},\ u_{n+1} < u_n$
    \end{itemize}
    \item Une suite $(u_n)$ converge vers $+\infty$ si :
    \begin{itemize}
        \item[A.] $\forall A > 0,\ \exists N \in \mathbb{N},\ \forall n \ge N,\ u_n > A$
        \item[B.] $\exists A > 0,\ \forall N \in \mathbb{N},\ \exists n \ge N,\ u_n > A$
        \item[C.] $\forall A > 0,\ \exists N \in \mathbb{N},\ \forall n \le N,\ u_n > A$
        \item[D.] $\exists A > 0,\ \forall n \in \mathbb{N},\ u_n > A$
    \end{itemize}
    \item Soit $(u_n)$ une suite convergente vers $\ell \in \mathbb{R}$. Alors :
    \begin{itemize}
        \item[A.] $(u_n)$ est nécessairement monotone.
        \item[B.] $(u_n)$ est nécessairement bornée.
        \item[C.] $(u_n)$ est nécessairement constante à partir d'un certain rang.
        \item[D.] $(|u_n|)$ diverge.
    \end{itemize}
    \item Pour étudier la monotonie d'une suite $(u_n)$ définie par $u_{n+1} = f(u_n)$ avec $f$ croissante sur un intervalle $I$, on compare généralement :
    \begin{itemize}
        \item[A.] $u_1$ et $u_0$
        \item[B.] $f(x)$ et $x$ sur $I$
        \item[C.] $u_n$ et $\ell$
        \item[D.] $f'(x)$ et $0$
    \end{itemize}
\end{enumerate}
\end{exercice}

\begin{exercice}[Vrai / Faux justifié]
Dire si les affirmations suivantes sont vraies ou fausses. Justifier brièvement chaque réponse (contre-exemple ou propriété).
\begin{enumerate}
    \item Toute suite bornée est convergente.
    \item Si $(u_n)$ et $(v_n)$ convergent vers $\ell$ et $\ell'$ respectivement, alors $(u_n + v_n)$ converge vers $\ell + \ell'$.
    \item Si $(u_n)$ est croissante et majorée, alors elle converge vers son majorant le plus petit.
    \item La suite définie par $u_n = (-1)^n$ admet deux valeurs d'adhérence.
    \item Si $\lim_{n\to +\infty} u_n = +\infty$ et $\lim_{n\to +\infty} v_n = -\infty$, alors la limite de $(u_n + v_n)$ est indéterminée.
\end{enumerate}
\end{exercice}

\begin{exercice}[Calculs directs de limites]
Calculer les limites suivantes (finies ou infinies) en justifiant par les opérations sur les limites ou les suites de référence.
\begin{enumerate}
    \item $\displaystyle \lim_{n\to +\infty} \frac{3n^2 - 5n + 2}{2n^2 + 7}$
    \item $\displaystyle \lim_{n\to +\infty} \left( \sqrt{n+1} - \sqrt{n} \right)$
    \item $\displaystyle \lim_{n\to +\infty} \frac{2^n + 3^n}{4^n}$
    \item $\displaystyle \lim_{n\to +\infty} \frac{n^3 + 2n}{5n^2 - 1}$
    \item $\displaystyle \lim_{n\to +\infty} \left(1 + \frac{1}{n}\right)^{2n}$
\end{enumerate}
\end{exercice}

\begin{exercice}[Raisonnement par récurrence : initiation]
Soit la suite $(u_n)$ définie par $u_0 = 3$ et $u_{n+1} = 2u_n - 1$ pour tout $n \in \mathbb{N}$.
\begin{enumerate}
    \item Calculer $u_1$, $u_2$, $u_3$. Conjecturer une expression explicite de $u_n$ en fonction de $n$.
    \item Démontrer par récurrence que pour tout $n \in \mathbb{N}$, $u_n = 2^{n+1} + 1$.
    \item En déduire la limite de la suite $(u_n)$.
\end{enumerate}
\end{exercice}

\courssub{Je m'entraîne}

\begin{exercice}[Suite explicite : étude complète]
Soit la suite $(u_n)$ définie pour tout $n \in \mathbb{N}^*$ par $u_n = \frac{4n - 1}{2n + 3}$.
\begin{enumerate}
    \item Étudier la monotonie de $(u_n)$ en calculant $u_{n+1} - u_n$.
    \item Montrer que la suite est bornée. Préciser un majorant et un minorant.
    \item Calculer $\lim_{n\to +\infty} u_n$.
    \item Déterminer le plus petit entier $N$ tel que pour tout $n \ge N$, $|u_n - 2| < 0,01$.
\end{enumerate}
\end{exercice}

\begin{exercice}[Suite récurrente linéaire : bornitude et convergence]
Soit la suite $(v_n)$ définie par $v_0 = 0$ et $v_{n+1} = 0,6 v_n + 4$ pour tout $n \in \mathbb{N}$.
\begin{enumerate}
    \item Démontrer par récurrence que pour tout $n \in \mathbb{N}$, $0 \le v_n \le 10$.
    \item Étudier la monotonie de $(v_n)$.
    \item En déduire que $(v_n)$ converge. Calculer sa limite $\ell$.
    \item Exprimer $v_n - \ell$ en fonction de $n$ et en déduire le plus petit $n$ tel que $|v_n - \ell| < 10^{-3}$.
\end{enumerate}
\end{exercice}

\begin{exercice}[Suite récurrente de type Héron : racine carrée]
Soit la suite $(w_n)$ définie par $w_0 = 1$ et $w_{n+1} = \frac{1}{2}\left(w_n + \frac{7}{w_n}\right)$ pour tout $n \in \mathbb{N}$.
\begin{enumerate}
    \item Montrer que pour tout $n \in \mathbb{N}$, $w_n > 0$.
    \item Étudier le signe de $w_{n+1} - w_n$ en fonction de $w_n^2 - 7$. En déduire la monotonie de la suite à partir du rang 1.
    \item Montrer que la suite converge et calculer sa limite.
    \item Interpréter géométriquement cette suite (méthode de Héron / Newton pour $\sqrt{7}$).
\end{enumerate}
\end{exercice}

\begin{popfigure}
\begin{tikzpicture}[scale=0.8, domain=0.5:4]
\draw[->] (0,0) -- (4.5,0) node[right] {$x$};
\draw[->] (0,0) -- (0,4.5) node[above] {$y$};
\draw[popcurve, thick] plot (\x, {7/\x}) node[right] {$y=7/x$};
\draw[popblue, thick] plot (\x, {\x}) node[above right] {$y=x$};
\draw[poporange, thick, dashed] plot (\x, {0.5*(\x+7/\x)}) node[above right] {$y=f(x)$};
% Cobweb for w0=1
\draw[popgreen, thick] (1,0) -- (1,1);
\draw[popgreen, thick] (1,1) -- (4,1);
\draw[popgreen, thick] (4,1) -- (4,2.875);
\draw[popgreen, thick] (4,2.875) -- (2.875,2.875);
\draw[popgreen, thick] (2.875,2.875) -- (2.875,2.65);
\node[popgreen] at (1,-0.3) {$w_0=1$};
\node[popgreen] at (4.3,1) {$w_1=4$};
\node[popgreen] at (2.5,3.2) {$w_2\approx 2,88$};
\draw[poppurple, fill=poppurple] (2.645,2.645) circle (2pt) node[above left] {$\sqrt{7}$};
\end{tikzpicture}
\legende{Construction graphique (toile d'araignée) de la suite de Héron pour $\sqrt{7}$.}
\end{popfigure}

\begin{exercice}[Modélisation : Épargne avec versements réguliers]
Un technicien place 2\,000\,000 FCFA sur un compte bloqué rapportant 5\% d'intérêts par an (capitalisation annuelle). Il verse en plus 300\,000 FCFA à la fin de chaque année.
On note $C_n$ le capital (en FCFA) au bout de $n$ années ($C_0 = 2\,000\,000$).
\begin{enumerate}
    \item Exprimer $C_{n+1}$ en fonction de $C_n$. Quelle est la nature de cette suite ?
    \item Démontrer que la suite $(C_n + 6\,000\,000)$ est géométrique. En déduire l'expression explicite de $C_n$.
    \item Calculer le capital au bout de 10 ans (arrondi au FCFA près).
    \item Au bout de combien d'années le capital dépasse-t-il 10\,000\,000 FCFA ?
\end{enumerate}
\end{exercice}

\courssub{Je me prépare au Probatoire}

\begin{exercice}[Probatoire Type 1 : Suite rationnelle — Étude complète]
\label{exo:proba1}
Soit la suite $(U_n)$ définie pour tout $n \in \mathbb{N}^*$ par :
\[ U_n = \frac{3n^2 - 2n + 1}{n^2 + 4} \]
\begin{enumerate}
    \item Calculer $U_1$, $U_2$, $U_3$. Quelle conjecture peut-on faire sur la monotonie ?
    \item Étudier la monotonie de la suite $(U_n)$ en calculant $U_{n+1} - U_n$ (ou le rapport). Préciser le sens de variation.
    \item Montrer que la suite $(U_n)$ est bornée. Donner un majorant et un minorant.
    \item Calculer $\lim\limits_{n\to +\infty} U_n$.
    \item Déterminer le plus petit entier naturel $N$ tel que pour tout $n \ge N$, $|U_n - 3| < 0,01$.
\end{enumerate}
\end{exercice}

\begin{exercice}[Probatoire Type 2 : Suite récurrente linéaire — Seuil et convergence]
\label{exo:proba2}
Soit la suite $(U_n)$ définie par $U_0 = 5$ et pour tout $n \in \mathbb{N}$ :
\[ U_{n+1} = \frac{U_n + 2}{3} \]
\begin{enumerate}
    \item Démontrer par récurrence que pour tout $n \in \mathbb{N}$, $1 \le U_n \le 5$.
    \item Étudier la monotonie de la suite $(U_n)$.
    \item En déduire que la suite converge. Calculer sa limite $\ell$.
    \item Exprimer $U_n - \ell$ en fonction de $n$.
    \item Déterminer le plus petit entier $n$ tel que $U_n < 1,1$.
\end{enumerate}
\end{exercice}

\begin{exercice}[Probatoire Type 3 : Modélisation — Pollution d'un lac (Suite par morceaux)]
\label{exo:proba3}
Une usine rejette chaque mois 100 kg de polluant dans un lac. Le courant d'évacuation renouvelle 10\% du volume d'eau du lac chaque mois, emportant avec lui 10\% du polluant présent.
On note $P_n$ la quantité de polluant (en kg) dans le lac au début du mois $n$ (avant le nouveau rejet), avec $P_0 = 0$.
\begin{enumerate}
    \item Justifier que pour tout $n \in \mathbb{N}$, $P_{n+1} = 0,9 P_n + 100$.
    \item Démontrer que la suite $(P_n - 1000)$ est géométrique. En déduire l'expression explicite de $P_n$.
    \item Étudier la convergence de $(P_n)$. Interpréter la limite dans le contexte.
    \item Au bout de combien de mois la quantité de polluant dépasse-t-elle 800 kg ?
    \item \textbf{Variante :} À partir du 12ème mois (i.e. pour $n \ge 12$), l'usine réduit son rejet à 50 kg/mois. La relation de récurrence devient :
    \[
    \begin{cases}
    P_{n+1} = 0,9 P_n + 100 & \text{pour } 0 \le n \le 11 \\
    P_{n+1} = 0,9 P_n + 50  & \text{pour } n \ge 12
    \end{cases}
    \]
    Calculer $P_{12}$. Étudier la convergence de la suite à partir de $n=12$. Quelle est la nouvelle limite ?
\end{enumerate}
\end{exercice}

\newpage

\coursec{Corrigés des exercices}

\begin{corrige}[Exercice 1 — QCM : Vocabulaire et propriétés de base]
\begin{enumerate}
    \item \textbf{Bonne réponse : B.}
    \emph{Justification :} Par définition, une suite $(u_n)$ est croissante si $\forall n \in \mathbb{N},\ u_{n+1} \ge u_n$. L'option A correspond à \emph{strictement} croissante. L'option C est trop faible (un seul terme). L'option D correspond à décroissante.
    \item \textbf{Bonne réponse : A.}
    \emph{Justification :} C'est la définition formelle de la limite $+\infty$ : $\forall A > 0,\ \exists N \in \mathbb{N},\ \forall n \ge N,\ u_n > A$. L'option B dit qu'il existe un $A$ (faux, ça doit marcher pour tout $A$). L'option C a $n \le N$ (inverse). L'option D dit que la suite est minorée par un $A>0$ fixe (bornitude, pas divergence).
    \item \textbf{Bonne réponse : B.}
    \emph{Justification :} C'est un théorème fondamental : toute suite convergente vers un réel est bornée. La réciproque est fausse (ex: $u_n = (-1)^n$). Une suite convergente n'est pas nécessairement monotone (ex: $u_n = \frac{(-1)^n}{n}$). Elle n'est pas constante à partir d'un rang (sauf cas trivial). $(|u_n|)$ converge vers $|\ell|$.
    \item \textbf{Bonne réponse : A.}
    \emph{Justification :} Pour une suite récurrente $u_{n+1}=f(u_n)$ avec $f$ croissante sur un intervalle $I$ contenant les termes, la monotonie de $(u_n)$ est celle de la comparaison entre $u_1$ et $u_0$ (ou signe de $f(x)-x$). Si $u_1 \ge u_0$, la suite est croissante ; sinon décroissante. Comparer $f(x)$ et $x$ (B) donne le sens mais la méthode standard au programme consiste à comparer $u_1$ et $u_0$ après avoir vérifié que $f$ est croissante et que les termes restent dans $I$.
\end{enumerate}
\end{corrige}

\begin{corrige}[Exercice 2 — Vrai / Faux justifié]
\begin{enumerate}
    \item \textbf{Faux.}
    Contre-exemple : $u_n = (-1)^n$. Elle est bornée (par $1$ et $-1$) mais diverge (deux valeurs d'adhérence : $1$ et $-1$).
    \item \textbf{Vrai.}
    C'est la propriété de la limite de la somme : si $\lim u_n = \ell$ et $\lim v_n = \ell'$ (finies), alors $\lim (u_n+v_n) = \ell+\ell'$.
    \item \textbf{Vrai.}
    C'est le théorème de la suite monotone bornée (exigible au Probatoire). Une suite croissante et majorée converge vers son supremum (majorant le plus petit).
    \item \textbf{Vrai.}
    La suite $u_n = (-1)^n$ prend les valeurs $1$ (rangs pairs) et $-1$ (rangs impairs). Les sous-suites $(u_{2n})$ et $(u_{2n+1})$ convergent respectivement vers $1$ et $-1$. Ce sont les deux valeurs d'adhérence (limites de sous-suites).
    \item \textbf{Vrai.}
    $\lim u_n = +\infty$ et $\lim v_n = -\infty$ donnent une forme indéterminée $\infty - \infty$ (ou $+\infty + (-\infty)$). La limite de la somme peut être $+\infty$, $-\infty$, un réel, ou ne pas exister. Exemples : $u_n=n, v_n=-n \to 0$ ; $u_n=n^2, v_n=-n \to +\infty$ ; $u_n=n, v_n=-n^2 \to -\infty$.
\end{enumerate}
\end{corrige}

\begin{corrige}[Exercice 3 — Calculs directs de limites]
\begin{enumerate}
    \item $\displaystyle \lim_{n\to +\infty} \frac{3n^2 - 5n + 2}{2n^2 + 7}$.
    Factoriser par $n^2$ (plus haut degré) au numérateur et dénominateur :
    \[ \frac{n^2(3 - \frac{5}{n} + \frac{2}{n^2})}{n^2(2 + \frac{7}{n^2})} = \frac{3 - \frac{5}{n} + \frac{2}{n^2}}{2 + \frac{7}{n^2}} \xrightarrow[n\to+\infty]{} \frac{3 - 0 + 0}{2 + 0} = \frac{3}{2}. \]
    \item $\displaystyle \lim_{n\to +\infty} \left( \sqrt{n+1} - \sqrt{n} \right)$.
    Forme $\infty - \infty$. Multiplier par le conjugué :
    \[ \sqrt{n+1} - \sqrt{n} = \frac{(\sqrt{n+1} - \sqrt{n})(\sqrt{n+1} + \sqrt{n})}{\sqrt{n+1} + \sqrt{n}} = \frac{1}{\sqrt{n+1} + \sqrt{n}}. \]
    Dénominateur $\to +\infty$, donc la limite est $0$.
    \item $\displaystyle \lim_{n\to +\infty} \frac{2^n + 3^n}{4^n}$.
    Réécrire : $\frac{2^n}{4^n} + \frac{3^n}{4^n} = \left(\frac{1}{2}\right)^n + \left(\frac{3}{4}\right)^n$.
    Comme $\frac{1}{2} < 1$ et $\frac{3}{4} < 1$, les deux suites géométriques tendent vers $0$. Limite = $0$.
    \item $\displaystyle \lim_{n\to +\infty} \frac{n^3 + 2n}{5n^2 - 1}$.
    Degré numérateur (3) $>$ degré dénominateur (2). Diviser numérateur et dénominateur par $n^2$ :
    \[ \frac{n + \frac{2}{n}}{5 - \frac{1}{n^2}} \sim \frac{n}{5} \xrightarrow[n\to+\infty]{} +\infty. \]
    \item $\displaystyle \lim_{n\to +\infty} \left(1 + \frac{1}{n}\right)^{2n}$.
    On connaît $\lim_{n\to+\infty} \left(1 + \frac{1}{n}\right)^n = e$.
    $\left(1 + \frac{1}{n}\right)^{2n} = \left[\left(1 + \frac{1}{n}\right)^n\right]^2 \xrightarrow[n\to+\infty]{} e^2$.
\end{enumerate}
\end{corrige}

\begin{corrige}[Exercice 4 — Raisonnement par récurrence : initiation]
\begin{enumerate}
    \item Calculs :
    $u_0 = 3$.
    $u_1 = 2\times 3 - 1 = 5$.
    $u_2 = 2\times 5 - 1 = 9$.
    $u_3 = 2\times 9 - 1 = 17$.
    On observe : $3 = 2^1+1$, $5 = 2^2+1$, $9 = 2^3+1$, $17 = 2^4+1$.
    Conjecture : $\forall n \in \mathbb{N},\ u_n = 2^{n+1} + 1$.
    \item \textbf{Démonstration par récurrence.}
    \textbf{Propriété $\mathcal{P}(n)$ :} $u_n = 2^{n+1} + 1$.
    \textbf{Initialisation ($n=0$) :} $u_0 = 3$ et $2^{0+1}+1 = 2+1=3$. $\mathcal{P}(0)$ vraie.
    \textbf{Hérédité :} Soit $k \in \mathbb{N}$. Supposons $\mathcal{P}(k)$ vraie, c.-à-d. $u_k = 2^{k+1} + 1$.
    Calculons $u_{k+1}$ par la relation de récurrence :
    \begin{align*}
    u_{k+1} &= 2u_k - 1 \\
    &= 2(2^{k+1} + 1) - 1 \quad \text{(hypothèse de récurrence)} \\
    &= 2^{k+2} + 2 - 1 \\
    &= 2^{(k+1)+1} + 1.
    \end{align*}
    Donc $\mathcal{P}(k+1)$ est vraie.
    \textbf{Conclusion :} Par le principe de récurrence, $\forall n \in \mathbb{N},\ u_n = 2^{n+1} + 1$.
    \item Limite : $u_n = 2^{n+1} + 1$. Comme $2 > 1$, $\lim_{n\to+\infty} 2^{n+1} = +\infty$. Donc $\lim_{n\to+\infty} u_n = +\infty$. La suite diverge vers $+\infty$.
\end{enumerate}
\end{corrige}

\begin{corrige}[Exercice 5 — Suite explicite : étude complète]
Soit $u_n = \frac{4n - 1}{2n + 3}$ pour $n \in \mathbb{N}^*$.
\begin{enumerate}
    \item \textbf{Monotonie :} Étudions $u_{n+1} - u_n$.
    \begin{align*}
    u_{n+1} - u_n &= \frac{4(n+1)-1}{2(n+1)+3} - \frac{4n-1}{2n+3} \\
    &= \frac{4n+3}{2n+5} - \frac{4n-1}{2n+3} \\
    &= \frac{(4n+3)(2n+3) - (4n-1)(2n+5)}{(2n+5)(2n+3)}.
    \end{align*}
    Calcul du numérateur :
    $(4n+3)(2n+3) = 8n^2 + 12n + 6n + 9 = 8n^2 + 18n + 9$.
    $(4n-1)(2n+5) = 8n^2 + 20n - 2n - 5 = 8n^2 + 18n - 5$.
    Différence : $(8n^2 + 18n + 9) - (8n^2 + 18n - 5) = 14$.
    Donc $u_{n+1} - u_n = \frac{14}{(2n+5)(2n+3)}$.
    Pour $n \ge 1$, dénominateur $> 0$, donc $u_{n+1} - u_n > 0$.
    La suite $(u_n)$ est \textbf{strictement croissante}.
    \item \textbf{Bornitude :}
    Comme la suite est croissante, elle est minorée par son premier terme : $u_1 = \frac{4-1}{2+3} = \frac{3}{5} = 0,6$.
    Pour le majorant, calculons la limite (qui sera le majorant le plus petit) :
    $u_n = \frac{4n-1}{2n+3} = \frac{4 - \frac{1}{n}}{2 + \frac{3}{n}} \xrightarrow[n\to+\infty]{} \frac{4}{2} = 2$.
    Comme la suite est croissante et converge vers $2$, on a $u_n < 2$ pour tout $n$.
    La suite est bornée : $\forall n \ge 1,\ 0,6 \le u_n < 2$.
    (Minorant : $0,6$ ; Majorant : $2$).
    \item \textbf{Limite :} $\lim_{n\to+\infty} u_n = 2$ (calculé ci-dessus).
    \item \textbf{Seuil :} On cherche $N$ tel que $\forall n \ge N,\ |u_n - 2| < 0,01$.
    Comme $u_n < 2$ (suite croissante vers 2), $|u_n - 2| = 2 - u_n$.
    $2 - u_n = 2 - \frac{4n-1}{2n+3} = \frac{2(2n+3) - (4n-1)}{2n+3} = \frac{4n+6-4n+1}{2n+3} = \frac{7}{2n+3}$.
    Inéquation : $\frac{7}{2n+3} < 0,01 = \frac{1}{100}$.
    $\iff 700 < 2n + 3 \iff 697 < 2n \iff n > 348,5$.
    Le plus petit entier $N$ est \textbf{349}.
    (Vérification : $n=349 \Rightarrow \frac{7}{698+3}=\frac{7}{701}\approx 0,00998 < 0,01$. $n=348 \Rightarrow \frac{7}{699}\approx 0,01001 > 0,01$).
\end{enumerate}
\end{corrige}

\begin{corrige}[Exercice 6 — Suite récurrente linéaire : bornitude et convergence]
Soit $v_0 = 0$, $v_{n+1} = 0,6 v_n + 4$.
\begin{enumerate}
    \item \textbf{Bornitude par récurrence.}
    \textbf{Propriété $\mathcal{P}(n)$ :} $0 \le v_n \le 10$.
    \textbf{Initialisation ($n=0$) :} $v_0 = 0$. $0 \le 0 \le 10$. Vrai.
    \textbf{Hérédité :} Supposons $0 \le v_k \le 10$ pour un $k \in \mathbb{N}$.
    $v_{k+1} = 0,6 v_k + 4$.
    Comme $0 \le v_k \le 10$, en multipliant par $0,6 > 0$ : $0 \le 0,6 v_k \le 6$.
    En ajoutant $4$ : $4 \le 0,6 v_k + 4 \le 10$.
    Donc $0 \le 4 \le v_{k+1} \le 10$. $\mathcal{P}(k+1)$ vraie.
    \textbf{Conclusion :} $\forall n \in \mathbb{N},\ 0 \le v_n \le 10$.
    \item \textbf{Monotonie :}
    $f(x) = 0,6x + 4$. $f'(x) = 0,6 > 0$, $f$ est strictement croissante sur $\mathbb{R}$.
    $v_1 = 0,6\times 0 + 4 = 4$. $v_1 = 4 > 0 = v_0$.
    Comme $f$ est croissante et $v_1 > v_0$, la suite $(v_n)$ est \textbf{strictement croissante}.
    \item \textbf{Convergence et limite :}
    La suite est croissante et majorée (par 10), donc elle converge (théorème de la suite monotone bornée).
    Soit $\ell$ sa limite. En passant à la limite dans $v_{n+1} = 0,6 v_n + 4$ :
    $\ell = 0,6 \ell + 4 \iff 0,4 \ell = 4 \iff \ell = 10$.
    \item \textbf{Expression de $v_n - \ell$ et seuil :}
    $v_{n+1} - 10 = 0,6 v_n + 4 - 10 = 0,6 v_n - 6 = 0,6(v_n - 10)$.
    La suite $(v_n - 10)$ est géométrique de raison $0,6$ et premier terme $v_0 - 10 = -10$.
    $v_n - 10 = -10 \times (0,6)^n \implies v_n = 10 - 10 \times (0,6)^n$.
    Inéquation : $|v_n - 10| < 10^{-3} \iff 10 \times (0,6)^n < 0,001 \iff (0,6)^n < 10^{-4}$.
    $\ln((0,6)^n) < \ln(10^{-4}) \iff n \ln(0,6) < -4 \ln(10)$.
    Comme $\ln(0,6) < 0$, en divisant : $n > \frac{-4 \ln(10)}{\ln(0,6)} \approx \frac{-9,2103}{-0,5108} \approx 18,03$.
    Le plus petit entier $n$ est \textbf{19}.
    (Vérif : $n=18 \to 10\times 0,6^{18} \approx 0,00101 > 0,001$ ; $n=19 \to \approx 0,0006 < 0,001$).
\end{enumerate}
\end{corrige}

\begin{corrige}[Exercice 7 — Suite récurrente de type Héron : racine carrée]
Soit $w_0 = 1$, $w_{n+1} = \frac{1}{2}\left(w_n + \frac{7}{w_n}\right)$.
\begin{enumerate}
    \item \textbf{Positivité par récurrence.}
    $\mathcal{P}(n) : w_n > 0$.
    $w_0 = 1 > 0$. Vrai.
    Si $w_k > 0$, alors $\frac{7}{w_k} > 0$, donc $w_{k+1} = \frac{1}{2}(w_k + \frac{7}{w_k}) > 0$.
    Par récurrence, $\forall n,\ w_n > 0$.
    \item \textbf{Monotonie :}
    $w_{n+1} - w_n = \frac{1}{2}\left(w_n + \frac{7}{w_n}\right) - w_n = \frac{1}{2}\left(\frac{7}{w_n} - w_n\right) = \frac{7 - w_n^2}{2w_n}$.
    Comme $w_n > 0$, le signe de $w_{n+1} - w_n$ est le signe de $7 - w_n^2$.
    $w_0 = 1 \Rightarrow w_0^2 = 1 < 7 \Rightarrow w_1 > w_0$.
    $w_1 = \frac{1}{2}(1+7) = 4 \Rightarrow w_1^2 = 16 > 7 \Rightarrow w_2 < w_1$.
    Pour $n \ge 1$, montrons que $w_n \ge \sqrt{7}$ (donc $w_n^2 \ge 7$).
    Pour $x > 0$, $f(x) = \frac{1}{2}(x + \frac{7}{x}) \ge \sqrt{7}$ (inégalité arithmético-géométrique : $\frac{x + 7/x}{2} \ge \sqrt{x \cdot 7/x} = \sqrt{7}$).
    Donc $\forall n \ge 1,\ w_n \ge \sqrt{7} \implies w_n^2 \ge 7 \implies w_{n+1} - w_n \le 0$.
    La suite est \textbf{croissante de $n=0$ à $n=1$}, puis \textbf{décroissante à partir du rang 1} (pour $n \ge 1$).
    \item \textbf{Convergence et limite :}
    Pour $n \ge 1$, la suite est décroissante et minorée par $\sqrt{7} > 0$. Elle converge.
    Soit $\ell$ la limite. $\ell = \frac{1}{2}(\ell + \frac{7}{\ell}) \iff 2\ell = \ell + \frac{7}{\ell} \iff \ell = \frac{7}{\ell} \iff \ell^2 = 7$.
    Comme $\ell \ge \sqrt{7} > 0$, $\ell = \sqrt{7}$.
    \item \textbf{Interprétation :} C'est la méthode de Héron (ou Newton) pour calculer $\sqrt{7}$. On cherche une racine de $x^2 - 7 = 0$. La tangente en $(w_n, w_n^2-7)$ coupe l'axe des abscisses en $w_{n+1} = w_n - \frac{w_n^2-7}{2w_n} = \frac{1}{2}(w_n + \frac{7}{w_n})$.
\end{enumerate}
\end{corrige}

\begin{corrige}[Exercice 8 — Modélisation : Épargne avec versements réguliers]
\begin{enumerate}
    \item Intérêts 5\,\% : multiplicateur $1,05$. Versement \num{300000} en fin d'année.
    $C_{n+1} = 1,05 C_n + 300\,000$.
    $C_0 = 2\,000\,000$.
    C'est une suite \textbf{affine} (récurrente linéaire d'ordre 1).
    \item On cherche $\alpha$ tel que $(C_n + \alpha)$ soit géométrique.
    $C_{n+1} + \alpha = 1,05 C_n + 300\,000 + \alpha = 1,05(C_n + \alpha) + 300\,000 + \alpha - 1,05\alpha$.
    On veut $300\,000 + \alpha - 1,05\alpha = 0 \iff 300\,000 - 0,05\alpha = 0 \iff \alpha = \frac{300\,000}{0,05} = 6\,000\,000$.
    Donc $(C_n + 6\,000\,000)$ est géométrique de raison $1,05$.
    Premier terme : $C_0 + 6\,000\,000 = 8\,000\,000$.
    $C_n + 6\,000\,000 = 8\,000\,000 \times (1,05)^n$.
    $C_n = 8\,000\,000 \times (1,05)^n - 6\,000\,000$.
    \item $n=10$ : $C_{10} = 8\,000\,000 \times (1,05)^{10} - 6\,000\,000$.
    $(1,05)^{10} \approx 1,6288946$.
    $C_{10} \approx 8\,000\,000 \times 1,6288946 - 6\,000\,000 = 13\,031\,156,8 - 6\,000\,000 = 7\,031\,157$ FCFA (arrondi).
    \item On cherche $n$ tel que $C_n > 10\,000\,000$.
    $8\,000\,000 \times (1,05)^n - 6\,000\,000 > 10\,000\,000 \iff 8\,000\,000 \times (1,05)^n > 16\,000\,000 \iff (1,05)^n > 2$.
    $n \ln(1,05) > \ln 2 \iff n > \frac{\ln 2}{\ln 1,05} \approx \frac{0,6931}{0,04879} \approx 14,2$.
    Le plus petit entier $n$ est \textbf{15} ans.
    (Vérif : $n=14 \to C_{14} \approx 9\,812\,000$ ; $n=15 \to C_{15} \approx 10\,602\,000$).
\end{enumerate}
\end{corrige}

\begin{corrige}[Exercice 9 — Probatoire Type 1 : Suite rationnelle — Étude complète]
\textbf{Barème indicatif (sur 20 pts) :} Q1 (2 pts), Q2 (5 pts), Q3 (3 pts), Q4 (4 pts), Q5 (6 pts).
\textbf{Points de vigilance :} Justifier le signe du numérateur de $U_{n+1}-U_n$ (factorisation ou étude trinôme). Pour la majoration, utiliser la monotonie et la limite. Pour Q5, bien gérer la valeur absolue en fonction de la position par rapport à la limite.
\begin{enumerate}
    \item $U_1 = \frac{3-2+1}{1+4} = \frac{2}{5} = 0,4$.
    $U_2 = \frac{12-4+1}{4+4} = \frac{9}{8} = 1,125$.
    $U_3 = \frac{27-6+1}{9+4} = \frac{22}{13} \approx 1,69$.
    Conjecture : la suite semble \textbf{croissante}.
    \item \textbf{Monotonie :} Pour $n \ge 1$,
    \begin{align*}
    U_{n+1} - U_n &= \frac{3(n+1)^2 - 2(n+1) + 1}{(n+1)^2 + 4} - \frac{3n^2 - 2n + 1}{n^2 + 4} \\
    &= \frac{3n^2+4n+2}{n^2+2n+5} - \frac{3n^2-2n+1}{n^2+4}.
    \end{align*}
    Mise au même dénominateur $(n^2+2n+5)(n^2+4) > 0$.
    Numérateur $N_n = (3n^2+4n+2)(n^2+4) - (3n^2-2n+1)(n^2+2n+5)$.
    Développement :
    $(3n^2+4n+2)(n^2+4) = 3n^4 + 12n^2 + 4n^3 + 16n + 2n^2 + 8 = 3n^4 + 4n^3 + 14n^2 + 16n + 8$.
    $(3n^2-2n+1)(n^2+2n+5) = 3n^4 + 6n^3 + 15n^2 - 2n^3 - 4n^2 - 10n + n^2 + 2n + 5 = 3n^4 + 4n^3 + 12n^2 - 8n + 5$.
    $N_n = (3n^4 + 4n^3 + 14n^2 + 16n + 8) - (3n^4 + 4n^3 + 12n^2 - 8n + 5) = 2n^2 + 24n + 3$.
    Pour $n \ge 1$, $N_n > 0$. Dénominateur $> 0$.
    Donc $U_{n+1} - U_n > 0$. La suite $(U_n)$ est \textbf{strictement croissante} sur $\mathbb{N}^*$.
    \item \textbf{Bornitude :}
    Croissante $\Rightarrow$ minorée par $U_1 = 0,4$.
    Limite (Q4) = 3. Comme elle est croissante vers 3, elle est majorée par 3.
    $\forall n \ge 1,\ 0,4 \le U_n < 3$. La suite est bornée.
    \item \textbf{Limite :}
    $U_n = \frac{3n^2 - 2n + 1}{n^2 + 4} = \frac{3 - \frac{2}{n} + \frac{1}{n^2}}{1 + \frac{4}{n^2}} \xrightarrow[n\to+\infty]{} \frac{3}{1} = 3$.
    \item \textbf{Seuil :} On veut $|U_n - 3| < 0,01$.
    $U_n - 3 = \frac{3n^2-2n+1 - 3(n^2+4)}{n^2+4} = \frac{-2n - 11}{n^2+4} < 0$ pour $n\ge 1$.
    Donc $|U_n - 3| = 3 - U_n = \frac{2n+11}{n^2+4}$.
    Inéquation : $\frac{2n+11}{n^2+4} < \frac{1}{100}$.
    $\iff 100(2n+11) < n^2 + 4 \iff n^2 - 200n - 1100 + 4 > 0 \iff n^2 - 200n - 1096 > 0$.
    Racines de $n^2 - 200n - 1096 = 0$ : $\Delta = 40000 + 4384 = 44384$. $\sqrt{\Delta} \approx 210,67$.
    $n = \frac{200 \pm 210,67}{2}$. Racine positive $\approx 205,3$.
    Le trinôme est positif pour $n > 205,3$.
    Le plus petit entier $N$ est \textbf{206}.
    (Vérif : $n=205 \to \frac{421}{42029} \approx 0,010017 > 0,01$ ; $n=206 \to \frac{423}{42440} \approx 0,009967 < 0,01$).
\end{enumerate}
\end{corrige}

\begin{corrige}[Exercice 10 — Probatoire Type 2 : Suite récurrente linéaire — Seuil et convergence]
\textbf{Barème indicatif (sur 20 pts) :} Q1 (4 pts), Q2 (3 pts), Q3 (4 pts), Q4 (4 pts), Q5 (5 pts).
\textbf{Points de vigilance :} Hérédité bien rédigée (encadrement conservé par $f$). Monotonie : citer $f$ croissante et comparer $U_1$ et $U_0$. Limite : passage à la limite justifié par convergence. Expression $U_n-\ell$ : suite géométrique. Seuil : utilisation correcte des logarithmes (inversion inéquation car $\ln q < 0$).
\begin{enumerate}
    \item \textbf{Bornitude par récurrence.}
    $\mathcal{P}(n) : 1 \le U_n \le 5$.
    \textbf{Initialisation ($n=0$) :} $U_0 = 5$. $1 \le 5 \le 5$. Vrai.
    \textbf{Hérédité :} Supposons $1 \le U_k \le 5$.
    $U_{k+1} = \frac{U_k + 2}{3}$.
    $1 \le U_k \le 5 \implies 3 \le U_k + 2 \le 7 \implies 1 \le \frac{U_k+2}{3} \le \frac{7}{3} \approx 2,33$.
    Donc $1 \le U_{k+1} \le \frac{7}{3} \le 5$. $\mathcal{P}(k+1)$ vraie.
    \textbf{Conclusion :} $\forall n \in \mathbb{N},\ 1 \le U_n \le 5$.
    \item \textbf{Monotonie :}
    $f(x) = \frac{x+2}{3}$. $f'(x) = \frac{1}{3} > 0$, $f$ strictement croissante sur $\mathbb{R}$.
    $U_1 = \frac{5+2}{3} = \frac{7}{3} \approx 2,33$.
    $U_1 = \frac{7}{3} < 5 = U_0$.
    Comme $f$ est croissante et $U_1 < U_0$, la suite $(U_n)$ est \textbf{strictement décroissante}.
    \item \textbf{Convergence et limite :}
    Suite décroissante et minorée (par 1), donc convergente (théorème monotone bornée).
    Soit $\ell$ la limite. $\ell = \frac{\ell+2}{3} \iff 3\ell = \ell+2 \iff 2\ell=2 \iff \ell = 1$.
    \item \textbf{Expression de $U_n - \ell$ :}
    $U_{n+1} - 1 = \frac{U_n+2}{3} - 1 = \frac{U_n - 1}{3}$.
    $(U_n - 1)$ est géométrique de raison $\frac{1}{3}$, premier terme $U_0 - 1 = 4$.
    $U_n - 1 = 4 \times \left(\frac{1}{3}\right)^n \implies U_n = 1 + 4 \times 3^{-n}$.
    \item \textbf{Seuil :} $U_n < 1,1 \iff 1 + 4 \times 3^{-n} < 1,1 \iff 4 \times 3^{-n} < 0,1 \iff 3^{-n} < 0,025 \iff 3^n > 40$.
    $n \ln 3 > \ln 40 \iff n > \frac{\ln 40}{\ln 3} \approx \frac{3,6889}{1,0986} \approx 3,35$.
    Le plus petit entier $n$ est \textbf{4}.
    (Vérif : $U_3 = 1 + 4/27 \approx 1,148 > 1,1$ ; $U_4 = 1 + 4/81 \approx 1,049 < 1,1$).
\end{enumerate}
\end{corrige}

\begin{corrige}[Exercice 11 — Probatoire Type 3 : Modélisation — Pollution d'un lac]
\textbf{Barème indicatif (sur 20 pts) :} Q1 (3 pts), Q2 (5 pts), Q3 (4 pts), Q4 (4 pts), Q5 (4 pts).
\textbf{Points de vigilance :} Q1 : modélisation claire (10\,\% évacué = 90\,\% reste + 100 kg ajoutés). Q2 : trouver le bon décalage (1000) pour faire apparaître géométrique. Q4 : inéquation sur géométrique, logarithmes. Q5 : suite par morceaux, calculer $P_{12}$ avec 1ère formule, puis nouvelle suite géométrique pour $n\ge 12$ avec nouvelle limite.
\begin{enumerate}
    \item Chaque mois, 10\,\% de l'eau (et du polluant) partent. Il reste 90\,\% du polluant : $0,9 P_n$.
    L'usine ajoute \SI{100}{kg}.
    Donc $P_{n+1} = 0,9 P_n + 100$. $P_0 = 0$.
    \item On cherche $\alpha$ tel que $(P_n - \alpha)$ soit géométrique.
    $P_{n+1} - \alpha = 0,9 P_n + 100 - \alpha = 0,9(P_n - \alpha) + 100 - 0,1\alpha$.
    On annule $100 - 0,1\alpha = 0 \iff \alpha = 1000$.
    $(P_n - 1000)$ est géométrique de raison $0,9$.
    Premier terme : $P_0 - 1000 = -1000$.
    $P_n - 1000 = -1000 \times (0,9)^n \implies P_n = 1000 - 1000 \times (0,9)^n = 1000(1 - 0,9^n)$.
    \item $0,9 \in ]0,1[ \Rightarrow \lim_{n\to+\infty} 0,9^n = 0$.
    $\lim_{n\to+\infty} P_n = 1000$ kg.
    \emph{Interprétation :} La quantité de polluant se stabilise à \SI{1000}{kg}. C'est l'équilibre où la quantité évacuée (10\,\% de 1000 = 100 kg) égale la quantité rejetée (100 kg).
    \item $P_n > 800 \iff 1000(1 - 0,9^n) > 800 \iff 1 - 0,9^n > 0,8 \iff 0,9^n < 0,2$.
    $n \ln 0,9 < \ln 0,2 \iff n > \frac{\ln 0,2}{\ln 0,9} \approx \frac{-1,6094}{-0,10536} \approx 15,27$.
    Le plus petit entier $n$ est \textbf{16} mois.
    (Au début du 16ème mois, soit après 15 mois écoulés, $P_{15} \approx 794$ ; $P_{16} \approx 815$).
    \item \textbf{Variante :}
    Calcul de $P_{12}$ avec l'ancienne formule : $P_{12} = 1000(1 - 0,9^{12})$.
    $0,9^{12} \approx 0,28243$. $P_{12} \approx 1000(1 - 0,28243) = 717,57$ kg.
    Pour $n \ge 12$, $P_{n+1} = 0,9 P_n + 50$.
    Nouvelle limite $\ell'$ : $\ell' = 0,9\ell' + 50 \iff 0,1\ell' = 50 \iff \ell' = 500$ kg.
    On peut écrire $P_n - 500 = 0,9(P_{n-1} - 500)$ pour $n \ge 13$.
    La suite converge vers \textbf{500 kg} (nouvel équilibre : 10\,\% de 500 = 50 kg évacués = 50 kg rejetés).
\end{enumerate}
\end{corrige}

\newpage

\coursec{Fiche bilan}

\begin{popfigure}
\begin{center}\tcbox[colback=popink,colframe=popink,coltext=white,arc=9pt,boxsep=4pt,left=6pt,right=6pt]{\bfseries\small Suites Numériques}\end{center}
\vspace{-2pt}
\begin{tcbraster}[raster columns=3,raster equal height=rows,raster column skip=6pt,raster row skip=6pt,enhanced,blankest]
\begin{tcolorbox}[enhanced,colback=white,colframe=poporange,boxrule=0.9pt,arc=3pt,title={Raisonnement par récurrence},fonttitle=\bfseries\scriptsize,coltitle=white,colbacktitle=poporange,left=4pt,right=4pt,top=2pt,bottom=2pt,boxsep=1pt,toptitle=1.5pt,bottomtitle=1.5pt,halign=left]
\begin{itemize}[nosep,leftmargin=*,itemsep=1pt,font=\scriptsize]
\item Initialisation
\item Hérédité
\item Conclusion
\end{itemize}
\end{tcolorbox}
\begin{tcolorbox}[enhanced,colback=white,colframe=popgreen,boxrule=0.9pt,arc=3pt,title={Monotonie},fonttitle=\bfseries\scriptsize,coltitle=white,colbacktitle=popgreen,left=4pt,right=4pt,top=2pt,bottom=2pt,boxsep=1pt,toptitle=1.5pt,bottomtitle=1.5pt,halign=left]
\begin{itemize}[nosep,leftmargin=*,itemsep=1pt,font=\scriptsize]
\item Croissante / Décroissante
\item Méthode $U_{n+1}-U_n$
\item Méthode $\frac{U_{n+1}}{U_n}$
\end{itemize}
\end{tcolorbox}
\begin{tcolorbox}[enhanced,colback=white,colframe=poppurple,boxrule=0.9pt,arc=3pt,title={Bornitude \& Convergence},fonttitle=\bfseries\scriptsize,coltitle=white,colbacktitle=poppurple,left=4pt,right=4pt,top=2pt,bottom=2pt,boxsep=1pt,toptitle=1.5pt,bottomtitle=1.5pt,halign=left]
\begin{itemize}[nosep,leftmargin=*,itemsep=1pt,font=\scriptsize]
\item Majorée / Minorée / Bornée
\item Th. Croissante majorée
\item Th. Gendarmes
\end{itemize}
\end{tcolorbox}
\begin{tcolorbox}[enhanced,colback=white,colframe=poppink,boxrule=0.9pt,arc=3pt,title={Calcul de limites},fonttitle=\bfseries\scriptsize,coltitle=white,colbacktitle=poppink,left=4pt,right=4pt,top=2pt,bottom=2pt,boxsep=1pt,toptitle=1.5pt,bottomtitle=1.5pt,halign=left]
\begin{itemize}[nosep,leftmargin=*,itemsep=1pt,font=\scriptsize]
\item Opérations sur limites
\item Suites de référence
\item Formes indéterminées
\end{itemize}
\end{tcolorbox}
\begin{tcolorbox}[enhanced,colback=white,colframe=popgold,boxrule=0.9pt,arc=3pt,title={Suites récurrentes $U_{n+1}=f(U_n)$},fonttitle=\bfseries\scriptsize,coltitle=white,colbacktitle=popgold,left=4pt,right=4pt,top=2pt,bottom=2pt,boxsep=1pt,toptitle=1.5pt,bottomtitle=1.5pt,halign=left]
\begin{itemize}[nosep,leftmargin=*,itemsep=1pt,font=\scriptsize]
\item Étude de $f$
\item Point fixe $l=f(l)$
\item Convergence vers $l$
\end{itemize}
\end{tcolorbox}
\begin{tcolorbox}[enhanced,colback=white,colframe=popdark,boxrule=0.9pt,arc=3pt,title={Applications \& Probatoire},fonttitle=\bfseries\scriptsize,coltitle=white,colbacktitle=popdark,left=4pt,right=4pt,top=2pt,bottom=2pt,boxsep=1pt,toptitle=1.5pt,bottomtitle=1.5pt,halign=left]
\begin{itemize}[nosep,leftmargin=*,itemsep=1pt,font=\scriptsize]
\item Modélisation technique
\item Intérêts composés
\item Rédaction rigoureuse
\end{itemize}
\end{tcolorbox}
\end{tcbraster}

\legende{Carte mentale : Suites numériques -- Terminale Technique}
\end{popfigure}

\begin{bilan}
\courssub{Définitions clés}
\begin{itemize}
\item \textbf{Suite numérique :} fonction $u : \mathbb{N} \to \mathbb{R}$ (ou $\mathbb{N}^* \to \mathbb{R}$), notée $(u_n)$.
\item \textbf{Raisonnement par récurrence :} prouver $P(n_0)$ vrai (initialisation) ; supposer $P(k)$ vrai pour $k \ge n_0$ et déduire $P(k+1)$ vrai (hérésité) $\implies P(n)$ vrai $\forall n \ge n_0$.
\item \textbf{Monotonie :} $(u_n)$ \emph{croissante} si $u_{n+1} \ge u_n\ \forall n$ ; \emph{décroissante} si $u_{n+1} \le u_n\ \forall n$. \emph{Stricte} si inégalités strictes.
\item \textbf{Bornitude :} Majorée ($\exists M, u_n \le M$), minorée ($\exists m, u_n \ge m$), bornée (majorée et minorée).
\item \textbf{Convergence :} $(u_n)$ converge vers $\ell \in \mathbb{R}$ si $\forall \varepsilon > 0,\ \exists N \in \mathbb{N},\ \forall n \ge N,\ |u_n - \ell| < \varepsilon$. Notation : $\lim_{n\to+\infty} u_n = \ell$.
\item \textbf{Divergence :} limite $+\infty$ ou $-\infty$, ou pas de limite.
\end{itemize}

\courssub{Théorèmes \& Propriétés à connaître par cœur}
\begin{center}
\begin{tabularx}{\linewidth}{|>{\bfseries}l|X|}\hline
\textbf{Théorème / Propriété} & \textbf{Énoncé / Condition d'application} \\\hline
Principe de récurrence & Si $P(n_0)$ vrai et $[P(k) \implies P(k+1)]$ pour $k\ge n_0$, alors $P(n)$ vrai $\forall n\ge n_0$. \\\hline
Monotonie par différence & $u_{n+1}-u_n \ge 0 \iff$ croissante ; $\le 0 \iff$ décroissante. \\\hline
Monotonie par quotient & Si $u_n > 0$ : $\frac{u_{n+1}}{u_n} \ge 1 \iff$ croissante ; $\le 1 \iff$ décroissante. \\\hline
Th. Convergence (Croissante majorée) & Suite croissante et majorée $\implies$ Convergente (limite = supremum). \\\hline
Th. Convergence (Décroissante minorée) & Suite décroissante et minorée $\implies$ Convergente (limite = infimum). \\\hline
Th. des Gendarmes & Si $v_n \le u_n \le w_n$ et $\lim v_n = \lim w_n = \ell$, alors $\lim u_n = \ell$. \\\hline
Opérations sur limites & Somme, produit, quotient (si denom $\neq 0$), passage à la limite dans inégalités larges. \\\hline
Suites de référence & $\lim n^\alpha = +\infty\ (\alpha>0)$ ; $\lim q^n = 0\ (|q|<1),\ 1\ (q=1),\ +\infty\ (q>1)$ ; $\lim \frac{\ln n}{n}=0$. \\\hline
Suite récurrente $u_{n+1}=f(u_n)$ & Si $f$ continue sur $I$, $(u_n)\subset I$ convergente vers $\ell$, alors $\ell = f(\ell)$ (point fixe). \\\hline
Critère convergence récurrente & Si $f$ dérivable, $|f'(\ell)| < 1$ et $u_0$ assez proche $\implies$ convergence vers $\ell$. \\\hline
\end{tabularx}
\end{center}

\courssub{Méthodes express}
\begin{enumerate}
\item \textbf{Démo par récurrence :} 1. Initialisation ($n=0$ ou $1$). 2. Hérédité : \og Soit $k\ge n_0$ tel que $P(k)$ vrai \fg, calculer $P(k+1)$ en utilisant $P(k)$. 3. Conclusion formelle.
\item \textbf{Étude monotonie :} Calculer $u_{n+1}-u_n$ (factoriser) ou $\frac{u_{n+1}}{u_n}$ (si signe constant). Conclure sur le signe.
\item \textbf{Étude convergence :} 1. Montrer monotonie. 2. Montrer bornitude (majorée si croissante, minorée si décroissante). 3. Conclure convergence par théorème. 4. Calculer limite $\ell$ en passant à la limite dans $u_{n+1}=f(u_n)$.
\item \textbf{Calcul limites :} Factoriser le terme de plus haut degré (polynômes), utiliser suites de référence, lever indéterminations ($+\infty-\infty$, $0\times\infty$, $\frac{\infty}{\infty}$, $\frac{0}{0}$) par factorisation/conjugaison/règle de l'hôpital (si autorisée) ou encadrement.
\end{enumerate}
\end{bilan}

\begin{aretenir}[Checklist avant le Probatoire]
$\square$ Savoir rédiger un raisonnement par récurrence complet (Initialisation, Hérédité, Conclusion) sans sauter d'étapes.\par
$\square$ Maîtriser les deux méthodes d'étude de monotonie (différence et quotient) et savoir choisir la plus simple.\par
$\square$ Connaître par cœur le théorème : \og Suite croissante majorée (ou décroissante minorée) converge \fg et savoir l'appliquer.\par
$\square$ Savoir encadrer une suite pour utiliser le théorème des gendarmes.\par
$\square$ Maîtriser les opérations sur les limites et les suites de référence (puissances, géométriques, logarithmes).\par
$\square$ Savoir lever les formes indéterminées classiques ($\infty-\infty$, $0\times\infty$, $\frac{\infty}{\infty}$, $\frac{0}{0}$).\par
$\square$ Pour $u_{n+1}=f(u_n)$ : savoir étudier $f$ (variations, point fixe), montrer que $(u_n)$ reste dans un intervalle stable, montrer la monotonie, conclure la convergence et calculer la limite.\par
$\square$ Vérifier les conditions d'application des théorèmes (continuité de $f$, intervalle stable, $|f'(\ell)|<1$ si nécessaire).\par
$\square$ Savoir modéliser un problème concret (intérêts composés, population, amortissement, physique) par une suite explicite ou récurrente.\par
$\square$ Soigner la rédaction : phrases complètes, notations précises ($\forall, \exists, \implies, \iff$), conclusion encadrée.\par
$\square$ Relire les sujets types Probatoire : identification du type de suite, pièges (suite non monotone, divergence, forme indéterminée masquée).
\end{aretenir}

\findecours

\end{document}
